%========================================================================
%  Computational Financial Engineering I
%  Reconstructing Retail FX and CFD Market Access Under Shariah
%  Muhammad Taha Asif -- School of Mathematics and Statistics, USyd
%========================================================================
\documentclass[preprint,12pt,authoryear]{elsarticle}

% ---------- Packages ----------
\usepackage[T1]{fontenc}
\usepackage[utf8]{inputenc}
\usepackage{lmodern}
\usepackage{amsmath, amssymb, amsthm, mathtools}
\usepackage{bm}
\usepackage{booktabs}
\usepackage{graphicx}
\usepackage{siunitx}
\usepackage{xcolor}
\usepackage[hidelinks,breaklinks=true]{hyperref}
\usepackage[capitalise,noabbrev]{cleveref}
\usepackage{natbib}
\usepackage{setspace}
\usepackage{microtype}
\usepackage{fancyhdr}
\usepackage{orcidlink}

% ---------- Math macros ----------
\newcommand{\RR}{\mathbb{R}}
\newcommand{\EE}{\mathbb{E}}
\newcommand{\PP}{\mathbb{P}}
\newcommand{\VV}{\mathbb{V}}
\newcommand{\FF}{\mathcal{F}}
\newcommand{\CC}{\mathcal{C}}
\newcommand{\ind}[1]{\mathbf{1}_{\{#1\}}}
\newcommand{\dd}{\mathrm{d}}
\newcommand{\ie}{\emph{i.e.}}
\newcommand{\eg}{\emph{e.g.}}
\DeclareMathOperator*{\argmin}{arg\,min}
\DeclareMathOperator*{\argmax}{arg\,max}

% ---------- Theorem environments ----------
\newtheorem{proposition}{Proposition}[section]
\newtheorem{theorem}{Theorem}[section]
\newtheorem{lemma}{Lemma}[section]
\newtheorem{corollary}{Corollary}[section]
\theoremstyle{definition}
\newtheorem{definition}{Definition}[section]
\newtheorem{assumption}{Assumption}[section]
\newtheorem{remark}{Remark}[section]

% (No journal/preprint status line; suppresses the elsarticle
% "Preprint submitted to ..." banner on page 1.)
\journal{}

\begin{document}

% ---------- Running header / footer (plagiarism-adjacent attribution) ----
\pagestyle{fancy}
\fancyhf{}
\fancyhead[L]{\small\itshape M.\,T.\ Asif}
\fancyhead[R]{\small\itshape Shariah CFE~I}
\fancyfoot[C]{\thepage}
\renewcommand{\headrulewidth}{0.4pt}
\renewcommand{\footrulewidth}{0pt}
\fancypagestyle{plain}{%
  \fancyhf{}%
  \fancyfoot[C]{\thepage}%
  \renewcommand{\headrulewidth}{0pt}%
}

% ========================================================================
%  FRONT MATTER
% ========================================================================
\begin{frontmatter}

\title{Shariah Computational Financial Engineering I ---\\
Reconstructing Retail FX and CFD Market Access:\\
A Structural Incompatibility Analysis and\\
First-Principles Alternative Architecture}

\author[usyd]{Muhammad Taha Asif%
  \,\orcidlink{0009-0003-5277-5908}\corref{cor1}}
\ead[url]{https://orcid.org/0009-0003-5277-5908}
\cortext[cor1]{The author is an undergraduate student. This paper
  is independent research: it is not affiliated with, funded by,
  or supervised through any University of Sydney research
  programme, and the views expressed do not necessarily reflect
  those of the University or its faculty.}

\address[usyd]{School of Mathematics and Statistics,
  The University of Sydney, Camperdown NSW 2006, Australia}

% ---------- Abstract ----------
\begin{abstract}
Retail contracts for difference (CFDs) and leveraged foreign exchange
(FX) trading constitute one of the fastest-growing segments of global
retail market access, with daily notional turnover exceeding several
trillion United States dollars across major brokerages. Despite the
size of the observant-Muslim retail population, no rigorous mathematical
treatment of the compatibility between the standard CFD/FX architecture
and Islamic contract law currently exists in the peer-reviewed
literature. Drawing on primary-source observations from a Capital
Markets and Trading Analyst residency at Bacera Co Pty Ltd, this paper
formalises the CFD payoff mechanism, the Tom-Next overnight financing
structure, and the leverage--margin dynamics as a coupled stochastic
system driven by a Bates (1996) jump-diffusion with stochastic
volatility. Within this formal model we document five structural
incompatibilities or concerns between the standard retail CFD and the
selected AAOIFI/OIC constraint framework: an interest-rate-linked
financing mechanism (\emph{rib\={a}} concern), a synthetic contingent
payoff without ownership or delivery (\emph{gharar} concern), a
counterparty-reciprocity structure under B-book execution
(\emph{maysir} concern), non-fit with the classical nominate contract
categories including absence of \emph{qab\d{d}}, and a
reciprocal-deferred-contingent obligation structure
(\emph{bay' al-k\={a}li' bi'l-k\={a}li'} concern). We further show
that the ``swap-free Islamic account'' variant marketed by conventional
brokers addresses at most the first of these concerns and only in
part. We then construct, subject to the same constraint framework,
three candidate alternative market-access architectures: (i) a
\emph{wa'd}-based FX exposure vehicle, whose scholarly permissibility
depends on whether the double-\emph{wa'd} structure is admitted under
the recognised authorities (this paper takes the AAOIFI restrictive
reading and records the divergence); (ii) a
commodity-\emph{mur\={a}ba\d{h}a} financing overlay, whose organised
form is impermissible under OIC Resolution 179; and (iii) a
\emph{shar\={\i}kat al-'in\={a}n} pooled-equity structure that is
formally consistent with the operational design constraints adopted
in this paper, subject in every case to review by a qualified Shariah
supervisory board. For each we derive the payoff dynamics and holding
cost under the same underlying process. A simulation-based
quantitative comparison, under illustrative Bates parameterisations
consistent with published FX and equity implied-volatility studies,
quantifies the fee-structure differential, tail-risk properties, and
margin-call trigger distributions of each architecture against a
conventional leveraged CFD baseline. Because the \emph{shar\={\i}kat}
architecture cannot be leveraged within the joint-ownership
construction adopted, the comparison is properly understood as a
quantification of the economic cost of the leverage cap imposed by
that construction: the compliant architecture forgoes proportional
leveraged upside and leveraged downside relative to a CFD of matched
client capital, while removing the structural concerns identified.
The paper does not issue fiqh rulings; it formalises the mathematical
structure and documents its consistency, or otherwise, with the
adopted constraint framework. Ultimate compliance certification
remains the province of qualified Shariah supervisory boards. An
open-source Python implementation is released with the paper,
together with a full empirical calibration pipeline outlined as
future work.
\end{abstract}

% ---------- Keywords ----------
\begin{keyword}
Islamic finance \sep contracts for difference \sep foreign exchange \sep
stochastic modelling \sep Bates model \sep AAOIFI \sep
\emph{shar\={\i}kat al-'in\={a}n} \sep \emph{wa'd} \sep
market microstructure \sep financial engineering

\JEL G13 \sep G14 \sep G23 \sep Z12
\end{keyword}

\end{frontmatter}

% ========================================================================
%  SECTION 1 --- INTRODUCTION
% ========================================================================
\section{Introduction}
\label{sec:introduction}

\subsection{The retail leverage market and its participants}
\label{sec:motivation}

The retail contracts-for-difference (CFD) and leveraged foreign
exchange (FX) segment has grown from an institutional niche in the
early 2000s into one of the largest retail-facing derivative
markets in the world. The Bank for International Settlements'
2025 Triennial Central Bank Survey reports global daily
over-the-counter FX turnover of approximately
USD~9.6~trillion in April 2025 \citep{BIS2025}; retail leveraged
CFD activity is a much smaller subset of this aggregate but is a
segment growing more rapidly than the OTC total. Regulators in
the European Union, United Kingdom, and Australia have documented
that a persistently large fraction, between roughly 74 and 89
percent, depending on jurisdiction and reporting period, of
retail CFD accounts lose money over a given year
\citep{ESMA2018,ASIC2020,FCA2022}. Despite these outcomes, retail CFD
participation continues to expand, driven by low account minimums,
smartphone-native execution platforms, and aggressive marketing.

A material and growing subset of this retail participant base is
observant Muslim. There are approximately 2.0~billion Muslims globally,
of whom a rising share are working-age and financially active
\citep{Pew2015}. Refinitiv's \emph{Islamic Finance Development Report
2023} estimates global Islamic finance assets at approximately
USD~4.5~trillion, projected to reach USD~6.7~trillion by 2027
\citep{Refinitiv2023}. Yet the overwhelming majority of these assets
sit in institutional vehicles, Islamic banking, \emph{\d{s}uk\={u}k},
\emph{tak\={a}ful}, and Shariah-compliant real estate, rather than
in retail market-access products for short-horizon leveraged synthetic
exposure of the kind associated with CFDs. Muslim retail participants
seeking such exposure currently face limited options: participate in
conventional CFD/FX markets under some implicit compromise of
religious observance, use a ``swap-free Islamic account'' variant
marketed by conventional brokers (which is examined and shown
structurally insufficient in
Section~\ref{sec:shariah-analysis}), or forgo leveraged synthetic
exposure and confine activity to permissible spot vehicles.
Alternatives such as Shariah-compliant equity funds, spot equity
investment under AAOIFI screening, permissible spot commodity
transactions, and \emph{\d{s}uk\={u}k} are available for other
investment purposes but do not substitute for the specific
short-horizon leveraged synthetic exposure use case.

\subsection{The scholarly gap}
\label{sec:gap}

The academic literature on Islamic finance is extensive on the topics
of \emph{\d{s}uk\={u}k} pricing \citep{Cakir2007,Jobst2007},
\emph{mur\={a}ba\d{h}a} contract structure \citep{Usmani2002}, Islamic
equity screening \citep{Derigs2008,Kunhibava2010}, and
\emph{tak\={a}ful} risk-sharing \citep{Iqbal2011}. It is comparatively
silent on the intersection of retail derivatives, market
microstructure, and Islamic contract law. The closest prior work
addresses institutional Islamic derivatives
\citep{DeLorenzo2007,Bacha1999} and Islamic hedging frameworks based on
\emph{wa'd} promises \citep{Bacha2013}, but does not formally treat the
retail CFD architecture, the Tom-Next financing mechanism, or the
specific structural violations they generate.

A structured literature search was conducted across Scopus, Web
of Science, Google Scholar, SSRN, and the Islamic Financial
Services Board and ISRA repositories, using combinations of the
terms \{``Contract for Difference'', ``CFD'', ``leveraged retail
FX'', ``Tom-Next'', ``rollover swap''\} with \{``Shariah'',
``Islamic finance'', ``AAOIFI'', ``\emph{ribā}'', ``\emph{gharar}'',
``\emph{maysir}'', ``\emph{qab\d{d}}''\}, covering publications
indexed through January 2026. No study matching all three of the
following criteria was identified: (i) formalising the retail
CFD contract as a fiqh-constrained stochastic system with explicit
price-process, financing, and margin dynamics; (ii) proving the
resulting structural concerns from primary fiqh sources at the
level of formal propositions; and (iii) constructing and
quantitatively comparing first-principles alternative
architectures under the constraint set implied by AAOIFI Shariah
Standards. This paper addresses that gap. The novelty claim is
narrowed accordingly: it is an assertion about the peer-reviewed
literature indexed in the named databases through the stated
date, not a universal claim about all extant scholarship.

\subsection{Research questions}
\label{sec:rqs}

The paper is organised around three research questions:
\begin{itemize}
  \item[\textbf{RQ1.}] What structural conditions must a leveraged
  retail market-access contract satisfy to be candidates for review
  as consistent with the requirements of AAOIFI Shariah Standards
  No.~1 (Trading in Currencies), No.~8 (\emph{Mur\={a}ba\d{h}a}),
  No.~12 (\emph{Shar\={\i}kah}), and No.~18 (\emph{Qab\d{d}})?
  \item[\textbf{RQ2.}] Under what conditions, if any, can these
  constraints be simultaneously satisfied for the economic use cases
  currently served by retail CFDs, namely short-horizon directional
  exposure, hedging of underlying commercial risk, and cash-efficient
  synthetic exposure?
  \item[\textbf{RQ3.}] What is the economic cost differential between a
  Shariah-compliant alternative architecture and a conventional
  leveraged CFD, and how does this differential vary with holding
  horizon, underlying volatility, and leverage regime?
\end{itemize}

\subsection{Contributions}
\label{sec:contributions}

The paper makes five contributions.
\begin{enumerate}
  \item A formal mathematical characterisation of the standard CFD/FX
  architecture as a stochastic contract system, including the
  price-process, Tom-Next financing rate, leverage--margin dynamics,
  and broker revenue decomposition
  (Section~\ref{sec:cfd-mechanics}).
  \item Formal propositions documenting five structural
  incompatibilities or concerns between the standard CFD and the
  adopted AAOIFI/OIC constraint framework, together with a proof
  that swap-free Islamic accounts address at most one of these and
  only in part (Section~\ref{sec:shariah-analysis}).
  \item Three first-principles alternative architectures,
  \emph{wa'd}-based FX exposure, commodity-\emph{mur\={a}ba\d{h}a}
  financing overlay, and \emph{shar\={\i}kat al-'in\={a}n}
  pooled-equity structure, with explicitly derived payoff dynamics
  under a Bates jump-diffusion price process
  (Section~\ref{sec:alternatives}).
  \item A Monte Carlo comparison of cost-of-carry, tail-risk, and
  margin-call trigger distributions across architectures under
  illustrative Bates parameterisations, honestly framed as a
  quantification of the compliance-cost of the Shariah-imposed
  leverage cap rather than a claim of intrinsic economic superiority
  of one architecture over another
  (Section~\ref{sec:simulation}).
  \item An open-source Python reference implementation, released with
  the paper under a permissive licence, enabling replication and
  extension.
\end{enumerate}

\subsection{Positioning and limitations}
\label{sec:positioning}

This paper is a work of mathematical financial engineering informed by
primary fiqh sources. It is not a work of fiqh. The author is a
mathematics and statistics student, not a qualified Shariah scholar.
Every claim of structural incompatibility in
Section~\ref{sec:shariah-analysis} is derived by citation to a specific
AAOIFI Shariah Standard, a resolution of the International Islamic Fiqh
Academy of the Organisation of Islamic Cooperation (OIC), or a widely
accepted classical or contemporary fiqh source, and is stated as a
formal proposition subject to that citation. The alternative
architectures proposed in Section~\ref{sec:alternatives} are candidate
structures whose ultimate certification for practical use remains the
province of qualified Shariah supervisory boards. Where scholarly
opinion differs between the recognised schools, the paper states the
divergence explicitly and does not attempt to resolve it.

The paper is further limited to the retail CFD/FX segment.
Institutional Islamic derivatives, \emph{\d{s}uk\={u}k},
\emph{tak\={a}ful}, and Islamic banking are treated only insofar as
they inform the market-access architectures under consideration. The
Monte Carlo analysis in Section~\ref{sec:simulation} operates under
illustrative Bates parameterisations rather than a live empirical
calibration; a full calibration pipeline against real tick data is
outlined in Appendix~\ref{app:data} and deferred to a subsequent
paper in this series. Extension to order-book microstructure,
adverse selection, and regime-switching dynamics is similarly
deferred.

\subsection{Structure of the paper}
\label{sec:structure}

Section~\ref{sec:literature} surveys the relevant Islamic finance and
market-microstructure literatures and identifies the specific gap this
paper addresses. Section~\ref{sec:cfd-mechanics} formalises the
standard CFD/FX architecture and its stochastic mechanics.
Section~\ref{sec:shariah-analysis} states and proves the five
structural incompatibilities or concerns with the adopted
AAOIFI/OIC constraint framework.
Section~\ref{sec:alternatives} constructs three candidate alternative
architectures from first principles. Section~\ref{sec:simulation}
presents the Monte Carlo comparison. Section~\ref{sec:discussion}
discusses limitations, regulatory and implementation considerations,
and open questions. Section~\ref{sec:conclusion} concludes and
identifies directions for subsequent papers in this series.

% ========================================================================
%  Main body --- Sections 2 to 8
% ========================================================================
\section{Literature Review}
\label{sec:literature}

This section surveys the three literatures on which the paper draws:
(i) Islamic commercial law and its foundational nominate contracts,
(ii) the extant scholarly and academic work on Islamic derivatives and
hedging, and (iii) mainstream retail market-microstructure and
trader-behaviour research. The three streams are treated in turn and
the specific gap the paper addresses is stated formally in
Section~\ref{sec:lit-gap}.

\subsection{Islamic commercial law and the foundational contracts}
\label{sec:lit-fiqh}

Islamic commercial law (\emph{fiqh al-mu'\={a}mal\={a}t}) is derived
from the four canonical sources of the Shariah, the Qur'\={a}n, the
Sunnah of the Prophet Mu\d{h}ammad, \emph{ijm\={a}'} (juristic
consensus), and \emph{qiy\={a}s} (analogical reasoning), and codified
across the received bodies of the four Sunn\={\i} schools
(\d{H}anaf\={\i}, M\={a}lik\={\i}, Sh\={a}fi'\={\i}, \d{H}anbal\={\i})
and the Ja'far\={\i} school \citep{Kamali2000,Vogel1998}. Contemporary
reference works synthesise this corpus into a manageable set of
nominate contracts that recur across modern Islamic finance practice
\citep{Usmani2002,ElGamal2006,Ayub2007,Iqbal2011}.

The nominate contracts most relevant to the market-access architectures
considered in this paper are the following.

\begin{itemize}
    \item \emph{Bay'} (sale), the exchange of one item of property
    for another with mutual consent; a valid \emph{bay'} requires
    \emph{m\={a}l} (property of recognised value), \emph{qab\d{d}}
    (constructive or physical possession by the seller prior to sale;
    see AAOIFI Shariah Standard No.~18), and clarity as to
    subject-matter, quantity and price \citep{AAOIFI2015}.
    \item \emph{Salam}: a forward sale in which the price is paid in
    full at contract inception against deferred delivery of a fungible
    commodity of specified quality and quantity.
    \item \emph{Isti\d{s}n\={a}'}, a commission-to-manufacture
    contract for a specified good to be produced or built to
    specification.
    \item \emph{Mur\={a}ba\d{h}a}, a cost-plus disclosed-markup
    sale, widely deployed as the workhorse financing structure of
    Islamic banking (AAOIFI Standard No.~8; \citealp{Usmani2002}).
    \item \emph{Mush\={a}raka}, a joint-venture partnership with
    proportional profit-and-loss sharing (AAOIFI Standard No.~12).
    \item \emph{Shar\={\i}kat al-'in\={a}n}, a subclass of
    \emph{mush\={a}raka} in which each partner contributes capital and,
    optionally, labour; profit sharing is by pre-agreed ratio, but loss
    sharing must be strictly proportional to capital contribution.
    \item \emph{Mu\d{d}\={a}raba}, a silent-partnership contract in
    which one party (\emph{rabb al-m\={a}l}) provides capital and the
    other (\emph{mu\d{d}\={a}rib}) provides labour and management.
    \item \emph{Ij\={a}ra}, a lease of the usufruct of an asset with
    ownership risk retained by the lessor.
    \item \emph{Wak\={a}la}, an agency contract, in which one party
    acts for another for a defined fee (\emph{ujra}).
    \item \emph{Wa'd}: a unilateral promise; classically treated as
    a moral rather than legal obligation, but binding on the promisor
    under contemporary AAOIFI and Bank Negara Malaysia rulings when
    structured as a formal, one-sided undertaking.
\end{itemize}

Every valid contract in this framework must simultaneously exclude
three prohibited features: \emph{rib\={a}} (usury and, by extension,
contractual interest), \emph{gharar} (excessive uncertainty as to
subject-matter, price, or execution), and \emph{maysir} (gambling,
understood as a zero-sum transaction in which wealth transfers by pure
chance). The prohibitions are grounded in explicit Qur'\={a}nic
verses, 2:275 for \emph{rib\={a}}, 5:90 for \emph{maysir}, and in
Prophetic \d{h}ad\={\i}th for \emph{gharar}, and are treated as absolute
constraints rather than as variables in a trade-off
\citep{Chapra1985,Warde2010}. Each prohibition admits formal
characterisation as a constraint on the parameters of a contract; this
observation is exploited in Section~\ref{sec:shariah-analysis} to state
structural violations of the standard CFD as formal propositions.

The consolidation of contemporary practice around a standardised set of
contract types has been driven principally by the Accounting and
Auditing Organization for Islamic Financial Institutions (AAOIFI),
whose \emph{Shari'ah Standards} \citep{AAOIFI2015} are adopted or
referenced in the regulatory frameworks of Bahrain, the United Arab
Emirates, Qatar, Oman, Jordan, Sudan, Syria, Nigeria, and, in modified
form, Malaysia and Indonesia; and by the resolutions of the
International Islamic Fiqh Academy of the Organisation of Islamic
Cooperation \citep{OICFiqh63,OICFiqh179}. This paper adopts AAOIFI
Standards as its reference framework, on the ground that they
represent the most widely subscribed international codification of
Islamic commercial law and are cross-referenced in nearly every
peer-reviewed article on the subject.

\subsection{The Islamic derivatives and hedging debate}
\label{sec:lit-derivatives}

The Shariah status of derivative instruments is one of the most
actively contested subjects in contemporary Islamic finance
scholarship. The debate divides broadly along two axes: (i) whether
risk-transfer instruments as a class are compatible with the Islamic
requirement that commercial gain be linked to real economic activity,
and (ii) whether specific instruments, futures, options, swaps,
contracts for difference, can be restructured so as to satisfy formal
fiqh constraints.

The \emph{restrictive} position, associated with \citet{Usmani2002} and
with the majority of scholars in the Gulf Cooperation Council (GCC),
holds that conventional derivatives violate \emph{gharar} through the
uncertainty of settlement embedded in their cash-settled payoff,
violate \emph{maysir} through their zero-sum counterparty structure,
and, in the case of leveraged products, violate the
\emph{rib\={a}} prohibition through their financing overlay. On this
view, most conventional derivatives cannot be Islamised through
structural modification: only fundamentally different instruments
(\emph{salam}, \emph{isti\d{s}n\={a}'}, and equity-based partnership
structures) satisfy the constraints. \citet{Kamali2000} provides the
most sustained scholarly defence of the restrictive position with
respect to futures and options specifically.

The \emph{permissive} position, associated principally with Bank Negara
Malaysia's Shariah Advisory Council and with a subset of scholars in
Malaysia, Indonesia, and the ISRA (International Shari'ah Research
Academy) tradition, holds that certain derivative economic functions,
hedging of underlying commercial exposure, transfer of specific
risks, may be permissibly discharged through \emph{wa'd}-based
structures and combinations of nominate contracts
\citep{Bacha2013,Dusuki2011,ISRA2016}. The Islamic Profit Rate Swap and
the Islamic FX Forward, both developed and deployed by Malaysian and
Gulf Islamic banks, are the two most prominent instruments to have
emerged from this school.

AAOIFI sits closer to the restrictive pole. Its Shariah Standard No.~1
explicitly prohibits currency forwards and options; Standard No.~20
restricts organised commodity markets; and Standard No.~27 governs
indices. The International Islamic Fiqh Academy of the OIC has issued
Resolution No.~63 (1/7) ruling against conventional forwards and
options, and Resolution No.~179 (19/5) ruling against organised
\emph{tawarruq} \citep{OICFiqh63,OICFiqh179}. Where a candidate
architecture proposed in this paper is more permissive on any dimension
than AAOIFI, the divergence is stated explicitly and the affected
result is qualified accordingly.

Prior academic treatments of specific instrument classes include:
\citet{DeLorenzo2007} on the total-return swap and its use of
\emph{wa'd}-based ``Shariah conversion technology'';
\citet{Bacha1999} on the general permissibility architecture of
derivatives; \citet{Jobst2007} on Islamic securitisation, with a formal
treatment of \emph{\d{s}uk\={u}k} structuring under a real-asset
requirement; \citet{Kunhibava2010}, whose comparative analysis of
Shariah and conventional-law objections to derivatives is the closest
general survey in the peer-reviewed literature; \citet{AlSuwailem2006}
on hedging under Islamic constraints; and \citet{Uberoi2008} on the
practitioner mechanics of \emph{wa'd}-based hedging. On the
retail-facing use case of leveraged FX and CFDs specifically, however,
the peer-reviewed literature is essentially silent, this is the
gap identified formally in Section~\ref{sec:lit-gap}.

\subsection{Retail market microstructure and trader behaviour}
\label{sec:lit-microstructure}

A rigorous analysis of retail CFD/FX market access requires grounding
in the mainstream financial-economics literature on market
microstructure, retail trader behaviour, and leverage regulation. This
grounding is necessary for two reasons. First, it anchors the paper's
conclusions on cost-of-carry, execution asymmetry, and margin-call
dynamics in an established empirical evidence base rather than in fiqh
reasoning alone. Second, it locates the paper within the mainstream
finance conversation, so that its findings are legible, and
falsifiable, to reviewers and readers outside the Islamic finance
subfield.

The foundational microstructure literature on informed and uninformed
order flow, adverse selection, and dealer market-making descends from
\citet{Kyle1985} and \citet{GlostenMilgrom1985}, with the standard
practitioner treatment given in \citet{Harris2003} and the theoretical
consolidation in \citet{OHara1995}. The retail-broker-specific
extension of this framework, developed over the last decade around the
economics of ``payment for order flow'' and the ``B-book'' business
model, is surveyed in \citet{Menkveld2016}. Together these works
provide the theoretical basis for the broker-revenue decomposition
introduced in Section~\ref{sec:cfd-mechanics}, in which a fraction
$\beta \in [0,1]$ of client order flow is internalised by the broker
(B-book execution) rather than externalised to an exchange or
liquidity provider (A-book execution), with the implication that
broker profit is, at least in part, the algebraic mirror of client
loss.

The retail-trader behaviour literature is even sharper.
\citet{BarberOdean2000, BarberOdean2013} established, using
account-level data on tens of thousands of retail brokerage clients,
that individual investors as a class systematically underperform
passive benchmarks net of transaction costs. \citet{GrinblattKeloharju2000}
confirm the pattern for the Finnish retail population.
\citet{Barber2014Taiwan}, using account-level data on the near-universe
of Taiwanese day traders, find that fewer than one per cent of day
traders earn returns sufficient to cover transaction costs over any
sustained period, and that the top and bottom performers exhibit
substantial persistence. \citet{KelleyTetlock2013} find that retail
order flow is predominantly liquidity-motivated and, on aggregate,
adversely selected. The consistency of these findings across markets
and decades is difficult to overstate.

Leverage-specific regulation has generated its own body of evidence.
\citet{HeimerSimsek2019} evaluate the 2010 U.S. Commodity Futures
Trading Commission limit on retail FX leverage, from 100:1 to 50:1,
and find that the intervention meaningfully reduced retail losses
without materially reducing participation. The European Securities and
Markets Authority intervention against retail CFDs
\citep{ESMA2018}, the Financial Conduct Authority policy statement
PS19/18 \citep{FCA2022}, and the Australian Securities and Investments
Commission Product Intervention Order \citep{ASIC2020} each report
retail CFD account loss rates in the 74--89 per cent range and impose
leverage caps, margin-close-out requirements, and negative-balance
protection accordingly. These findings are treated in this paper as
external constraints on the design space for any candidate retail
market-access architecture, Shariah-compliant or otherwise: a
compliant architecture that leaves retail participants materially worse
off in expectation than the conventional benchmark is not a socially
defensible design, whatever its jurisprudential merits.

\subsection{The identified gap}
\label{sec:lit-gap}

The three literatures surveyed above have not, to the best of the
author's knowledge, been synthesised in the peer-reviewed record. The
Islamic commercial-law and derivatives literatures address the
contractual and jurisprudential dimensions but do not treat the retail
CFD/FX architecture as a formal stochastic system, do not derive its
specific structural violations at the level of formal propositions,
and do not construct alternative architectures whose payoff and
cost-of-carry dynamics are explicitly modelled under a common
underlying price process. The market-microstructure and
retail-behaviour literatures treat the empirical mechanics of retail
CFD and FX markets rigorously but make no reference to Islamic
contract law, or to the substantial observant-Muslim retail
participant base that intersects those markets. Practitioner treatments
of ``Islamic FX'' and ``swap-free accounts'', where they exist in
the trade press and in broker marketing material, are unsourced,
and are shown in Section~\ref{sec:shariah-analysis} to be structurally
insufficient.

The specific gap this paper addresses is therefore stated as follows:
\emph{no prior peer-reviewed work treats the retail CFD/FX architecture
as a formal fiqh-constrained stochastic system, proves its structural
violations at the level of formal propositions grounded in AAOIFI
Shariah Standards and OIC Fiqh Academy resolutions, or constructs and
quantitatively compares first-principles alternative architectures
under the resulting constraint set.} The remainder of the paper
addresses this gap directly:
Section~\ref{sec:cfd-mechanics} formalises the CFD/FX architecture;
Section~\ref{sec:shariah-analysis} documents its five structural
violations; Section~\ref{sec:alternatives} constructs three candidate
alternative architectures; and Section~\ref{sec:simulation}
compares them quantitatively.

\section{Formal Characterisation of the CFD/FX Architecture}
\label{sec:cfd-mechanics}

This section formalises the standard retail CFD/FX architecture as a
coupled stochastic system, comprising: an underlying price process
(Section~\ref{sec:price-process}), a contract-payoff mapping
(Section~\ref{sec:cfd-payoff}), a financing mechanism
(Section~\ref{sec:tomnext}), a margin and leverage dynamic
(Section~\ref{sec:margin}), and a broker-revenue decomposition
(Section~\ref{sec:broker}). The section concludes with primary-source
institutional context from a four-week Capital Markets and Trading
Analyst residency at Bacera Co Pty Ltd
(Section~\ref{sec:primary-source}). The characterisation established
here is the object of the Shariah analysis in
Section~\ref{sec:shariah-analysis}.

\subsection{Preliminaries and the underlying price process}
\label{sec:price-process}

Throughout the paper we work on a filtered probability space
$(\Omega, \FF, \{\FF_t\}_{t \geq 0}, \PP)$ satisfying the usual
conditions of right-continuity and $\PP$-completeness. All stochastic
processes are adapted to $\{\FF_t\}_{t \geq 0}$. The measure $\PP$
denotes the physical (real-world) probability measure; risk-neutral
counterparts, where required, are denoted $\PP^{*}$ or $\mathbb{Q}$.

Let $(S_t)_{t \geq 0}$ denote the spot price of the underlying asset
at time~$t$. For an FX pair CCY1/CCY2 (\eg\ EUR/USD), $S_t$ denotes the
price of one unit of the base currency CCY1 quoted in units of the quote
currency CCY2. For an equity or index underlying, $S_t$ denotes the
spot price in the reference currency.

\begin{assumption}[Underlying price process]
\label{ass:bates}
The underlying price process $(S_t)_{t \geq 0}$ follows the
\citet{Bates1996} jump-diffusion model with stochastic volatility. The
dynamics under $\PP$ are given by the coupled system
\begin{align}
    \frac{\dd S_t}{S_{t^{-}}}
    &= \bigl( \mu - \lambda \kappa \bigr) \, \dd t
       + \sqrt{v_t}\, \dd W^{S}_{t}
       + \bigl( J_t - 1 \bigr) \, \dd N_t,
    \label{eq:bates-S} \\[4pt]
    \dd v_t
    &= \kappa_v \bigl( \theta - v_t \bigr) \, \dd t
       + \sigma_v \sqrt{v_t}\, \dd W^{v}_{t},
    \label{eq:heston-v}
\end{align}
where:
\begin{itemize}
    \item $\mu \in \RR$ is the drift under $\PP$; under an equivalent
    martingale measure $\PP^{*}$ the drift is replaced by $r - q$,
    with $r$ the domestic risk-free rate and $q$ the dividend yield
    (for equity) or foreign risk-free rate (for FX);
    \item $(v_t)_{t \geq 0}$ is the instantaneous variance, following
    the \citet{Heston1993} square-root process with mean-reversion
    speed $\kappa_v > 0$, long-run variance $\theta > 0$, and
    volatility-of-volatility $\sigma_v > 0$;
    \item $W^{S}$ and $W^{v}$ are standard Brownian motions with
    correlation $\dd \langle W^{S}, W^{v} \rangle_{t} = \rho \, \dd t$,
    $\rho \in [-1, 1]$;
    \item $N_t$ is a Poisson process with intensity $\lambda > 0$,
    independent of $W^{S}$ and $W^{v}$;
    \item $(J_t)_{t \geq 0}$ is a sequence of i.i.d.\ strictly positive
    relative jump sizes, with $\ln J_t \sim \mathcal{N}(\mu_J, \sigma_J^2)$;
    \item $\kappa \coloneqq \EE[J_t - 1]
    = \exp\!\bigl(\mu_J + \tfrac{1}{2}\sigma_J^{2}\bigr) - 1$
    is the mean relative jump size, subtracted from the drift so that
    the compensated jump term is a $\PP$-martingale.
\end{itemize}
The strict Feller condition $2\kappa_v \theta > \sigma_v^{2}$ is
imposed throughout to guarantee that the variance process $v_t$
almost surely does not attain zero (\citealp{Heston1993}); the
boundary case $2\kappa_v \theta = \sigma_v^{2}$ is excluded.
\end{assumption}

\begin{remark}[Measure choice for the Shariah and simulation analysis]
\label{rem:measure-choice}
The Shariah structural analysis (Section~\ref{sec:shariah-analysis})
depends only on the pathwise almost-sure properties of $(S_t)$, which
are invariant under equivalent measure change and therefore hold under
both $\PP$ and any equivalent $\PP^{*}$. The Monte Carlo simulations
(Section~\ref{sec:simulation}) are conducted under the physical measure
$\PP$ throughout, using the parameters of Table~\ref{tab:calibration},
to reproduce empirically observed loss distributions, margin-call
frequencies and Sharpe ratios of the sort a retail client would face.
Where a risk-neutral valuation is required
(Proposition~\ref{prop:wad-payoff}), we invoke the standard equivalent
martingale measure for the incomplete jump-diffusion market, noting
that under $\PP^{*}$ the jump intensity $\lambda$ and the jump-size
law of $\ln J_t$ may differ from their $\PP$-counterparts; the choice
of $\PP^{*}$ within the family of equivalent measures is not unique in
an incomplete market (\citealp{Cont2004}). No claim in the paper turns
on a specific such choice.
\end{remark}

\begin{remark}[Choice of price process]
\label{rem:model-choice}
The Bates model is selected over the plain \citet{BlackScholes1973}
geometric Brownian motion for three reasons load-bearing for the
analysis of Sections~\ref{sec:shariah-analysis}--\ref{sec:simulation}.
First, retail CFD/FX markets display well-documented jump behaviour
around scheduled macroeconomic releases (US non-farm payrolls, FOMC
and ECB decisions, RBA cash-rate announcements) and around unscheduled
events (central bank interventions, flash crashes such as the CHF
un-pegging of 15 January 2015); these jumps drive the clustering of
margin-call trigger events, which the diffusion-only benchmark cannot
reproduce (\citealp{Merton1976}). Second, the leverage effect,
empirically, $\rho < 0$ for equity underlyings, is quantitatively
material for tail-risk assessment and is captured by the correlated
Brownian structure. Third, stochastic volatility is required to
reproduce the implied volatility smiles observed in traded FX and
equity option markets against which the model would be calibrated
(\citealp{Heston1993,Bjork2009}).
\end{remark}

\subsection{The contract for difference and its payoff process}
\label{sec:cfd-payoff}

We now formalise the CFD contract as a mapping from the underlying
price path to a cash-settled payoff to the client, together with a
financing overlay to be specified in Section~\ref{sec:tomnext}.

\begin{definition}[Contract for Difference]
\label{def:cfd}
A \emph{contract for difference} (CFD) on the underlying $S$ is a
bilateral cash-settled derivative contract between a client (retail
participant) and a broker (counterparty), specified by the tuple
$\bigl( t_0, t_1, \delta, N, L, s, m \bigr)$ where:
\begin{itemize}
    \item $t_0 \in [0, T]$ is the entry time,
    \item $t_1 \in (t_0, T]$ is the exit time (chosen by the client
    or triggered exogenously by margin-call; see
    Section~\ref{sec:margin}),
    \item $\delta \in \{+1, -1\}$ is the directional indicator ($+1$
    for a long position, $-1$ for a short position),
    \item $N \in \RR_{+}$ is the notional size, measured in units of
    the underlying,
    \item $L \geq 1$ is the leverage ratio,
    \item $(s_t)_{t \geq 0}$ is the bid--ask spread process, and
    \item $(m_t)_{t \geq 0}$ is the broker's markup process over the
    interbank financing rate (see Section~\ref{sec:tomnext}).
\end{itemize}
At exit time $t_1$ the broker settles cash to the client in the
amount of the net payoff $\Pi_{\text{net}}(t_0, t_1)$ defined in
equation~\eqref{eq:cfd-net-payoff} below.
\end{definition}

\begin{definition}[Gross and net payoff]
\label{def:payoff}
The client's \emph{gross payoff} on a CFD held from $t_0$ to $t_1$ is
\begin{equation}
    \Pi_{\text{gross}}(t_0, t_1)
    = \delta \cdot N \cdot \bigl( S_{t_1} - S_{t_0} \bigr).
    \label{eq:cfd-gross-payoff}
\end{equation}
The corresponding \emph{net payoff}, accounting for accumulated
financing and transaction costs, is
\begin{equation}
    \Pi_{\text{net}}(t_0, t_1)
    = \delta \cdot N \cdot \bigl( S_{t_1} - S_{t_0} \bigr)
      - \int_{t_0}^{t_1} F_s \, \dd s
      - C(t_0, t_1),
    \label{eq:cfd-net-payoff}
\end{equation}
where $(F_s)_{s \geq 0}$ is the instantaneous financing rate derived in
Section~\ref{sec:tomnext}, and $C(t_0, t_1)$ is the aggregate
transaction cost from spreads and fixed commissions,
\begin{equation}
    C(t_0, t_1)
    = N \cdot \frac{s_{t_0} + s_{t_1}}{2}
      + c_{\text{fixed}},
    \label{eq:cfd-transaction-cost}
\end{equation}
with $c_{\text{fixed}} \geq 0$ a broker-specific per-trade commission
(often zero on retail FX; positive on single-name equity CFDs).
\end{definition}

\begin{remark}[Cash settlement and the absence of delivery]
\label{rem:no-delivery}
Critical to the fiqh analysis of Section~\ref{sec:shariah-analysis}: at
no point during $[t_0, t_1]$ does the client acquire ownership of, or
any legal or beneficial claim to, $N$ units of the underlying $S$.
The client's exposure is exclusively to the price difference
$S_{t_1} - S_{t_0}$. This distinguishes the CFD from (i) a spot
purchase of the underlying, which confers ownership; (ii) a forward or
futures contract requiring physical delivery; and (iii) an equity share
carrying voting and residual-claim rights. The absence of any
transfer of ownership is the empirical foundation of the \emph{qab\d{d}}
violation proved as Proposition~\ref{prop:qabd} in
Section~\ref{sec:shariah-analysis}.
\end{remark}

\subsection{The Tom-Next swap: formal derivation of the financing rate}
\label{sec:tomnext}

The overnight rollover fee, industry-termed the ``Tom-Next swap''
(from ``tomorrow-next-day''), is the mechanism by which brokers
charge (or credit) the client for holding a leveraged position across
the daily settlement cutoff, conventionally 5:00~PM New York time.
This is the empirical mechanism underlying the \emph{rib\={a}}
violation formalised in Section~\ref{sec:shariah-analysis}. We derive
it here from first principles under covered interest rate parity.

\subsubsection{Derivation for FX pairs}
\label{sec:tomnext-fx}

Consider a long CFD position of notional $N$ on the FX pair CCY1/CCY2,
held from settlement date $T_d$ to the next business settlement date
$T_{d+1}$. Under covered interest rate parity (\citealp{Bjork2009},
Ch.~24), the no-arbitrage forward exchange rate satisfies
\begin{equation}
    F(T_d, T_{d+1})
    = S_{T_d} \cdot
      \frac{1 + r^{\text{CCY2}}_{T_d} \tau}
           {1 + r^{\text{CCY1}}_{T_d} \tau},
    \label{eq:cip-forward}
\end{equation}
where $r^{\text{CCY}}_{T_d}$ is the overnight interbank lending rate
for currency CCY at $T_d$, SOFR for USD, ESTR for EUR, SONIA for
GBP, TONA for JPY, and the RBA cash rate for AUD, and
$\tau = 1/365$ is the actual/365 day-count fraction.

Rolling a spot position forward by one day requires the counterparty to
pay (or receive) the difference $F(T_d, T_{d+1}) - S_{T_d}$. For small
$\tau$, first-order Taylor expansion of \eqref{eq:cip-forward} gives
\begin{equation}
    F(T_d, T_{d+1}) - S_{T_d}
    = S_{T_d} \cdot
      \bigl( r^{\text{CCY2}}_{T_d} - r^{\text{CCY1}}_{T_d} \bigr) \tau
      + O(\tau^{2}).
    \label{eq:cip-expansion}
\end{equation}
The broker charges the client this differential plus a proprietary
markup $m_{T_d} \geq 0$. The Tom-Next swap fee for a long CFD position
held across the cutoff is therefore
\begin{equation}
    F^{\text{swap}}_{T_d}
    = N \cdot S_{T_d} \cdot
      \frac{\bigl( r^{\text{CCY2}}_{T_d} - r^{\text{CCY1}}_{T_d} \bigr)
             + m_{T_d}}{365}.
    \label{eq:tomnext-discrete}
\end{equation}
For a short position ($\delta = -1$), the sign of the interbank
differential reverses (since the client is now the funder of the base
currency), but the broker's markup is charged with a sign that always
disadvantages the client. Empirically, retail brokers post separate
long and short swap rates in which the markup component is applied
so that both sides of the market bear a non-negative broker margin
(see, e.g., MetaTrader~5 symbol specifications, which display distinct
\texttt{swap\_long} and \texttt{swap\_short} figures with the broker
markup embedded in each). We encode this sign convention explicitly
in the continuous-time definition below by separating the
CIP-implied differential (whose direction depends on $\delta$) from
the markup (which is always a cost to the client).

\begin{definition}[Instantaneous financing rate]
\label{def:instantaneous-financing}
Passing to the continuous-time limit, the instantaneous financing rate
$F_s$ that appears in the net payoff~\eqref{eq:cfd-net-payoff} is
\begin{equation}
    F_s
    = N \cdot S_s \cdot
      \Bigl[ \delta \bigl( r^{\text{CCY2}}_{s} - r^{\text{CCY1}}_{s} \bigr)
             + m_s \Bigr],
    \label{eq:financing-rate-continuous}
\end{equation}
where $r^{\text{CCY1}}_{s}$ and $r^{\text{CCY2}}_{s}$ are the
instantaneous overnight interbank rates for the base and quote
currencies respectively, $m_s \geq 0$ is the broker's continuous-time
markup process, and $\delta \in \{+1,-1\}$ is the position direction.
Under this sign convention, both the long-side and short-side
financing charges include a non-negative markup contribution,
consistent with observed retail broker pricing.
\end{definition}

\begin{proposition}[Interbank-rate measurability of the
    financing rate]
\label{prop:financing-measurable}
Under Definition~\ref{def:instantaneous-financing}, the
Tom-Next financing rate $F_s$ is measurable with respect to
the $\sigma$-algebra
$\mathcal{G}_s \coloneqq \sigma\bigl( S_s,\,
r^{\text{CCY1}}_{s},\, r^{\text{CCY2}}_{s},\, m_s \bigr)$
generated by the underlying price and the two interbank
overnight rates. In particular, conditional on
$\{S_s, m_s, \delta\}$, $F_s$ is an affine function of the interbank
rate differential $r^{\text{CCY2}}_{s} - r^{\text{CCY1}}_{s}$
with sensitivity:
\begin{equation}
    \frac{\partial F_s}{\partial
      (r^{\text{CCY2}}_{s} - r^{\text{CCY1}}_{s})}
    = \delta \cdot N \cdot S_s.
    \label{eq:financing-measurable}
\end{equation}
\end{proposition}

\begin{proof}
Direct from equation~\eqref{eq:financing-rate-continuous}. The
coefficients $\delta$, $N$, $S_s$, $m_s$ are $\FF_s$-measurable by
construction. The rates $r^{\text{CCY1}}_{s}$ and $r^{\text{CCY2}}_{s}$
are observable at time $s$ (published by the respective central banks
or interbank rate administrators). The affine-with-unit-sensitivity
property is immediate from differentiation of
\eqref{eq:financing-rate-continuous} with respect to the differential.
\qed
\end{proof}

\begin{remark}[The \emph{rib\={a}} structural implication, foreshadowed]
\label{rem:riba-foreshadow}
Proposition~\ref{prop:financing-measurable} formalises what is
empirically obvious to any practitioner: the CFD financing rate is,
by construction, a function of interbank interest rates. Under the
general \emph{rib\={a}} prohibition articulated in \citet{Usmani2002}
and \citet{Kamali2000}, namely the prohibition of a fixed increment
charged for the deferral of exchange, and taken together with
AAOIFI Shariah Standard No.~1 (on Trading in Currencies) and No.~8
(on \emph{Mur\={a}ba\d{h}a}), which respectively restrict deferred
currency transactions and the addition of stipulated increments to
credit-sale contracts, the Tom-Next financing structure is at
minimum a serious \emph{rib\={a}} concern under the majority
contemporary reading; the argument is developed formally in
Section~\ref{sec:shariah-analysis} and its full authority basis is
discussed there.
\end{remark}

\subsection{Margin, leverage, and margin-call dynamics}
\label{sec:margin}

Leverage is central to the retail CFD product: it is the feature that
converts small-notional retail capital into large-notional market
exposure, and it is the feature that drives the empirically observed
loss rates cited in Section~\ref{sec:motivation}. We formalise its
mechanics.

\subsubsection{Initial and maintenance margin}
\label{sec:margin-levels}

For a CFD position of notional value $N \cdot S_{t_0}$ opened at
leverage $L$, the client posts \emph{initial margin} equal to
\begin{equation}
    M_0 = \frac{N \cdot S_{t_0}}{L},
    \label{eq:initial-margin}
\end{equation}
which is the fraction $1/L$ of the notional value of the exposure. The
broker requires a \emph{maintenance margin} threshold
\begin{equation}
    M_{\text{maint}} = \alpha \cdot M_0,
    \qquad \alpha \in (0, 1],
    \label{eq:maintenance-margin}
\end{equation}
with $\alpha$ typically in $[0.25, 0.75]$ depending on regulatory
regime and asset class. The client's account equity at time $t$ is
\begin{equation}
    E_t = M_0 + \delta \cdot N \cdot \bigl( S_t - S_{t_0} \bigr)
          - \int_{t_0}^{t} F_s \, \dd s.
    \label{eq:equity}
\end{equation}

\subsubsection{The margin-call condition}
\label{sec:margin-call}

The margin call is triggered at the first stopping time $\tau_{\text{MC}}$
such that account equity falls at or below the maintenance threshold:
\begin{equation}
    \tau_{\text{MC}}
    \coloneqq \inf \bigl\{
      t > t_0 : E_t \leq M_{\text{maint}}
    \bigr\}.
    \label{eq:tau-mc-definition}
\end{equation}
On short holding horizons the financing term
$\int_{t_0}^{t} F_s \, \dd s$ is small relative to price innovations
and may be neglected in a leading-order analysis. Under this
approximation, the barrier condition~\eqref{eq:tau-mc-definition}
simplifies to
\begin{equation}
    \delta \cdot N \cdot \bigl( S_t - S_{t_0} \bigr)
    \leq -(1 - \alpha) M_0,
    \label{eq:mc-price-barrier}
\end{equation}
which for a long position ($\delta = +1$), combined with
\eqref{eq:initial-margin}, yields the explicit price barrier
\begin{equation}
    S_t \leq S_{t_0} \cdot \left( 1 - \frac{1 - \alpha}{L} \right).
    \label{eq:mc-explicit-barrier}
\end{equation}
For a short position the barrier lies symmetrically above $S_{t_0}$.

\begin{proposition}[Leverage sensitivity of margin-call probability]
\label{prop:leverage-sensitivity}
Suppose $S_t$ follows a geometric Brownian motion with constant
volatility $\sigma$ (Assumption~\ref{ass:bates} with $\lambda = 0$
and $v_t \equiv \sigma^{2}$). For a long CFD position with leverage
$L$ and holding horizon $h \coloneqq t_1 - t_0$, the margin-call
probability admits the approximation
\begin{equation}
    \PP\bigl( \tau_{\text{MC}} < t_1 \mid \FF_{t_0} \bigr)
    \approx 2 \, \Phi\!\left(
      -\frac{1 - \alpha}{L \, \sigma \sqrt{h}}
    \right),
    \label{eq:mc-probability}
\end{equation}
where $\Phi$ denotes the standard normal cumulative distribution
function. Consequently
\begin{equation}
    \frac{\partial}{\partial L}
    \PP\bigl( \tau_{\text{MC}} < t_1 \mid \FF_{t_0} \bigr) > 0,
    \qquad
    \PP\bigl( \tau_{\text{MC}} < t_1 \mid \FF_{t_0} \bigr)
    = 1 - \sqrt{\tfrac{2}{\pi}} \cdot
        \frac{1 - \alpha}{L \, \sigma \sqrt{h}}
      + O\!\bigl( L^{-3} \bigr),
    \label{eq:mc-leverage-scaling}
\end{equation}
so that
$\PP\bigl( \tau_{\text{MC}} < t_1 \mid \FF_{t_0} \bigr) \to 1$
as $L \to \infty$.
\end{proposition}

\begin{proof}
The result follows from the reflection principle applied to Brownian
motion with drift, using the standard first-passage-time density
(\citealp{Shreve2004}, \S3.7.1). Under geometric Brownian
motion the log-price barrier corresponding to
\eqref{eq:mc-explicit-barrier} is
\[
    b \coloneqq \ln\!\left( 1 - \frac{1 - \alpha}{L} \right)
    = -\frac{1 - \alpha}{L} + O(L^{-2}),
\]
and the first-passage probability for the downward barrier is
\[
    \PP(\tau_{\text{MC}} < h)
    = \Phi\!\left( \frac{b - \mu h}{\sigma \sqrt{h}} \right)
      + e^{2 \mu b / \sigma^{2}} \,
      \Phi\!\left( \frac{b + \mu h}{\sigma \sqrt{h}} \right).
\]
Here $\mu$ denotes the drift of the log-price process
$X_t \coloneqq \log(S_t/S_{t_0})$ (i.e.\ $\nu = \mu_S - \tfrac12 \sigma^{2}$
in the notation of Appendix~\ref{app:barrier}, where $\mu_S$
denotes the price drift). Setting $\mu = 0$,
the leading-order approximation of a negligible log-price
drift over the short horizons $h$ considered here, and
substituting the leading term of $b$ yields
\eqref{eq:mc-probability}. Writing
$x \coloneqq (1-\alpha)/(L\sigma\sqrt{h})$, the scaling
in~\eqref{eq:mc-leverage-scaling} follows from the Taylor
expansion $\Phi(-x) = \tfrac{1}{2} - x/\sqrt{2\pi} + O(x^3)$
valid as $x \to 0$; here $L \to \infty$ drives $x \to 0$ so this
is the appropriate expansion (the more familiar Mills-ratio
asymptotic $\Phi(-x) \sim \phi(x)/x$ applies as $x \to \infty$
and would be wrong in this regime).
The limit $\PP \to 1$ is intuitive: at very high leverage the
maintenance-margin barrier is arbitrarily close to $S_{t_0}$, so
any downward Brownian excursion within horizon $h$ triggers a
margin call. \qed
\end{proof}

\begin{remark}[Effect of jumps on the barrier-crossing probability]
\label{rem:jump-amplification}
The diffusion-only expression~\eqref{eq:mc-probability} is not,
in general, a strict lower bound on the true Bates margin-call
probability. Under the full Bates dynamics
(Assumption~\ref{ass:bates}) the log-jump-size distribution
$\ln J_t \sim \mathcal{N}(\mu_J, \sigma_J^2)$ admits both
$J_t > 1$ (upward jumps, which move the price away from the
downward barrier) and $J_t < 1$ (downward jumps, which move
towards or through it). The net effect of jump inclusion on
first-passage probability depends on the calibrated
$(\lambda, \mu_J, \sigma_J)$ and the compensator adjustment
$-\lambda\kappa$; it is a genuinely two-sided perturbation and
must be quantified numerically. The empirical simulation in
Section~\ref{sec:simulation} finds that under the illustrative
parameterisations adopted, the jump contribution is material at
holding horizons above approximately 24~hours. The
constant-intensity Poisson process used here does not model
event-time clustering (macroeconomic release windows,
central-bank announcements); the empirical literature documents
elevated jump activity around such events, and modelling
event-time clustering via a state-dependent intensity
$\lambda(t)$ is left to future work.
\end{remark}

\subsection{Broker revenue decomposition and B-book internalisation}
\label{sec:broker}

We decompose the broker's revenue from a single CFD position into
three components, following the market microstructure framework of
\citet{OHara1995}, \citet{Harris2003}, and \citet{Menkveld2016}. This
decomposition is essential for the \emph{maysir} analysis of
Section~\ref{sec:shariah-analysis}, which turns on the algebraic
relationship between broker P\&L and client P\&L.

\subsubsection{Spread revenue}

For a completed round-trip trade of notional $N$, the broker earns
\begin{equation}
    R^{\text{spread}}
    = N \cdot \frac{s_{t_0} + s_{t_1}}{2},
    \label{eq:r-spread}
\end{equation}
capturing half the bid--ask spread on entry and half on exit. This
component is deterministic conditional on execution and is independent
of the price path $\{S_t\}_{t \in [t_0, t_1]}$.

\subsubsection{Financing revenue}

Aggregate financing revenue over the holding period equals the
client's paid financing:
\begin{equation}
    R^{\text{swap}}
    = \int_{t_0}^{t_1} F_s \, \dd s,
    \label{eq:r-swap}
\end{equation}
with $F_s$ given by Definition~\ref{def:instantaneous-financing}. This
component grows linearly in holding time under stationary interbank
rates and constitutes a large fraction of broker revenue on positions
held across multiple settlement cutoffs.

\subsubsection{Internalisation P\&L and the A-book / B-book distinction}

Retail brokers operate on a spectrum of internalisation strategies.
Let $\beta \in [0, 1]$ denote the fraction of client notional that the
broker holds on its own book (``B-book'' execution) rather than
externalising to an exchange or third-party liquidity provider
(``A-book'' execution).
\begin{itemize}
    \item At $\beta = 0$ the broker externalises every client order and
    earns only spread and financing revenue; its P\&L is independent
    of the client's directional P\&L.
    \item At $\beta = 1$ the broker takes the opposite side of every
    client order and holds the resulting exposure. Its P\&L on that
    exposure is the negative of the client's gross P\&L.
    \item Intermediate $\beta$ corresponds to selective
    internalisation, in which the broker retains client flow assessed
    as unlikely to be informed on its own book and externalises flow
    assessed as informed; this is the standard practitioner
    description of retail broker order-management practice, and the
    associated adverse-selection framework in microstructure is
    surveyed in the market-maker inventory literature
    (\citealp{OHara1995}; \citealp{Harris2003}).
\end{itemize}
Formally, the internalisation P\&L component is
\begin{equation}
    R^{\text{internal}}
    = -\, \beta \cdot \Pi_{\text{gross}}(t_0, t_1),
    \label{eq:r-internal}
\end{equation}
so that the broker's total revenue from a single position is
\begin{equation}
    R^{\text{total}}
    = R^{\text{spread}} + R^{\text{swap}} + R^{\text{internal}}
    = \underbrace{
        N \cdot \tfrac{s_{t_0} + s_{t_1}}{2}
      }_{\text{spread}}
    + \underbrace{
        \int_{t_0}^{t_1} F_s \, \dd s
      }_{\text{financing}}
    - \underbrace{
        \beta \, \delta N (S_{t_1} - S_{t_0})
      }_{\text{internalisation}}.
    \label{eq:r-total}
\end{equation}

\begin{proposition}[Zero-sum structure under full B-book]
\label{prop:zero-sum}
Suppose the broker operates fully B-book, so that $\beta = 1$, and
consider the internalisation P\&L in isolation from spread and
financing components. Then the client's gross payoff and the broker's
internalisation P\&L sum to zero pathwise:
\begin{equation}
    \Pi_{\text{gross}}(t_0, t_1) + R^{\text{internal}} = 0
    \qquad \text{a.s.},
    \label{eq:zero-sum}
\end{equation}
and in particular
$\EE[\Pi_{\text{gross}}(t_0, t_1)] = -\EE[R^{\text{internal}}]$.
\end{proposition}

\begin{proof}
Substitute $\beta = 1$ into \eqref{eq:r-internal} and add
\eqref{eq:cfd-gross-payoff}: the sum vanishes identically, pathwise, on
$\Omega$. Expectation preservation is then immediate. \qed
\end{proof}

\begin{remark}[Empirical relevance of $\beta$]
\label{rem:beta-empirical}
Regulatory disclosure regimes in the European Union
(\citealp{ESMA2018}), the United Kingdom (\citealp{FCA2022}), and
Australia (\citealp{ASIC2020}) require retail CFD brokers to disclose
the fraction of client accounts that lose money over the reporting
period. The reported loss ranges (74--89\%, cited in
Section~\ref{sec:motivation}) are broadly consistent with
substantial B-book internalisation but do not identify $\beta$:
they measure client account outcomes, not the fraction of client
notional the broker warehouses on its own book, and external
counterparties may also profit systematically from adversely
selected retail flow. The paper does not estimate $\beta$ and
does not claim any industry-wide value.
Proposition~\ref{prop:maysir} below is therefore a conditional
structural result: wherever $\beta > 0$, the internalised
fraction has the pathwise direct-counterparty zero-sum property.
The strength of the \emph{maysir} concern in a specific
commercial implementation depends on the value of $\beta$ that
actually obtains there, which is a question of broker order
management and outside the scope of this paper.
\end{remark}

The zero-sum structure of Proposition~\ref{prop:zero-sum} is the
mathematical content of the \emph{maysir} violation formalised in
Section~\ref{sec:shariah-analysis} as Proposition~\ref{prop:maysir}.
The subject of Section~\ref{sec:alternatives} is precisely to
construct market-access architectures under which the analogous
revenue decomposition satisfies the Shariah constraint set, in
particular, structures under which $R^{\text{swap}}$ is not a function
of interbank rates and under which $R^{\text{internal}}$ is not the
pathwise negation of client P\&L.

\subsection{Primary-source institutional context}
\label{sec:primary-source}

The formal characterisation established in
Sections~\ref{sec:price-process}--\ref{sec:broker} draws on public
regulatory documentation, the academic literature surveyed in
Section~\ref{sec:literature}, and standard market-microstructure
references. It is corroborated by primary-source observations gathered
during the author's four-week residency as a Capital Markets and
Trading Analyst at Bacera Co Pty Ltd (Sydney, Australia), a global
CFD and FX brokerage regulated by the Australian Securities and
Investments Commission under Australian Financial Services Licence
328794. The observations recorded here are institutional context and
are not relied upon for any specific quantitative claim in the paper;
all such claims are grounded in the public sources cited above.

\paragraph{Product mix.}
The observed product suite comprised leveraged CFDs on major, minor,
and exotic FX pairs; spot metals (XAU/USD, XAG/USD); energy
(WTI, Brent); major equity indices (SPX, NDX, DAX, ASX~200, N225);
and a subset of individual equities across the US, UK, and Australian
markets. Client-facing leverage was capped in accordance with the ASIC
2020 product intervention order for accounts domiciled in Australia:
up to 30:1 on major FX pairs, 20:1 on minor FX and gold, 10:1 on
non-major indices and commodities, 2:1 on cryptocurrencies, and 5:1 on
individual equities.

\paragraph{Financing mechanics.}
Tom-Next swaps were applied daily at 5:00~PM New York time (which
corresponds to 7:00~AM the following calendar day in Sydney AEDT
during Australian summer). A triple-swap charge was applied on
Wednesdays, following the T+2 spot settlement convention that shifts
the effective rollover date across the weekend on that day of the
week. Published swap rates were consistent in sign and order of
magnitude with the interbank differential plus markup structure
derived formally in equation~\eqref{eq:tomnext-discrete} and were
observed to adjust in response to policy actions by the Federal
Reserve, the European Central Bank, and the Reserve Bank of Australia
during the residency period.

\paragraph{Client outcomes.}
Consistent with the published ASIC disclosures cited in
Section~\ref{sec:motivation}, the majority of retail accounts observed
during the residency period were in aggregate loss over the reporting
horizons for which internal management dashboards were made available.
No individual account-level data are reproduced in this paper.

\paragraph{Islamic account variant.}
The broker offers a swap-free account variant marketed toward
observant-Muslim clients. The variant removes the Tom-Next swap charge
for holding periods up to a specified duration (typically between
three and fourteen calendar days depending on instrument), after which
a fixed administration fee is applied in lieu. Whether this variant
resolves the five structural concerns identified in
Section~\ref{sec:shariah-analysis} is analysed formally in
Section~\ref{sec:swap-free}; the answer, foreshadowed, is that it
addresses only the first (and only partially), leaving the remaining
three structurally intact.

\section{Shariah Analysis: Structural Concerns in the Standard
    CFD Under the Adopted Framework}
\label{sec:shariah-analysis}

This section applies the adopted AAOIFI/OIC constraint framework
to the formal characterisation of Section~\ref{sec:cfd-mechanics}
and documents five structural incompatibilities or concerns
between the standard retail CFD and the operational constraint
set adopted in Section~\ref{sec:aaoifi-constraints}.
Section~\ref{sec:aaoifi-constraints} formalises the constraint
set; Sections~\ref{sec:riba-analysis}--\ref{sec:kali-analysis}
state and argue each concern as a formal proposition, invoking the
mathematical objects introduced in
Section~\ref{sec:cfd-mechanics}; Section~\ref{sec:swap-free} shows
that the ``swap-free Islamic account'' variant offered by
conventional brokers resolves at most one of these concerns, and
only partially; and Section~\ref{sec:violation-summary}
consolidates the results in a summary table. Throughout the
section, the term \emph{violation} is used only where a
proposition establishes an unconditional failure under an explicit
assumption; where a conclusion is conditional on a scholarly
reading (as, e.g., for \emph{maysir} conditional on
$\beta > 0$, or \emph{bay' al-k\={a}li' bi'l-k\={a}li'}
conditional on the classification of margin and broker services),
the paper uses \emph{concern}, \emph{structural feature}, or
\emph{contested classification} instead.

\subsection{The AAOIFI/OIC operational constraint set adopted in
    this paper}
\label{sec:aaoifi-constraints}

We operationalise the requirements for a bilateral commercial
contract to be reviewable as Shariah-consistent by extracting a
working constraint set from the primary \citet{AAOIFI2015}
Shariah Standards and the resolutions of the International
Islamic Fiqh Academy of the OIC \citep{OICFiqh63,OICFiqh179}. The
constraint set below is the paper's operationalisation; it is
\emph{not} claimed as a universal necessary-and-sufficient
mathematical definition of Shariah validity in general, which
belongs to qualified Shariah supervisory boards and involves
additional considerations (contract combination, prohibited
underlying activities, procedural conditions, and jurisdictional
practice) beyond what a mathematical framework can capture.

\begin{definition}[Operational Shariah constraint set adopted in
    this paper]
\label{def:shariah-valid}
For the purposes of this paper, let $\CC_{\text{Shariah}}$
denote the following operational constraint set used to
evaluate a bilateral commercial contract $\mathcal{K}$,
characterised by the tuple
$(\text{parties},\ \text{subject-matter},\ \text{consideration},\
  \text{terms of exchange})$, against the selected AAOIFI/OIC
framework. A contract is said to be \emph{consistent with
$\CC_{\text{Shariah}}$ under the operationalisation adopted in
this paper} whenever it satisfies each of:
\begin{itemize}
    \item[(C1)] \emph{Rid\={a}}, mutual consent: both parties
    enter the contract freely and with full knowledge of the terms.
    \item[(C2)] \emph{Ma'l\={u}miyya}, the subject matter, the
    consideration, and the terms of exchange are known and
    unambiguous at the time of contracting.
    \item[(C3)] Absence of \emph{rib\={a}}: no stipulated
    increment is charged on a loan or extension of credit
    (\emph{qar\d{d}} plus \emph{ziy\={a}dah}), and no
    increment on a credit-sale or deferred exchange is priced
    off a conventional interest-rate benchmark (\citealp{Usmani2002};
    \citealp{Kamali2000}; AAOIFI Shariah Standards No.~1 on
    Trading in Currencies and No.~8 on \emph{Mur\={a}ba\d{h}a},
    \citealp{AAOIFI2015}).
    \item[(C4)] Absence of \emph{maysir}: the contract does not
    constitute a zero-sum wager whose outcome is determined by
    chance and to which neither counterparty contributes productive
    economic activity (OIC Fiqh Academy Resolution No.~63).
    \item[(C5)] Absence of excessive \emph{gharar}
    (\emph{gharar~f\={a}\d{h}ish}): the uncertainty embedded in the
    subject matter or the execution is not so great as to render
    the contract speculative in essence (AAOIFI Standard~No.~31;
    \citealp{Kamali2000}).
    \item[(C6)] \emph{Qab\d{d}}: where the contract involves the
    sale of an asset, the seller has constructive possession of
    the asset prior to sale (AAOIFI Standard~No.~18).
    \item[(C7)] Absence of \emph{bay' al-k\={a}li' bi'l-k\={a}li'}:
    the contract does not constitute the sale of a deferred
    obligation for a deferred obligation (AAOIFI Standard~No.~8
    \S2/2; Prophetic prohibition reported in the
    \d{h}ad\={\i}th of Ibn Umar).
\end{itemize}
\end{definition}

Constraints (C1) and (C2) are procedural and, at the level of
contract formation, are satisfied by the mutual and voluntary
nature of the retail CFD relationship: the client consents to the
broker's terms and the terms themselves are documented in the
account agreement and product-disclosure statement. The substantive
constraints (C3)--(C7) are the object of
Propositions~\ref{prop:riba}--\ref{prop:bay-kali} below. Denote the
standard CFD contract of Definition~\ref{def:cfd} by
$\mathcal{K}_{\text{CFD}}$. Under the operationalisation adopted
in this paper, this section documents
\begin{equation}
    \mathcal{K}_{\text{CFD}} \notin \CC_{\text{Shariah}},
    \label{eq:cfd-not-shariah}
\end{equation}
via failure of each of (C3), (C4), (C5), (C6), and (C7) in turn.
Whether these failures amount to definitive Shariah-non-compliance
under a particular scholar's chain of \emph{ijtih\={a}d} is a
distinct question that this paper does not resolve; the paper's
contribution is the formal mathematical mapping from contract
mechanics to the operational constraints.

\subsection{\emph{Rib\={a} al-nas\={\i}'ah}: the interest-rate
    financing violation}
\label{sec:riba-analysis}

\begin{proposition}[\emph{Rib\={a}} concern arising from the CFD
    financing mechanism]
\label{prop:riba}
Let $\mathcal{K}_{\text{CFD}}$ be the standard CFD contract of
Definition~\ref{def:cfd} with instantaneous financing rate $F_s$
specified by Definition~\ref{def:instantaneous-financing}. Then
$\mathcal{K}_{\text{CFD}}$ fails to satisfy constraint (C3) of
Definition~\ref{def:shariah-valid} at the mechanical level: the
financing consideration is priced off interbank benchmark rates.
Whether this failure amounts to a definitive determination of
\emph{rib\={a} al-nas\={\i}'ah} is a matter of legal
characterisation for a qualified Shariah Supervisory Board.
\end{proposition}

\begin{proof}
By Proposition~\ref{prop:financing-measurable}, the instantaneous
financing rate $F_s$ is measurable with respect to the
$\sigma$-algebra generated by the interbank overnight rates
$r^{\text{CCY1}}_{s}$ and $r^{\text{CCY2}}_{s}$, with unit
sensitivity to the differential
$r^{\text{CCY2}}_{s} - r^{\text{CCY1}}_{s}$
(equation~\eqref{eq:financing-rate-continuous}).

The classical characterisation of \emph{rib\={a} al-nas\={\i}'ah},
a stipulated fixed increment paid for the deferral of an
exchange, over and above the principal
(\citealp{Kamali2000}, ch.~5; \citealp{Usmani2002}, ch.~3), is
supported by the CFD architecture through two mechanisms whose
combined weight is a matter of scholarly characterisation.
\emph{First} (the mechanical concern), the increment $F_s$ is
priced off the interbank interest-rate differential
(Proposition~\ref{prop:financing-measurable}); AAOIFI Shariah
Standard No.~8 does not permit an increment on a credit or
deferred exchange to be linked to a conventional interest-rate
benchmark (\citealp{AAOIFI2015}), a position echoed in the
general \emph{rib\={a}} treatment of \citet{Usmani2002} and
codified across the subsequent AAOIFI standards. \emph{Second}
(the structural characterisation), the leveraged CFD position
may be characterised as a credit-like extension of exposure by
the broker to the client in nominal amount
$(L-1)/L \cdot N \cdot S_{t_0}$ for the duration $[t_0, t_1]$,
on which $F_s$ is a time-proportional periodic charge. Whether
that characterisation rises to a \emph{qar\d{d}} (loan) in the
strict fiqh sense, so that the increment becomes definitional
\emph{ziy\={a}dah}, is a legal question: a CFD broker does not
necessarily hand the client cash equal to
$(L-1)/L \cdot N \cdot S_{t_0}$, but grants contractual market
exposure against margin. The credit-like-exposure reading is
defended in the contemporary Islamic-finance literature
(\citealp{Usmani2002}; \citealp{Kamali2000}) but requires SSB
determination in each concrete implementation.

The compliance concern is not resolvable by adjusting the sign of
$\delta$, by netting $F_s$ against other cash flows, or by
relabelling the fee as a service charge or ``financing
adjustment''. The mechanical concern (fee priced off interbank
rates) follows directly from
Proposition~\ref{prop:financing-measurable} and requires no
additional \emph{ijtih\={a}d}. The stronger structural
characterisation (leveraged CFD as \emph{qar\d{d}} plus
\emph{ziy\={a}dah}) is offered as a defensible fiqh overlay for
SSB confirmation, not as a mathematical corollary. The paper
therefore documents a mechanical (C3) failure and a structural
concern flagged for scholarly assessment. \qed
\end{proof}

\begin{remark}[Consequence for institutional practice]
\label{rem:institutional-consequence}
The permissibility of Islamic FX Forwards and Islamic Profit Rate
Swaps, offered by Islamic-banking desks at CIMB Islamic, Maybank
Islamic, and Standard Chartered Saadiq, is defended in the Bank
Negara Malaysia tradition on the grounds that they replace $F_s$
with a \emph{wa'd}-based structure whose consideration is not
benchmarked against interbank rates. Under
Proposition~\ref{prop:riba}, this defence is sound only to the
extent that the \emph{wa'd}-based structure's pricing is genuinely
independent of the process $r^{\text{CCY}}_s$.
Section~\ref{sec:alternatives} constructs such structures explicitly
and states the conditions under which the independence is preserved.
\end{remark}

\begin{remark}[Two distinct \emph{rib\={a}} readings]
\label{rem:riba-two-concerns}
Proposition~\ref{prop:riba} distinguishes a \emph{mechanical}
concern (the CFD financing rate $F_s$ is priced off interbank
rates, which follows directly from
Proposition~\ref{prop:financing-measurable}) from a
\emph{structural} concern (the leveraged CFD may be
characterisable as a credit-like exposure whose increment is
\emph{ziy\={a}dah}). The mechanical concern is a mathematical
consequence of the CFD's financing specification and requires no
further \emph{ijtih\={a}d}; the structural characterisation is a
scholarly overlay defensible in the classical \emph{rib\={a}}
literature (\citealp{Usmani2002}; \citealp{Kamali2000}) but
depends on the legal treatment of the CFD's margin-vs-loan
character, which is the province of a qualified SSB. The paper
depends only on the mechanical concern for the (C3) failure.
\end{remark}

\subsection{\emph{Gharar f\={a}\d{h}ish}: the excessive uncertainty
    violation}
\label{sec:gharar-analysis}

\emph{Gharar} denotes uncertainty inherent to a contract, and
\emph{gharar f\={a}\d{h}ish} is the threshold at which such
uncertainty renders the contract invalid. It is essential to
distinguish this concept from ordinary commercial risk: every real
commercial transaction, the purchase of a farm, a manufacturing
plant, an equity share, involves uncertainty about future value.
Such uncertainty is not per se prohibited under Shariah. The
threshold for prohibited \emph{gharar f\={a}\d{h}ish} is
uncertainty that is (a) \emph{material} to the essence of the
contract rather than to accidents of its performance; (b)
\emph{unbounded or unresolvable} at contract formation, in the sense
that essential terms cannot be pinned down; and (c) unaccompanied by
a countervailing \emph{commercial function} that justifies bearing
the risk. Ordinary business risk fails at least the third leg,
the risk-bearing serves the productive purpose of enabling
commerce. Under \citet{Kamali2000} this three-part test is
operationalised through three factors: (i) uncertainty as to the
existence, ownership, or specification of the subject matter;
(ii) uncertainty as to the price or consideration; and
(iii) uncertainty as to the time or manner of delivery. The CFD is
argued below to fail each of these three tests through independent
mechanisms; the strength of the resulting classification as
\emph{gharar f\={a}\d{h}ish} is, ultimately, a fiqh matter that
the propositions in this section document rather than resolve.

\begin{definition}[Execution slippage asymmetry]
\label{def:slippage}
Let $S^{\text{quote}}_{t_0}$ denote the price displayed on the
client's trading terminal at the moment an order is submitted at
time $t_0$, and $S^{\text{fill}}_{t_0}$ the price at which the order
is actually executed by the broker's matching engine. The
\emph{execution slippage} is
\begin{equation}
    \Xi_{t_0} \coloneqq
    \delta \cdot \bigl( S^{\text{fill}}_{t_0}
                       - S^{\text{quote}}_{t_0} \bigr),
    \label{eq:slippage-def}
\end{equation}
adopting the sign convention that $\Xi_{t_0} > 0$ is adverse to the
client. Execution slippage is \emph{asymmetric against the client}
if $\EE[\Xi_{t_0}] > 0$.
\end{definition}

\begin{proposition}[\emph{Gharar f\={a}\d{h}ish} concern]
\label{prop:gharar}
Let $\mathcal{K}_{\text{CFD}}$ be the standard CFD contract of
Definition~\ref{def:cfd}. Then $\mathcal{K}_{\text{CFD}}$
exhibits three independent structural features that together
constitute a substantial \emph{gharar f\={a}\d{h}ish} concern
under the framework of Definition~\ref{def:shariah-valid}:
(a) exposure to a synthetic contingent payoff without ownership,
delivery, or commercial anchoring;
(b) forced-liquidation risk under leveraged maintenance-margin
mechanics; and
(c) execution ambiguity through asymmetric slippage between
quoted and filled prices.
\end{proposition}

\begin{proof}
\emph{Source (a): synthetic contingent payoff without
ownership.} By Remark~\ref{rem:no-delivery}, the client acquires
no rights in the underlying and no commercial exposure other
than a claim to the future price difference $S_{t_1} - S_{t_0}$.
The subject matter of the contract is therefore a synthetic
contingent payoff, not an underlying commodity, currency, or
equity. Under \citet{Kamali2000}, the classical \emph{gharar}
threshold turns not on nonzero variance \emph{per se}
(future prices of farms, businesses, and any real asset are
uncertain) but on the conjunction of (i) contractual
uncertainty as to the essence of the subject matter,
(ii) absence of underlying ownership or delivery, and
(iii) absence of a commercial function that anchors the
risk-bearing to genuine economic activity. The CFD structure
satisfies (i)--(iii) simultaneously: the subject matter is a
scalar $S_{t_1} - S_{t_0}$ rather than a delivered asset; no
ownership of the underlying is transferred at any point
(Remark~\ref{rem:no-delivery}); and the contract does not
involve the acquisition, production, or use of the underlying
in commerce. Under Assumption~\ref{ass:bates}, the payoff
$S_{t_1} - S_{t_0}$ has strictly positive conditional
variance,
\begin{equation}
    \VV\bigl[ S_{t_1} - S_{t_0} \bigm| \FF_{t_0} \bigr] > 0,
    \label{eq:variance-positive}
\end{equation}
with support on $(-S_{t_0}, +\infty)$; but the paper's
\emph{gharar} argument rests on the conjunction (i)--(iii)
rather than on the positive variance alone. This is the
configuration that AAOIFI Standard No.~31 and the Kamali
three-factor test identify as characteristic of prohibited
\emph{gharar} in a bilateral commercial contract.

\emph{Source (b): termination-time uncertainty from
maintenance-margin close-out.} By
Proposition~\ref{prop:leverage-sensitivity}, the client's realised
holding horizon is not entirely within the client's control: with
positive probability $\PP(\tau_{\text{MC}} < t_1) > 0$, the position
is closed out at $\tau_{\text{MC}}$, a stopping time whose
realisation is not observable at contract inception $t_0$. This
creates a termination-time uncertainty. A cautious scholarly
reading is that this is not the same as classical uncertainty
about the time or manner of \emph{delivery of the subject matter}
(a scholar could reasonably object that the contract specifies the
maintenance-margin trigger clearly, and the uncertainty concerns
only whether that trigger occurs, not what is being delivered).
The paper therefore treats source~(b) as a
\emph{termination-time uncertainty} contributing to the broader
\emph{gharar} concern rather than as a self-sufficient
demonstration of \citet{Kamali2000}'s third factor; the
combination with sources~(a) and~(c) supplies the substantive
weight of the (C5) concern.

\emph{Source (c): execution ambiguity through asymmetric slippage.}
Definition~\ref{def:slippage} makes precise the notion of
asymmetric slippage $\Xi = \delta(S^{\text{fill}} - S^{\text{quote}})$
that arises in retail CFD execution. The proposition depends on
the analytical condition $\EE[\Xi] > 0$; that condition is
motivated by, but not directly estimated from, the retail
execution-quality material regulators have published (e.g.\
\citealp{FCA2022}, \citealp{ESMA2018}, which document adverse
execution outcomes, rejected orders and price-improvement
asymmetries without publishing the expectation of the exact
$\Xi$-statistic defined here). The condition $\EE[\Xi] > 0$ is
therefore introduced as an analytical assumption consistent with
the direction of that regulatory evidence, not as an
empirically estimated industry parameter. Under the assumption,
the client cannot observe $\Xi_{t_0}$ at order submission and
bears uncertainty as to the true price of the
transaction, the second of \citet{Kamali2000}'s three
factors.

Each of sources (a), (b), (c) documents a structural feature that
the AAOIFI Standard~No.~31 and OIC Fiqh Academy Resolution No.~63
frameworks identify as characteristic of \emph{gharar
f\={a}\d{h}ish}. Their combination \emph{a fortiori} constitutes an
argument that the CFD fails constraint (C5). Whether this argument
is dispositive under classical fiqh criteria is a scholarly
determination outside the scope of the present paper; the
propositions above establish the structural conditions on which such
a determination would rest. \qed
\end{proof}

\subsection{\emph{Maysir}: the zero-sum wager violation}
\label{sec:maysir-analysis}

\begin{proposition}[Conditional \emph{maysir} argument under
    B-book execution]
\label{prop:maysir}
Let $\mathcal{K}_{\text{CFD}}$ be the standard CFD contract of
Definition~\ref{def:cfd} with internalisation ratio $\beta > 0$.
Then, on the internalised fraction $\beta \cdot N$ of client
notional, $\mathcal{K}_{\text{CFD}}$ satisfies the structural
conditions for \emph{maysir} under \citet{OICFiqh63}. The overall
compliance verdict under constraint (C4) is discussed in
Remarks~\ref{rem:maysir-beyond-beta}--\ref{rem:beta-uncertainty}
below.
\end{proposition}

\begin{proof}
By Proposition~\ref{prop:zero-sum}, under $\beta = 1$ the client's
gross P\&L and the broker's internalisation P\&L sum to zero
pathwise:
\begin{equation}
    \Pi_{\text{gross}}(t_0, t_1) + R^{\text{internal}} = 0
    \quad \text{a.s.}
    \label{eq:pathwise-zero-sum-repeat}
\end{equation}
For intermediate $\beta \in (0, 1)$, the same identity holds on the
internalised fraction $\beta \cdot N$ of the notional exposure. On
that internalised fraction, the transaction is a direct and
pathwise mutually exclusive transfer of wealth between the two
identified counterparties, whose outcome is determined by the
future realisation of the price process $(S_t)_{t \geq 0}$, over
which neither counterparty exercises control.

The paper operationalises \emph{maysir} using the following
three-element structural test, informed by IIFA Resolution
No.~63 (\citealp{OICFiqh63}, which discusses financial-market
structures and characterises certain cash-settled derivative
constructions, including index trading, as impermissible on
the ground that they constitute pure gambling on a fictitious
subject-matter) and the contemporary treatments of
\citet{Kamali2000}, \citet{Usmani2002} and \citet{ElGamal2006}:
a transaction exhibits the structural features of \emph{maysir}
when (i) wealth transfers directly and mutually exclusively
between identified counterparties, (ii) the transfer is
contingent on a future event whose outcome is uncertain at
inception, and (iii) neither counterparty contributes productive
economic activity to the outcome. The internalised component of
the CFD satisfies each element:
(i) the transfer is direct and mutually exclusive between the two
identified counterparties (Prop.~\ref{prop:zero-sum});
(ii) it depends on $S_{t_1} - S_{t_0}$, unresolved at $t_0$;
(iii) neither counterparty contributes to the productive activity
underlying $S$ (Remark~\ref{rem:no-delivery}).
Hence, on the internalised fraction, the structural test above
is satisfied; whether this is dispositive of \emph{maysir} under
a specific SSB reading is a matter outside the paper's scope. \qed
\end{proof}

\begin{remark}[Beyond B-book: the counterparty-agnostic reciprocity
    argument]
\label{rem:maysir-beyond-beta}
Proposition~\ref{prop:maysir} is deliberately conditional on
$\beta > 0$: the pathwise zero-sum identity of
Proposition~\ref{prop:zero-sum} rests on the broker being the
client's direct counterparty for the internalised notional. A
separate, counterparty-agnostic, argument may be constructed: the
CFD is a cash-settled bilateral contract in which for every dollar
the client gains, some counterparty, broker (B-book), external
liquidity provider (A-book), or mixture thereof, loses that
dollar, because no productive real-world activity generates the
P\&L. Whether this counterparty-agnostic reciprocity itself
constitutes \emph{maysir} under classical criteria is a distinct
scholarly question. Some contemporary scholars
(\citealp{Usmani2002,ElGamal2006}) treat the reciprocity as
sufficient regardless of the identity of the counterparty; others
restrict the classification to direct-counterparty structures. This
paper proves only the direct-counterparty case; the broader
argument is noted as a scholarly extension awaiting SSB
determination.
\end{remark}

\begin{remark}[Empirical uncertainty about $\beta$]
\label{rem:beta-uncertainty}
Any claim that retail CFD brokers exhibit large $\beta$ values is
\emph{inferential} rather than directly measured: regulatory
disclosures report the fraction of client accounts that lose
money, not the internalisation ratio. A large disclosed loss rate
is consistent with substantial B-book internalisation but does not
identify $\beta$ uniquely, external counterparties may also
profit systematically from adversely selected retail flow, a
possibility surveyed in the retail-microstructure literature
(\citealp{BarberOdean2000}). The empirical claim underlying the
(C4) analysis is therefore only that $\beta$ is materially
positive for the retail segment; the proposition above is
careful to depend solely on this qualitative claim rather than
on any specific point estimate.
\end{remark}

\begin{remark}[Distinction from commercial hedging]
\label{rem:hedging-comparison}
A common defence of derivative instruments in Islamic finance is
that they may permissibly discharge a hedging function: a farmer
using a futures contract to hedge a harvest, an exporter using an
FX forward to hedge foreign-currency receivables. This defence is
inapplicable to the retail CFD use case, in which the participant
has no underlying commercial exposure to hedge; the CFD position
\emph{is} the exposure. See Section~\ref{sec:alternatives} for
treatment of the genuine hedging use case under a \emph{wa'd}-based
framework consistent with the \citet{Bacha2013} tradition.
\end{remark}

\subsection{The contract-classification problem}
\label{sec:qabd-analysis}

An objection to the standard \emph{qab\d{d}} argument is that
AAOIFI Standard~No.~18's constructive-possession requirement
applies specifically to contracts of sale (\emph{bay'}); a
respondent might argue that a cash-settled CFD is not a
\emph{bay'} of the underlying at all, and that the \emph{qab\d{d}}
requirement is therefore misapplied. Proposition~\ref{prop:qabd}
below reframes the analysis around the more fundamental question
of \emph{contract classification}: the CFD, on inspection, does
not correspond to any classical nominate contract. Non-fit with
the nominate contracts is not by itself a determination of
impermissibility (classical fiqh admits new contract forms
subject to the substantive prohibitions); it means the
substantive prohibitions of \emph{rib\={a}}, \emph{gharar}, and
\emph{maysir} apply directly to the novel contract without the
mediating discipline of a recognised nominate category, which is
the deeper scholarly concern.

\begin{proposition}[Contract-classification: the CFD is not any
    classical nominate contract]
\label{prop:qabd}
The standard CFD contract of Definition~\ref{def:cfd} does not
correspond to any of the classical nominate contract categories
enumerated in \citet[Sec.~2.1]{AAOIFI2015} and reviewed in
Section~\ref{sec:lit-fiqh}: it is not \emph{bay'} (no
\emph{qab\d{d}} of the underlying; no transfer of ownership rights
characteristic of a sale); not \emph{salam} (no immediate full
payment of price against a deferred fungible commodity); not
\emph{isti\d{s}n\={a}'} (no commissioning of manufacture);
not \emph{shar\={\i}kah} (no capital pooling for joint ownership);
not \emph{ij\={a}ra} (no lease of usufruct);
not \emph{wak\={a}la} (no defined agency for a service fee); and
not \emph{wa'd} (no unilateral binding promise). It is,
structurally, a novel contract of cash-settled contingent
obligations. Under the general principles surveyed by
\citet{Kamali2000}, \citet{Usmani2002}, and \citet{ElGamal2006},
contracts falling outside the recognised nominate categories are
evaluated against the substantive prohibitions of
\emph{rib\={a}}, \emph{gharar}, and \emph{maysir}, the analyses
of Sections~\ref{sec:riba-analysis}--\ref{sec:maysir-analysis}
above.
\end{proposition}

\begin{proof}
Verification of non-fit is by inspection of
Definitions~\ref{def:cfd} and \ref{def:instantaneous-financing}
against the classical categories catalogued in the literature
review. The CFD confers no ownership of, or legal claim to, $N$
units of the underlying (Remark~\ref{rem:no-delivery}); it
requires no upfront payment of the notional (only initial margin
$M_0$); it involves no manufacture, lease, agency, partnership, or
unilateral promise. The residual category is that of
cash-settled contingent obligation, which is not a classical
nominate category. Under the general-principles fallback of
\citet{Kamali2000}, the substantive prohibitions apply directly,
which is the subject of Sections~\ref{sec:riba-analysis}--\ref{sec:maysir-analysis}.
\qed
\end{proof}

\begin{remark}[On the classical \emph{qab\d{d}} formulation]
\label{rem:qabd-classical}
Proposition~\ref{prop:qabd} is a reformulation of the more
familiar claim that the CFD violates the \emph{qab\d{d}}
requirement of AAOIFI Standard~No.~18. That familiar claim
assumes, implicitly, that the CFD is (or purports to be) a
\emph{bay'} of the underlying, and then observes that
\emph{qab\d{d}} is absent. The present reformulation avoids the
assumption: rather than asking whether the CFD satisfies the
\emph{qab\d{d}} requirement of a \emph{bay'}, it asks whether the
CFD is a \emph{bay'} in the first place. The answer is no, and
also no to every other nominate category. Constraint (C6) is
therefore best understood not as a violation of \emph{qab\d{d}}
in isolation but as a symptom of the deeper non-fit problem,
which is itself the foundational scholarly concern.
\end{remark}

\subsection{\emph{Bay' al-k\={a}li' bi'l-k\={a}li'}: the
    deferred-for-deferred concern}
\label{sec:kali-analysis}

\begin{proposition}[Deferred-for-deferred structural feature]
\label{prop:bay-kali}
The standard CFD contract of Definition~\ref{def:cfd} with
$t_0 < t_1$ exhibits the reciprocal-deferred-contingent structure
that the classical prohibition of
\emph{bay' al-k\={a}li' bi'l-k\={a}li'} addresses.
\end{proposition}

\begin{proof}[Proof sketch of the structural feature]
The Prophetic prohibition against
\emph{bay' al-k\={a}li' bi'l-k\={a}li'}, reported in the
\d{h}ad\={\i}th of Ibn Umar and codified in AAOIFI Shariah
Standard~No.~8 \S2/2, forbids the exchange of one deferred
obligation for another. Under \emph{the reading in which the
initial margin $M_0$ is classified as a security deposit rather
than as consideration}, and \emph{the reading in which broker
execution, platform access and counterparty performance do not
themselves constitute immediate counter-consideration}, the
following structural feature is present. At contract inception
$t_0$, the client posts only $M_0 = N \cdot S_{t_0} / L$
(equation~\eqref{eq:initial-margin}), which under the assumed
reading is a security deposit held against margin-call events;
the broker under the same reading delivers no goods, no cash,
and no asset at $t_0$. Both parties' substantive obligations,
settlement of $\Pi_{\text{net}}(t_0, t_1)$ per
equation~\eqref{eq:cfd-net-payoff}, are deferred to $t_1$ and
are moreover contingent on the terminal state $S_{t_1}$. Neither
party under the assumed reading takes immediate \emph{qab\d{d}}
(Proposition~\ref{prop:qabd}) or renders immediate
counter-consideration. Under those readings, the
reciprocal-deferred-contingent structural feature that the
classical prohibition addresses is present. Each reading is
plausible and each is contested: a scholar could take the
alternative view (that margin is partial consideration; that
broker services constitute immediate counter-consideration), in
which case the structural feature would not be present in the
same form. The paper documents the structural concern under the
broad reading and leaves the definitive classification to SSB
review. \qed
\end{proof}

\begin{remark}[Scope of the classical prohibition]
\label{rem:kali-scope}
The classical prohibition of \emph{bay' al-k\={a}li' bi'l-k\={a}li'}
addresses, at its narrowest reading, the sale of a debt for a debt.
Its extension to non-sale contracts exhibiting the same
reciprocal-deferred structure is a matter of scholarly extension:
\citet{Usmani2002} and the majority of GCC scholars take the broad
view, extending the prohibition wherever both counter-values are
deferred; a subset of Malaysian scholars restrict the classical
prohibition to formal sales. Combined with the non-fit result of
Proposition~\ref{prop:qabd}, however, the CFD is on any reading
outside the class of contracts recognised as permissible under
classical fiqh: it is neither a permitted deferred sale (e.g.\
\emph{salam}, in which the price is paid in full at inception) nor
a permitted spot sale (in which both counter-values are exchanged
immediately). Whether its exclusion is best framed as violation of
(C7) narrowly or as violation of the general non-fit principle of
Proposition~\ref{prop:qabd} is a scholarly question; the paper
records the structural feature and both framings.
\end{remark}

\begin{remark}[The margin-as-consideration defence and the
    broker-service defence]
\label{rem:margin-consideration}
Two defences may be raised against
Proposition~\ref{prop:bay-kali}. The \emph{margin-as-consideration
defence} holds that $M_0$ constitutes immediate consideration by
the client; against this, $M_0$ is fully returnable to the client
at $t_1$ absent adverse price movement of magnitude exceeding
$(1-\alpha) M_0$, which is more characteristic of a security
deposit than of a payment for goods or services. The
\emph{broker-service defence} holds that the broker's execution,
platform access, counterparty performance, and liquidity
provision are themselves immediate counter-consideration; against
this, these services are not the primary subject-matter of the
CFD (which is the price-differential settlement at $t_1$) and are
in any event provided against separate transaction fees where
they are charged. Neither defence is definitively refuted; both
are plausible scholarly positions the paper acknowledges. The
paper documents the structural concern that survives \emph{if} one
takes the broad reading of the classical prohibition together with
the assumed reading of margin and broker services, and leaves the
definitive classification to SSB review.
\end{remark}

\subsection{The ``Islamic account'' variant: a formal analysis}
\label{sec:swap-free}

We now analyse whether the swap-free ``Islamic account'' variant
offered by conventional brokers (Section~\ref{sec:primary-source})
resolves the five violations proved in
Sections~\ref{sec:riba-analysis}--\ref{sec:kali-analysis}.

\begin{definition}[Swap-free variant]
\label{def:swap-free}
The swap-free Islamic-account variant
$\mathcal{K}_{\text{CFD-SF}}$ modifies the standard CFD contract
$\mathcal{K}_{\text{CFD}}$ by replacing the Tom-Next financing rate
of Definition~\ref{def:instantaneous-financing} with
\begin{equation}
    \widetilde{F}_s =
    \begin{cases}
        0
        & \text{if } s - t_0 \leq \Delta_{\text{grace}}, \\[4pt]
        \phi_{\text{admin}}
        & \text{if } s - t_0 > \Delta_{\text{grace}},
    \end{cases}
    \label{eq:swap-free-rate}
\end{equation}
where $\Delta_{\text{grace}}$ is a broker-specified grace period
(typically in the range 3--14 calendar days) and
$\phi_{\text{admin}} > 0$ is a fixed per-diem administrative fee
accruing beyond the grace window. All other elements of the
contract tuple $(t_0, t_1, \delta, N, L, s, m)$ are unchanged.
\end{definition}

\begin{proposition}[Structural insufficiency of the swap-free
    variant]
\label{prop:swap-free}
The swap-free variant $\mathcal{K}_{\text{CFD-SF}}$ resolves
constraint (C3) partially, only for holding periods
$s - t_0 \leq \Delta_{\text{grace}}$, and leaves constraints
(C4), (C5), (C6), and (C7) structurally unresolved.
\end{proposition}

\begin{proof}
\emph{(i) Constraint (C3).} For
$s - t_0 \leq \Delta_{\text{grace}}$, the modified financing rate
$\widetilde{F}_s = 0$ carries no interest-rate dependence and
Proposition~\ref{prop:financing-measurable} does not apply. For
$s - t_0 > \Delta_{\text{grace}}$, the flat fee
$\phi_{\text{admin}}$ may or may not violate (C3) depending on
whether it is calibrated against prevailing interbank rates in the
broker's pricing model. Brokers do not typically disclose this
calibration, and the scholarly consensus is that flat administrative
fees are permissible only if their magnitude is not indexed to
interest-rate benchmarks either directly or by regular recalibration
\citep{AlSuwailem2006,Uberoi2008}. The (C3) resolution is therefore
partial and conditional.

\emph{(ii) Constraint (C4).} The internalisation ratio $\beta$ and
the pathwise zero-sum relation of Proposition~\ref{prop:zero-sum}
are independent of the specification of the financing rate.
Proposition~\ref{prop:maysir} applies verbatim to
$\mathcal{K}_{\text{CFD-SF}}$. The \emph{maysir} violation is
unresolved.

\emph{(iii) Constraint (C5).} The three sources of \emph{gharar}
identified in Proposition~\ref{prop:gharar}, unbounded terminal
price uncertainty, forced-liquidation risk, and execution ambiguity,
are all independent of the financing rate. The
\emph{gharar~f\={a}\d{h}ish} violation is unresolved.

\emph{(iv) Constraint (C6).} The absence of \emph{qab\d{d}}
(Proposition~\ref{prop:qabd}) is a property of the CFD contract
structure itself and is independent of the specification of
$\widetilde{F}_s$. Unresolved.

\emph{(v) Constraint (C7).} The
\emph{bay' al-k\={a}li' bi'l-k\={a}li'} structure
(Proposition~\ref{prop:bay-kali}) is likewise independent of the
financing rate. Unresolved. \qed
\end{proof}

\begin{remark}[Marketing versus structural compliance]
\label{rem:marketing-vs-structural}
The swap-free variant is therefore best understood as a marketing
product that addresses the most conspicuous fiqh objection,
interest-linked financing, while leaving the deeper structural
objections intact. Its usage may be defended only under a chain of
\emph{ijtih\={a}d} that treats the \emph{qab\d{d}}, \emph{maysir},
\emph{gharar}, and \emph{k\={a}li'-bi'l-k\={a}li'} requirements as
non-operative for cash-settled derivative contracts, a position
that is not the one adopted in the AAOIFI/OIC framework used in
this paper, nor in the classical fiqh treatments of
\citet{Kamali2000}, \citet{ElGamal2006}, and \citet{Usmani2002}.
Alternative permissive positions on cash-settled derivative
structures exist in the Bank Negara Malaysia and ISRA scholarly
tradition (see Section~\ref{sec:alt-wad} and
Remark~\ref{rem:wad-debate}) and remain the subject of active
scholarly disagreement.
\end{remark}

\subsection{Summary of structural violations}
\label{sec:violation-summary}

The results of
Sections~\ref{sec:riba-analysis}--\ref{sec:swap-free} are
consolidated in Table~\ref{tab:violation-summary}. The table
records, for each of the five substantive Shariah constraints, the
proposition establishing the violation, the primary AAOIFI Standard
or OIC resolution invoked, and the compliance status under the
standard retail CFD (Definition~\ref{def:cfd}), the swap-free
variant (Definition~\ref{def:swap-free}), and each of the three
alternative architectures constructed in
Section~\ref{sec:alternatives}.

\begin{table}[!ht]
\centering
\small
\caption{Documented Shariah concerns for retail CFD products and
their treatment across contract variants. Entries describe the
paper's characterisation of each concern under the standard CFD,
the ``swap-free Islamic account'' variant (SF variant), and the
\emph{shar\={\i}kat al-'in\={a}n} alternative architecture of
Section~\ref{sec:alternatives}. All classifications are
provisional and subject to review by qualified Shariah supervisory
boards.}
\label{tab:violation-summary}
\begin{tabular}{@{}p{2.6cm}p{2.6cm}p{1.9cm}p{2.4cm}p{2.0cm}p{2.4cm}@{}}
\toprule
\textbf{Concern} & \textbf{Primary source}
    & \textbf{Argued in}
    & \textbf{CFD} & \textbf{SF variant}
    & \textbf{Alt.\ arch.} \\
\midrule
(C3) \emph{rib\={a}}
    & AAOIFI SS 1, 8; Usmani, Kamali
    & Prop.~\ref{prop:riba}
    & structural concern
    & partially addressed
    & satisfied by design \\
(C4) \emph{maysir}
    & OIC Res.~63 (1/7)
    & Prop.~\ref{prop:maysir}
    & conditional on $\beta > 0$
    & unchanged
    & satisfied by design \\
(C5) \emph{gharar}
    & AAOIFI SS~31; Kamali
    & Prop.~\ref{prop:gharar}
    & documented concern
    & unchanged
    & largely mitigated \\
(C6) \emph{qab\d{d}} / non-fit
    & AAOIFI SS~18
    & Prop.~\ref{prop:qabd}
    & non-fit with nominate contracts
    & unchanged
    & joint ownership from $t_0$ \\
(C7) \emph{k\={a}li'}
    & AAOIFI SS~8
    & Prop.~\ref{prop:bay-kali}
    & structural concern
    & unchanged
    & simultaneous consideration \\
\bottomrule
\end{tabular}
\end{table}

Table~\ref{tab:violation-summary} makes the paper's central
negative finding precise: the standard retail CFD violates five
foundational Shariah constraints, and the ``swap-free Islamic
account'' variant offered by conventional brokers resolves at most
one of these, and only partially. The task of
Section~\ref{sec:alternatives} is to construct market-access
architectures that resolve all five simultaneously, and
Section~\ref{sec:simulation} quantifies the economic cost of doing
so.

\section{First-Principles Alternative Architectures}
\label{sec:alternatives}

Sections~\ref{sec:cfd-mechanics}--\ref{sec:shariah-analysis}
established that the standard retail CFD violates five foundational
constraints of Islamic commercial law and that the swap-free variant
marketed by conventional brokers resolves at most one of these, and
only partially. This section undertakes the constructive task of the
paper: to design market-access architectures that satisfy all five
constraints simultaneously and to derive their payoff dynamics under
the same jump-diffusion process (Assumption~\ref{ass:bates}) used for
the standard CFD.

Section~\ref{sec:design-constraints} restates the constraint set of
Section~\ref{sec:shariah-analysis} as design requirements on any
candidate architecture.
Sections~\ref{sec:alt-wad}--\ref{sec:alt-sharika} introduce and
analyse three architectures,
\emph{wa'd}-based FX exposure,
organised \emph{tawarruq} for financing, and
\emph{shar\={\i}kat al-'in\={a}n} pooled-equity market access, with
the third being the paper's principal constructive contribution.
Section~\ref{sec:arch-comparison} consolidates the architectures in a
comparative table. Section~\ref{sec:no-equivalent} acknowledges the
retail CFD use cases for which no permissible equivalent exists.

\subsection{Design constraints on a compliant architecture}
\label{sec:design-constraints}

The Shariah constraint set $\CC_{\text{Shariah}}$ of
Definition~\ref{def:shariah-valid} translates, at the level of
contract engineering, into five design conditions on any candidate
market-access architecture.

\begin{definition}[Operational design constraints adopted in this
    paper]
\label{def:compliant-arch}
The design conditions (D1)--(D5) below are the operational
constraints adopted in this paper for evaluating candidate
market-access architectures against the framework of
Definition~\ref{def:shariah-valid}. They are not claimed as
necessary and sufficient conditions for Shariah compliance in
general (a claim which properly belongs to qualified Shariah
supervisory boards and involves additional considerations
including agency, custody, contract combination, prohibited
underlying activities, contractual conditions, and settlement
rights). They are a researcher's working framework that
operationalises the five concerns of
Definition~\ref{def:shariah-valid} in a form amenable to
mathematical characterisation.

\begin{itemize}
    \item[(D1)] \emph{No interest-linked consideration.} The
    periodic financing rate $F_s^{\mathcal{K}}$, if any, has no
    functional dependence on any interbank rate process. Let
    $\mathcal{G}^{\mathcal{K}}_{s}$ denote the $\sigma$-algebra
    generated by the underlying price and any other
    non-rate variables that may enter $F_s^{\mathcal{K}}$; the
    formal requirement is that $F_s^{\mathcal{K}}$ is
    $\mathcal{G}^{\mathcal{K}}_{s}$-measurable, or equivalently
    $\partial F_s^{\mathcal{K}} / \partial r^{\text{CCY1}}_{s}
     = \partial F_s^{\mathcal{K}} / \partial r^{\text{CCY2}}_{s}
     = \partial F_s^{\mathcal{K}} / \partial r_s^{\text{IBOR}}
     = 0$
    holds jointly. Addresses (C3).
    \item[(D2)] \emph{Constructive possession
    (\emph{qab\d{d}}) where a sale is present.} For any contract
    element construed as a sale of an underlying asset, at least
    one counterparty holds legal or constructive ownership of the
    traded quantity prior to sale. Addresses (C6).
    \item[(D3)] \emph{Real-economic-activity linkage.} The payoff
    derives from the acquisition, production, transformation, or
    delivery of an identifiable real asset, not from cash-settled
    price differences alone. Addresses (C4) and (C7).
    \item[(D4)] \emph{Proportional risk sharing.} Where the
    architecture is constructed as a partnership,
    \emph{losses} must be allocated strictly in proportion to
    capital contribution (AAOIFI Standard~No.~12 \S3/1/5/2);
    \emph{profits} may be distributed under the pre-agreed
    permissible profit-sharing ratio, which need not equal the
    capital ratio. Addresses (C4).
    \item[(D5)] \emph{Immediate substantive consideration where
    required.} At contract inception, at least one counterparty
    renders immediate non-contingent consideration in
    circumstances where the applicable contract category so
    requires. Addresses (C7).
\end{itemize}
The remaining substantive concern, excessive
\emph{gharar} (C5), is treated as a parameter constraint
rather than a binary design condition: candidate parameterisations
that introduce forced-liquidation risk on the client, unbounded
terminal exposure without commercial backing, or execution
ambiguity in the sense of Definition~\ref{def:slippage} are
flagged for scrutiny rather than automatically excluded.
\end{definition}

Sections~\ref{sec:alt-wad}--\ref{sec:alt-sharika} evaluate three
candidate architectures against Definition~\ref{def:compliant-arch},
recording explicitly where each falls short of one or more of
(D1)--(D5) under alternative scholarly readings.

\subsection{Architecture A: \emph{Wa'd}-based FX exposure}
\label{sec:alt-wad}

The \emph{wa'd}-based FX exposure architecture is a well-documented
Shariah-compliant substitute for the conventional FX forward and FX
swap in institutional Islamic banking. It has been permitted by the
Shariah Advisory Council of Bank Negara Malaysia
since 2005 (\citealp{Bacha2013}) and forms the basis of the Islamic
FX Forward and Islamic Profit Rate Swap products offered by a
number of Malaysian, Middle Eastern, and international Islamic
banking desks; comprehensive market-share statistics for the
Islamic derivatives universe are not publicly available. It is
contested by a substantial portion of Gulf Cooperation Council
scholars, whose objections we record honestly.

\begin{definition}[\emph{Wa'd}-based FX exposure]
\label{def:wad-fx}
The \emph{wa'd}-based FX exposure architecture
$\mathcal{K}_{\text{wa'd}}$ on the FX pair CCY1/CCY2 comprises the
following.

\emph{At contract inception $t_0$:}
\begin{enumerate}
    \item Client $A$ issues a unilateral binding promise
    (\emph{wa'd mulzim}) to purchase $N$ units of CCY1 from
    Broker $B$ at the pre-agreed rate $K$ at delivery time $T$.
    \item Broker $B$ issues a corresponding unilateral binding
    promise to sell $N$ units of CCY1 to Client $A$ at rate $K$ at
    delivery time $T$.
\end{enumerate}

\emph{At delivery time $T$:}
\begin{enumerate}
    \setcounter{enumi}{2}
    \item Both parties execute a spot FX transaction: $A$ pays
    $N \cdot K$ units of CCY2 to $B$, and $B$ delivers $N$ units of
    CCY1 to $A$. Both parties take immediate \emph{qab\d{d}} of the
    exchanged currencies.
\end{enumerate}
The client may thereafter sell the received $N$ units of CCY1 at
the prevailing spot rate $S_T$, for a net cash outcome in CCY2 as
in Proposition~\ref{prop:wad-payoff} below.
\end{definition}

\begin{proposition}[\emph{Wa'd} payoff dynamics under
    Assumption~\ref{ass:bates}]
\label{prop:wad-payoff}
Under Definition~\ref{def:wad-fx} and Assumption~\ref{ass:bates} on
the underlying spot process $(S_t)$, the client's terminal cash
outcome in units of CCY2, net of the broker's structuring fee
$u^{\text{wa'd}} \geq 0$, is
\begin{equation}
    \Pi^{\text{wa'd}}_{c}(t_0, T)
    = N \cdot \bigl( S_T - K \bigr) - u^{\text{wa'd}}.
    \label{eq:wad-payoff}
\end{equation}
The value at $t \in [t_0, T]$ is
$\Pi^{\text{wa'd}}_{c}(t) = N \cdot (S_t - K)\,e^{-r(T-t)}
- u^{\text{wa'd}}\,e^{-r(T-t)}$
under the risk-neutral measure.
\end{proposition}

\begin{proof}
By construction of Definition~\ref{def:wad-fx}, both promises are
executed at $T$ at rate $K$: the client receives $N$ units of CCY1
in exchange for $N \cdot K$ units of CCY2. Immediate sale of the $N$
units of CCY1 at the prevailing spot $S_T$ yields $N \cdot S_T$
units of CCY2. Netting the two cash flows and the ujrah gives
\eqref{eq:wad-payoff}. Measurability of $\Pi^{\text{wa'd}}_c$ with
respect to $(S_t)$ is immediate; no dependence on interbank rates
enters through the transaction structure. The time-$t$ value is a
standard risk-neutral valuation
(\citealp{Bjork2009}, ch.~7). \qed
\end{proof}

\paragraph{Cost of carry: \emph{wa'd} versus conventional NDF.}
Under the risk-neutral measure, the expected terminal spot equals
the forward implied by covered interest rate parity,
$\EE^{*}[S_T] = F^{\text{CIP}}_0 = S_{t_0}\,
e^{(r^{\text{CCY2}} - r^{\text{CCY1}})(T - t_0)}$
(\citealp{Bjork2009}, ch.~24). If the \emph{wa'd} strike is set to
$K = F^{\text{CIP}}_0$, the expected payoff net of ujrah equals zero,
economically equivalent to a conventional non-deliverable
forward (NDF). The scholarly debate turns on precisely this point:
if $K$ is calibrated to the CIP-implied forward rate, the
\emph{wa'd} structure inherits an economic dependence on interbank
rates even though its contractual mechanics do not reference them.
Remark~\ref{rem:institutional-consequence} foreshadowed this
concern.

\begin{remark}[The scholarly debate on the double-\emph{wa'd}
    structure]
\label{rem:wad-debate}
The most serious concern with Definition~\ref{def:wad-fx} is
not (C3) but the more fundamental question of whether a
\emph{bilateral} binding \emph{wa'd} structure, both parties
issuing binding promises against each other, is permissible
at all. AAOIFI Shariah Standard No.~1 draws an express
distinction between a single binding unilateral promise (which
may be permissible under specified conditions) and a binding
bilateral promise (\emph{muw\={a}'adah}) to purchase and sell
currencies against each other (\citealp{AAOIFI2015}). AAOIFI
Shariah Standard No.~49 (\citealp{AAOIFI-SS49}) provides the
standards-body's dedicated treatment of the distinction and of
the conditions under which each may operate. The International
Islamic Fiqh Academy Resolution No.~238
(\citealp{OICFiqh238}) discusses mutual binding commitments in
future currency exchange along similar lines. Under the restrictive reading of these
authorities, Definition~\ref{def:wad-fx}, which requires both
counterparties to issue binding promises, would not be
permissible even before the (C3) pricing question arises.

The permissive tradition associated with the Bank Negara Malaysia
Shariah Advisory Council and ISRA scholarship is more nuanced
than a blanket endorsement of double-\emph{wa'd} structures.
BNM's earlier rulings on forward FX permitted the use of a
\emph{unilateral} binding \emph{wa'd} by one counterparty and,
notably, took the position that no fee could be charged on the
promisee for the promise itself (a fee on the promise being one
of the features that converts a unilateral into a de facto
bilateral commitment). Later BNM rulings addressed
\emph{muw\={a}'adah} more directly and admitted certain bilateral
structures subject to specified governance conditions. The exact
architecture of Definition~\ref{def:wad-fx}, a double binding
\emph{wa'd} plus a structuring fee, should not therefore be
assumed to inherit any of these BNM rulings automatically; each
specific implementation would need SSB assessment under whichever
authority governs it. The \emph{restrictive} position associated
with \citet{Usmani2002}, \citet{ElGamal2006}, and the majority of
GCC scholars treats such bilateral binding structures as a form
of \emph{\d{h}iyal} that reproduces the economic substance of a
forward contract. This paper adopts the restrictive AAOIFI
reading; the divergence is recorded, and consequently
Definition~\ref{def:wad-fx} is \emph{not} presented as an
unconditionally compliant architecture.

A separate concern arises with the structuring fee
$u^{\text{wa'd}} \geq 0$. Both AAOIFI SS~No.~49
(\citealp{AAOIFI-SS49}, in its general treatment of promise-based
structures) and IIFA Resolution No.~238
(\citealp{OICFiqh238}) disfavour consideration paid for the
binding commitment itself; a fee-on-promise is one of the features that transforms
a unilateral promise into an effectively bilateral one. The
paper's treatment of $u^{\text{wa'd}}$ therefore requires that
the fee be clearly identified, in a separately executed
management-services agreement, as consideration for a
distinguishable service (execution, documentation, custody),
never as payment for the promise itself; whether that
separation is respected in practice is a matter for SSB review
of any specific implementation.
\end{remark}

\begin{proposition}[Contested compliance of
    $\mathcal{K}_{\text{wa'd}}$]
\label{prop:wad-comply}
Under Definition~\ref{def:wad-fx}, the compliance status of
$\mathcal{K}_{\text{wa'd}}$ with the design constraints
(D1)--(D5) of Definition~\ref{def:compliant-arch} is contested:
\begin{itemize}
    \item[(a)] The bilateral binding-promise structure of clauses
    1--2 of Definition~\ref{def:wad-fx} sits in tension with the
    AAOIFI restrictive reading of Standard No.~1 on
    \emph{muw\={a}'adah} in currencies, and with the IIFA
    Resolution No.~238 position; under those authorities the
    architecture is not permissible in its stated form.
    \item[(b)] Under the permissive reading (Bank Negara Malaysia
    Shariah Advisory Council), the architecture may be
    permissible subject to governance conditions on the
    unilateral character of each promise, actual delivery at $T$,
    and separately identifiable service consideration for
    $u^{\text{wa'd}}$ per Remark~\ref{rem:wad-debate}.
    \item[(c)] Even on the permissive reading, (D1) is respected
    only if the strike $K$ and the fee $u^{\text{wa'd}}$ are set
    independently of the interbank-rate processes; calibration
    to the CIP-implied forward $F^{\text{CIP}}_0$ reintroduces
    the concern of Proposition~\ref{prop:financing-measurable}.
\end{itemize}
\end{proposition}

\begin{proof}[Proof sketch]
Clauses (a)--(c) follow by inspection of Definition~\ref{def:wad-fx}
against, respectively, the AAOIFI SS No.~1 restriction on
binding bilateral currency promises, the Bank Negara Malaysia
permissive framework, and
Proposition~\ref{prop:financing-measurable}. The paper does not
adjudicate between the permissive and restrictive readings;
both are recorded, and consequently the architecture is not
carried forward in this paper as a strictly recommended option.
\qed
\end{proof}

\begin{remark}[Consequence for the paper]
\label{rem:wad-outcome}
Because Architecture A is contested at the fundamental level
(not only in its calibration under (D1)), this paper's
constructive contribution is centred on Architecture C
(\emph{shar\={\i}kat al-'in\={a}n}), Section~\ref{sec:alt-sharika}.
Architecture A is retained as an intermediate reference against
which the constructive contribution is contrasted, not as a
standalone recommendation.
\end{remark}

\begin{proof}[Proof sketch]
(C4): the transaction is not zero-sum, both parties genuinely
exchange currencies to their commercial benefit; there is no
mutually exclusive transfer conditional on price. (C5): the
contract terms and delivery are fixed at $t_0$; there is no
forced-liquidation risk in the sense of
Proposition~\ref{prop:leverage-sensitivity}. (C6): both parties
take immediate \emph{qab\d{d}} at $T$. (C7): the client's payment
and the broker's delivery at $T$ are simultaneous, not both
deferred. (C3): the contract mechanics of
Definition~\ref{def:wad-fx} do not reference interbank rates, so
(C3) is satisfied at the mechanical level; but if $K$ or
$u^{\text{wa'd}}$ is calibrated to $F^{\text{CIP}}_0$, the
economic exposure inherits interbank dependence and (C3) is
violated at the calibration level. \qed
\end{proof}

\subsection{Architecture B: Organised \emph{tawarruq} financing
    overlay}
\label{sec:alt-tawarruq}

The \emph{tawarruq} structure synthesises cash financing through a
coordinated sequence of commodity transactions. It is included for
completeness, organised \emph{tawarruq} is widely deployed
across Middle Eastern and Malaysian Islamic financing desks, though
without publicly available market-share statistics, but is
\emph{not} recommended by this paper, on the grounds of the
scholarly consensus recorded in
Remark~\ref{rem:tawarruq-criticism}.

\begin{definition}[Organised \emph{tawarruq} financing overlay]
\label{def:tawarruq}
The organised \emph{tawarruq} financing overlay
$\mathcal{K}_{\text{taw}}$ comprises the following sequence.

\emph{At inception $t_0$:}
\begin{enumerate}
    \item Bank purchases a quantity $Q$ of a fungible commodity,
    typically London Metal Exchange copper or aluminium warrants
   , at spot price $P_c$ per unit.
    \item Bank sells the $Q$ units to the client at the deferred
    price $P_d = P_c \cdot (1 + \mu)$ per unit, payable at maturity
    $t_1$, where $\mu \geq 0$ is the \emph{mur\={a}ba\d{h}a}
    markup.
    \item Client sells the $Q$ units immediately at spot to a
    third-party counterparty at price $P_r \leq P_c$ per unit,
    receiving cash $Q \cdot P_r$.
\end{enumerate}

\emph{At maturity $t_1$:}
\begin{enumerate}
    \setcounter{enumi}{3}
    \item Client pays $Q \cdot P_d$ to the bank in cash.
\end{enumerate}
\end{definition}

\begin{proposition}[\emph{Tawarruq} financing cost]
\label{prop:tawarruq-cost}
The client's net financing cost over $[t_0, t_1]$ under
Definition~\ref{def:tawarruq} is
\begin{equation}
    C^{\text{taw}}(t_0, t_1)
    = Q \cdot \bigl( P_d - P_r \bigr)
    = Q \cdot P_c \cdot (1 + \mu) - Q \cdot P_r.
    \label{eq:tawarruq-cost}
\end{equation}
The effective annualised cost of the synthesised cash financing
$Q \cdot P_r$ is
\begin{equation}
    r^{\text{eff}}
    = \frac{Q \cdot (P_d - P_r)}{Q \cdot P_r \cdot (t_1 - t_0)}
    = \frac{P_c (1 + \mu) - P_r}{P_r \cdot (t_1 - t_0)}.
    \label{eq:tawarruq-effective-rate}
\end{equation}
\end{proposition}

\begin{proof}
Direct from the cash flows of Definition~\ref{def:tawarruq}: the
client receives $Q \cdot P_r$ at $t_0$ and pays $Q \cdot P_d$ at
$t_1$. The net cost is the difference; the effective rate follows
by dividing the cost by the principal and the tenor. \qed
\end{proof}

\begin{remark}[Scholarly criticism of organised \emph{tawarruq}]
\label{rem:tawarruq-criticism}
The International Islamic Fiqh Academy of the OIC has explicitly
ruled that organised \emph{tawarruq}, in which the bank
coordinates all four transactions and the third-party purchaser is
pre-arranged, is impermissible \citep{OICFiqh179}.
\citet{Usmani2002} and the majority of contemporary GCC scholars
have taken the same position, on the ground that the structure is
\emph{\d{h}iyal} designed to synthesise interest-bearing financing.
The practical usage of organised \emph{tawarruq} in Islamic banking
is defended by some Malaysian scholars and, on necessity grounds,
by \citet{AlSuwailem2006}; but the consensus AAOIFI and OIC
position remains restrictive. In institutional practice, the
markup $\mu$ is almost invariably calibrated against a prevailing
interbank benchmark, which additionally invokes the (C3)
concerns of Proposition~\ref{prop:riba}.
\end{remark}

\begin{proposition}[Contested compliance of
    $\mathcal{K}_{\text{taw}}$]
\label{prop:tawarruq-comply}
Under the organised form of Definition~\ref{def:tawarruq},
$\mathcal{K}_{\text{taw}}$ is impermissible under
\citet{OICFiqh179} and AAOIFI Shariah Standard~No.~30 \S5. Under a
non-coordinated form, in which the client independently sources
the third-party buyer and the sequence is not pre-arranged by the
bank, permissibility is contested and depends on the
\emph{ijtih\={a}d} followed. Where the markup $\mu$ is calibrated
against an interbank benchmark, (C3) is additionally violated
mechanically by Proposition~\ref{prop:riba}.
\end{proposition}

This paper does not endorse organised \emph{tawarruq}. It is
recorded here for completeness of the survey of alternative
architectures, and to enable the reader to place the third
architecture, introduced in Section~\ref{sec:alt-sharika}, in
context.

\subsection{Architecture C: \emph{Shar\={\i}kat al-'in\={a}n}
    pooled-equity market access}
\label{sec:alt-sharika}

This subsection introduces and analyses the paper's principal
constructive contribution: a market-access architecture built on the
classical Islamic contractual joint-equity partnership
(\emph{shar\={\i}kat al-'in\={a}n}, a species of
\emph{shar\={\i}kat al-'aqd}) that is formally consistent with the
operational design constraints of
Definition~\ref{def:compliant-arch} while providing the retail
client with proportional economic exposure to a chosen
underlying~$S$.

\begin{remark}[Categorical classification: \emph{shar\={\i}kat
    al-'aqd}, not \emph{shar\={\i}kat al-milk}]
\label{rem:sharikat-al-aqd}
Classical \emph{fiqh} distinguishes two categories of partnership:
\emph{shar\={\i}kat al-'aqd} (contractual partnership, formed by
mutual undertaking to pool capital and share profits) and
\emph{shar\={\i}kat al-milk} (proprietary co-ownership of a
specific asset, which may arise incidentally, e.g. through joint
inheritance). AAOIFI Shariah Standard No.~12 is expressly framed
around \emph{shar\={\i}kat al-'aqd} and its provisions on
profit-and-loss-sharing, management, and dissolution presuppose
that categorisation. The construction below is deliberately a
\emph{shar\={\i}kat al-'aqd}: a distinct legal entity (the
partnership) formed by contractual undertaking between client and
broker, into which capital is contributed, which then acquires and
holds the venture's assets through a segregated custody
arrangement. The parties' interests are proportional
beneficial-and-profit shares in the partnership itself, not direct
individual co-ownership shares in the underlying asset. This
distinction matters both for the applicability of AAOIFI
SS~No.~12 and for the qab\d{d} analysis in
Proposition~\ref{prop:sharika-comply}.
\end{remark}

\begin{definition}[\emph{Shar\={\i}kat al-'in\={a}n} market-access
    architecture]
\label{def:sharika}
The \emph{shar\={\i}kat al-'in\={a}n} market-access architecture
$\mathcal{K}_{\text{shr}}$ on underlying $S$ comprises the
following. The construction is a \emph{shar\={\i}kat al-'aqd}
(Remark~\ref{rem:sharikat-al-aqd}) supported by an independent
management-services agreement
(Remark~\ref{rem:sharika-two-contracts}).

\emph{At contract inception $t_0$:}
\begin{enumerate}
    \item Client contributes capital $K_c > 0$ and broker
    contributes capital $K_b > 0$ to a contractual partnership
    (\emph{shar\={\i}kat al-'in\={a}n}, a species of
    \emph{shar\={\i}kat al-'aqd}) formed by executed partnership
    deed. The strict positivity of both contributions is what
    makes the structure a genuine \emph{shar\={\i}kah}; the
    limiting case $K_b = 0$ is treated separately in
    Remark~\ref{rem:sharika-limit-wakala}.
    \item Using the pooled capital $K_c + K_b$, the partnership
    acquires $N$ units of the underlying at spot price
    $S_{t_0}$, where
    \begin{equation}
        N = \frac{K_c + K_b}{S_{t_0}}.
        \label{eq:sharika-N}
    \end{equation}
    \item The $N$ units are held by a segregated
    nominee/custodian acting on behalf of the partnership under
    an arrangement that satisfies the operational
    constructive-possession criteria of
    Proposition~\ref{prop:sharika-comply}. Each partner holds a
    proportional beneficial-and-profit-share interest in the
    partnership itself (fractions $K_c / (K_c + K_b)$ and
    $K_b / (K_c + K_b)$) rather than a direct individual
    co-ownership share of the underlying. Realising this in
    practice requires: (i) a custody arrangement that
    safeguards partnership assets separately from the broker's
    proprietary balance sheet, under the applicable domestic
    safe-custody regime (in Australia, the ASIC provisions
    governing custodial services, including RG~133 on funds
    management and custodial services; in the United Kingdom,
    the FCA safe-custody-assets regime, principally CASS~6),
    (ii) a bilateral partnership deed under the governing law
    that formally identifies the shares as
    \emph{shar\={\i}kat al-'in\={a}n} interests in the
    partnership, and (iii) alignment between the accounting
    treatment on the broker's books and the partnership's
    substantive economic reality. Whether any particular
    jurisdictional implementation satisfies all three is a
    matter of legal analysis outside the paper's scope; see
    Section~\ref{sec:disc-practical}.
    \item Profit-share ratio $(\pi_c, \pi_b)$ is pre-agreed with
    $\pi_c + \pi_b = 1$; per AAOIFI Shariah Standard No.~12,
    \emph{losses} must be allocated strictly in proportion to
    capital contribution,
    \begin{equation}
        \ell_c = \frac{K_c}{K_c + K_b},
        \qquad \ell_b = \frac{K_b}{K_c + K_b},
        \label{eq:sharika-loss-share}
    \end{equation}
    while \emph{profits} may be distributed under the agreed
    permissible profit-sharing ratio $(\pi_c, \pi_b)$, which
    need not equal the capital ratio.
    \item Broker acts as \emph{managing partner} of the
    \emph{shar\={\i}kah}. Any compensation for broker
    management services is contracted for and paid under a
    separate, independent management-services agreement (see
    Remark~\ref{rem:sharika-two-contracts}), not as internal
    remuneration under the partnership deed itself. Under that
    agreement the broker acts as \emph{wak\={\i}l} for a
    defined \emph{ujrah} rate $u \geq 0$ per unit time per unit
    of client capital, with no functional dependence on
    interbank benchmark rates. The cumulative \emph{ujrah}
    payable by the client over the holding period $[t_0, t_1]$
    of length $\tau = t_1 - t_0$ (in years) is
    \begin{equation}
        u^{\text{shr}}_c \coloneqq K_c \cdot u \cdot \tau.
        \label{eq:sharika-ujrah}
    \end{equation}
\end{enumerate}

\emph{At exit time $t_1$:}
\begin{enumerate}
    \setcounter{enumi}{5}
    \item Position is liquidated at spot price $S_{t_1}$; the
    total proceeds $N \cdot S_{t_1}$ are distributed to the
    parties in accordance with the pre-agreed profit-and-loss
    rules (steps~4--5).
\end{enumerate}

\noindent The eligible-underlying set is any asset that admits
transfer to the partnership's segregated custody at spot:
exchange-listed equities (subject to AAOIFI Standard No.~21
sectoral and financial-ratio screening), spot commodities held via
allocated account or physical custody, and spot FX delivered to
segregated bank accounts. Cash-settled indices, non-deliverable
forwards, volatility instruments, and futures contracts (which are
not themselves the underlying) fall outside the eligible set; a
post-physical-delivery construction on deliverable-underlying
futures is a distinct architecture that is not analysed here.
\end{definition}

\begin{remark}[Two-contract construction: partnership plus
    independent management-services agreement]
\label{rem:sharika-two-contracts}
AAOIFI Shariah Standard No.~12 does not permit fixed remuneration
for management to be embedded directly in a
\emph{shar\={\i}kah} partnership contract; it does, however,
permit a partner to provide management services under a
\emph{separately executed} contract of agency
(\emph{wak\={a}la bi'l-ujrah}) or \emph{ij\={a}ra}, under which
remuneration is payable. Definition~\ref{def:sharika} adopts this
two-contract construction:
\begin{itemize}
    \item[(a)] a \emph{partnership deed} that documents capital
    contributions, profit-and-loss shares, decision rights, and
    dissolution;
    \item[(b)] a \emph{separate management-services agreement}
    executed alongside the partnership, under which the broker
    (in the role of \emph{wak\={\i}l}) provides defined
    execution, custody-coordination, and administration services
    to the partnership for the \emph{ujrah}
    $u^{\text{shr}}_c$ of~\eqref{eq:sharika-ujrah}.
\end{itemize}
The management-services agreement must specify at least: scope
of services, fee schedule, agency duties owed to the
partnership, termination conditions, treatment of conflicts of
interest, liability for breach, whether the fee is payable
independent of venture profit, and an express undertaking not to
guarantee partnership capital. Adopting this two-contract
construction is what allows the \emph{ujrah} to be legitimate
compensation for a separately identifiable service (as required
by SS~No.~12) rather than a disguised guaranteed return on the
broker's partnership capital. It also addresses the concern that
the broker's aggregate compensation combines partner returns and
manager fees: the two flows arise from formally distinct
contracts.
\end{remark}

\begin{remark}[Limiting case: $K_b = 0$ is not a
    \emph{shar\={\i}kah}]
\label{rem:sharika-limit-wakala}
If $K_b = 0$, the broker contributes no partnership capital,
the structure of Definition~\ref{def:sharika} is not a
\emph{shar\={\i}kat al-'in\={a}n} but rather a \emph{wak\={a}la}
(agency): the client owns the entire venture, and the broker
acts purely as agent-manager for a fee. That arrangement is
also permissible under Shariah, but under a different
nominate-contract category with different governance and
disclosure requirements. The propositions in this subsection
are stated for the genuine $K_b > 0$ \emph{shar\={\i}kah} case;
the $K_b = 0$ limit is discussed as a distinct architectural
variant.
\end{remark}

We now derive the client's net payoff and establish the paper's
formal consistency claims.

\begin{remark}[Governance concentration: broker as managing
    partner, \emph{wak\={\i}l}, and custody-coordinator]
\label{rem:sharika-governance}
Definition~\ref{def:sharika} concentrates three distinct roles in a
single institution: the broker is simultaneously a capital partner
in the \emph{shar\={\i}kah} (contributing $K_b > 0$), the appointed
managing partner acting as \emph{wak\={\i}l} under the separate
management-services agreement of
Remark~\ref{rem:sharika-two-contracts} (receiving the \emph{ujrah}
$u^{\text{shr}}_c$), and, via its coordination of the segregated
custody arrangement, an intermediary whose relationship with the
custodian may allow it effective operational influence over the
partnership's assets. This concentration creates well-known
conflict-of-interest concerns: the broker as managing partner may
have discretion over trading, hedging, or liquidation decisions
whose outcomes affect its own partner return, and its coordination
role on custody exposes the partnership to operational risk if the
custody arrangement is not properly independent. Meaningful
mitigation of these conflicts requires at minimum (a) written
governance-rights allocation in the partnership deed (voting rights
on material decisions; scope of managing-partner discretion;
grounds for removal), (b) independent Shariah Supervisory Board
(SSB) oversight of managing-partner conduct per AAOIFI Governance
Standard No.~1 (\citealp{AAOIFI-GS1}), (c) genuine separation of
the custodian from the broker's balance sheet under the applicable
safe-custody regime, and (d) transparent disclosure to the client
of the broker's combined economic exposure across partnership and
management-services agreements. Whether the resulting governance
package adequately discharges the fiduciary duties expected of a
managing partner and \emph{wak\={\i}l} in classical \emph{fiqh} is
a matter for the SSB reviewing the specific implementation, not
for this paper.
\end{remark}

\begin{proposition}[\emph{Shar\={\i}kat al-'in\={a}n} payoff
    dynamics]
\label{prop:sharika-payoff}
Under Definition~\ref{def:sharika} with proportional
profit-and-loss sharing $\pi_c = \ell_c = K_c / (K_c + K_b)$, the
client's net payoff is
\begin{equation}
    \Pi^{\text{shr}}_c(t_0, t_1)
    = K_c \cdot \left( \frac{S_{t_1}}{S_{t_0}} - 1 \right)
      - u^{\text{shr}}_c.
    \label{eq:sharika-payoff}
\end{equation}
\end{proposition}

\begin{proof}
The joint venture's gross P\&L at $t_1$ is
$N \cdot (S_{t_1} - S_{t_0})$. Under proportional sharing, the
client receives fraction $K_c / (K_c + K_b)$ of this quantity.
Substituting for $N$ from~\eqref{eq:sharika-N} gives
\[
    \pi_c \cdot N \cdot (S_{t_1} - S_{t_0})
    = \frac{K_c}{K_c + K_b} \cdot \frac{K_c + K_b}{S_{t_0}}
      \cdot (S_{t_1} - S_{t_0})
    = \frac{K_c}{S_{t_0}} \cdot (S_{t_1} - S_{t_0})
    = K_c \cdot \left( \frac{S_{t_1}}{S_{t_0}} - 1 \right).
\]
Subtracting the broker's ujrah gives \eqref{eq:sharika-payoff}. \qed
\end{proof}

\begin{remark}[Economic equivalence to a spot purchase, at the
    gross level]
\label{rem:sharika-equivalent-spot}
Proposition~\ref{prop:sharika-payoff} establishes that, under
proportional profit-and-loss sharing, the client's \emph{gross}
exposure to the underlying is exactly that of a direct spot
purchase of $K_c / S_{t_0}$ units of $S$. The ujrah
$u^{\text{shr}}_c = K_c u \tau$ is the incremental cost of the
partnership structure and is bounded below by zero. At the
\emph{gross venture} level, the client and broker are on the
same side of the position and gross P\&L is shared in
proportion to capital contribution. At the \emph{net individual}
level after ujrah, small positive venture returns can leave the
client with negative net P\&L while the broker receives the
ujrah on top of its share of the venture return; the gross-level
partnership alignment is what matters for the \emph{maysir}
analysis, and the fee is separately treated as legitimate
compensation for the agency service. This is a substantively
different situation from the CFD case
(Proposition~\ref{prop:zero-sum}), where the client's gain is
the broker's loss by construction and there is no shared venture.
\end{remark}

\begin{remark}[No leverage \emph{within this architecture};
    lower bound on net payoff]
\label{rem:sharika-no-leverage}
Definition~\ref{def:sharika} does not admit leverage: the
client's effective exposure to the underlying is bounded above
by $K_c \cdot S_{t_1} / S_{t_0}$, and the client's net payoff
$\Pi^{\text{shr}}_c$ is bounded below by
$-K_c - u^{\text{shr}}_c$ (attained in the limit
$S_{t_1} \to 0$; the underlying investment loss alone is
bounded below by $-K_c$ and the accrued ujrah adds to that).
Introducing leverage into this specific joint-ownership
architecture would require either (a) debt financing of the
leveraged fraction (a \emph{qar\d{d}} with \emph{ziy\={a}dah}
concern), (b) non-ownership of the leveraged fraction (a
\emph{qab\d{d}} concern), or (c) a deferred contingent
obligation structure (a
\emph{bay' al-k\={a}li' bi'l-k\={a}li'} concern). Leverage is
therefore structurally unavailable within
Definition~\ref{def:sharika} \emph{as constructed}; this is not
a general theorem that no Shariah-compliant financing anywhere
can produce leverage, but a property of this specific
joint-ownership architecture under the paper's adopted
constraints. Section~\ref{sec:no-equivalent} addresses the CFD
use cases for which this absence of leverage is a binding
constraint.
\end{remark}

\begin{proposition}[Formal consistency of
    $\mathcal{K}_{\text{shr}}$ with the operational design
    constraints]
\label{prop:sharika-comply}
Under Definition~\ref{def:sharika} with the loss-sharing
constraint $\ell_c = K_c/(K_c+K_b)$ of clause~4, $K_b > 0$, $u$
$\mathcal{G}^{\mathcal{K}}_{s}$-measurable per (D1) below, the
two-contract construction of
Remark~\ref{rem:sharika-two-contracts}, and \emph{provided the
segregated custody arrangement of clause~3 satisfies the
operational constructive-possession criteria stated in the (D2)
step below}, $\mathcal{K}_{\text{shr}}$ satisfies each of the
operational design constraints (D1)--(D5) of
Definition~\ref{def:compliant-arch}. Whether this formal
consistency amounts to Shariah compliance in general is a
question outside the scope of the proposition and belongs to a
qualified Shariah supervisory board.
\end{proposition}

\begin{proof}
We verify each design constraint in turn.

\emph{(D1) No interest-linked consideration.} The only periodic
consideration in $\mathcal{K}_{\text{shr}}$ is the ujrah rate
$u$, which by hypothesis is measurable with respect to the
non-rate $\sigma$-algebra $\mathcal{G}^{\mathcal{K}}_{s}$ of
Definition~\ref{def:compliant-arch}, i.e. no functional
dependence on any of $r^{\text{CCY1}}_s, r^{\text{CCY2}}_s,
r^{\text{IBOR}}_s$.

\emph{(D2) Constructive possession where a sale is present.}
This step is the qab\d{d} analysis. By clause~3 of
Definition~\ref{def:sharika}, the partnership acquires the $N$
units of $S$ at $t_0$ and holds them via a segregated
nominee/custodian on behalf of the partnership. Ownership does
not by itself entail \emph{qab\d{d}} in the sense of AAOIFI
Shariah Standard No.~18; constructive possession additionally
requires legally recognised control and entitlement over the
asset. We therefore state constructive possession as a
\emph{conditional} conclusion contingent on the custody
arrangement satisfying the following operational criteria: (i)
the assets are held in a segregated custody account whose title
identifies the partnership (not the broker's proprietary
balance sheet), (ii) the partnership has an enforceable
entitlement to direct the disposal of the assets on demand,
(iii) the assets are ring-fenced from the broker's creditors in
the event of broker insolvency (as required by the applicable
safe-custody regime), and (iv) the nominee owes fiduciary
duties to the partnership rather than to the broker in its
proprietary capacity. Under any custody implementation
satisfying (i)--(iv), the partnership-level control and custody
conditions required for constructive possession at the level of
the \emph{partnership} are met at $t_0$. Whether that
partnership-level \emph{qab\d{d}} is sufficient to establish the
requisite possession for each partner's partnership interest, in
the sense demanded by AAOIFI Standard No.~18 for the relevant
transaction under the particular governing law, is a distinct
fiqh-and-legal question that this paper does not resolve; it is
left to SSB and legal review of any specific implementation
(Section~\ref{sec:disc-practical}). Whether a given
jurisdictional implementation genuinely delivers (i)--(iv) is
similarly a matter of legal analysis.

\emph{(D3) Real-economic-activity linkage.} The joint venture's
payoff derives from the acquisition of $N$ units of the
underlying asset $S$ at $t_0$ and its subsequent disposal at
$t_1$; the payoff is not a cash-settled contingent claim on a
price difference alone.

\emph{(D4) Proportional risk sharing.} By clause~4 of
Definition~\ref{def:sharika} and the proportional condition
$\pi_c = \ell_c$, gross venture P\&L is shared strictly in
proportion to capital contribution. At the gross level, both
counterparties are on the same side of the underlying position
(gains and losses accrue jointly); the ujrah $u^{\text{shr}}_c$
is a separately identifiable agency fee, so at the net
individual level after fees the parties' P\&Ls need not have
the same sign, but the gross-level partnership structure is
what the \emph{maysir} concern addresses (see
Remark~\ref{rem:sharika-equivalent-spot}).

\emph{(D5) Immediate substantive consideration.} At $t_0$, both
$K_c$ and $K_b$ are contributed in cash to the joint venture;
both contributions are non-contingent, immediate, and
irreversible absent liquidation. Substantive consideration is
therefore rendered by both parties at inception.

\emph{Parameter constraint on \emph{gharar} (C5).} Under
Definition~\ref{def:sharika}, no leverage is introduced (see
Remark~\ref{rem:sharika-no-leverage}) and consequently there is
no price-triggered maintenance-margin liquidation of the kind
analysed in Proposition~\ref{prop:leverage-sensitivity}
(contractual liquidation, redemption, or fund termination may
still occur under separately specified conditions). The
underlying is a deliverable real asset with observable spot
pricing, eliminating the specific execution-ambiguity concern
of Definition~\ref{def:slippage} under a standard
exchange-traded custody arrangement. The terminal exposure of
the underlying investment is bounded below by $-K_c$, so the
client's net payoff is bounded below by $-K_c - u^{\text{shr}}_c$.

Combining the five verifications gives the stated formal
consistency. \qed
\end{proof}

\paragraph{Restriction on eligible underlyings.}
Definition~\ref{def:sharika} requires that $N$ units of the
underlying $S$ be actually deliverable to the partnership's
segregated custody. This restricts the eligible underlying set
to instruments admitting such transfer:
\begin{itemize}
    \item Spot equities via nominee custody (with additional
    Shariah-compliance screening on the issuer under AAOIFI
    Standard~No.~21);
    \item Spot commodities via warrant, allocated account, or
    physical custody;
    \item Spot FX via delivery to segregated bank accounts.
\end{itemize}
It excludes: cash-settled indices (S\&P~500, DAX, VIX);
non-deliverable forwards on restricted currencies; futures
contracts (which are not themselves the underlying deliverable
asset); and synthetic instruments such as variance swaps,
dispersion products, and volatility indices. A distinct
post-physical-delivery construction on deliverable-underlying
futures is imaginable but is a separate architecture not analysed
here. Section~\ref{sec:no-equivalent} addresses the resulting
design gap.

\subsection{Comparative architecture map}
\label{sec:arch-comparison}

Table~\ref{tab:architecture-comparison} summarises the three
alternative architectures against a common set of dimensions.

\begin{table}[!ht]
\centering
\small
\caption{Comparative properties of the three alternative
architectures. ``Effective leverage'' refers to notional exposure
per unit of client capital; ``scholarly status'' summarises the
consensus recorded in
Sections~\ref{sec:alt-wad}--\ref{sec:alt-sharika}.}
\label{tab:architecture-comparison}
\begin{tabular}{@{}p{2.8cm}p{3.2cm}p{2.6cm}p{2cm}p{2.8cm}@{}}
\toprule
\textbf{Architecture}
    & \textbf{Use case served}
    & \textbf{Eligible underlyings}
    & \textbf{Effective leverage}
    & \textbf{Scholarly status} \\
\midrule
$\mathcal{K}_{\text{wa'd}}$
    & FX forward exposure; commercial hedging
    & Deliverable FX pairs
    & 1$\times$
    & Contested (BNM permits, GCC/AAOIFI restrictive) \\[6pt]
$\mathcal{K}_{\text{taw}}$
    & Cash financing synthesis
    & LME-listed commodities
    & N/A (financing, not exposure)
    & Impermissible in organised form (OIC 179); contested otherwise \\[6pt]
$\mathcal{K}_{\text{shr}}$
    & Spot directional exposure; buy-and-hold
    & Deliverable equities, spot FX, spot commodities
    & 1$\times$
    & Formally consistent with (D1)--(D5)
      (Prop.~\ref{prop:sharika-comply}); subject to SSB review \\
\bottomrule
\end{tabular}
\end{table}

The comparison exposes two structural features common to all three
candidate architectures: (i) none admits leverage in excess of
$1\times$ on the client's capital, and (ii) all require
deliverable, real underlyings and are therefore incompatible with
purely synthetic instruments. The economic cost of these features
against the conventional leveraged CFD baseline is quantified in
Section~\ref{sec:simulation}.

\subsection{CFD use cases with no permissible equivalent}
\label{sec:no-equivalent}

We conclude the section by stating honestly which retail CFD use
cases have \emph{no} permissible equivalent under the constraint
set $\CC_{\text{Shariah}}$. This is a finding, not a limitation of
the paper: the incompatibility is structural, and identifying it
precisely is itself a research contribution.

The following use cases have no permissible substitute in
Definition~\ref{def:compliant-arch}:
\begin{enumerate}
    \item \emph{Leveraged short-horizon directional speculation.}
    Any use case requiring effective exposure in excess of $1\times$
    client capital is precluded by
    Remark~\ref{rem:sharika-no-leverage}: leverage is structurally
    incompatible with the joint constraint set (C3), (C6), (C7).
    \item \emph{Cash-settled synthetic index exposure.} Any use
    case requiring exposure to a non-deliverable underlying (index,
    volatility measure, variance swap) is precluded by
    Definition~\ref{def:sharika} clauses~2--3, which require an
    actually-deliverable and jointly-ownable underlying.
    \item \emph{Short-selling of non-owned assets.} The classical
    short sale, selling an asset one does not own with the
    intention of buying it back later at a lower price,
    violates Standard~No.~18 (\emph{qab\d{d}}) and cannot be
    reconstructed under the joint-equity framework.
    \item \emph{Pure intraday price-difference wagers on
    unbacked references.} Retail contracts that transact on a
    reference price with no underlying commercial rationale and no
    contemplation of delivery, often marketed as ``binary
    options'' or ``exotic CFDs'', lie outside the domain of any
    of the three architectures.
\end{enumerate}

The paper's position on these use cases is that their
unavailability under the joint-ownership architecture of
Definition~\ref{def:sharika} is a feature of that specific
construction rather than a limitation. The retail-outcome
evidence surveyed in Section~\ref{sec:motivation}, retail
CFD loss rates of 74--89 per cent across major jurisdictions,
is associated with substantial realised losses among retail
participants in these product categories. Reasonable readers
may reach different normative conclusions about the wider
policy implications; the paper records the empirical evidence
and its convergence with the leverage-cap feature of the
proposed architecture without extrapolating to broader
normative claims.

Section~\ref{sec:simulation} quantifies the economic cost of the
candidate architectures, principally
$\mathcal{K}_{\text{shr}}$, against a conventional leveraged
CFD baseline, under illustrative Bates parameterisations.

\section{Quantitative Comparison via Monte Carlo Simulation}
\label{sec:simulation}

Sections~\ref{sec:cfd-mechanics}--\ref{sec:alternatives} established
the theoretical basis for the paper's central claim: the standard
retail CFD exhibits five structural concerns under the adopted
AAOIFI/OIC constraint framework, and a
\emph{shar\={\i}kat al-'in\={a}n} pooled-equity architecture is
formally consistent with the operational design constraints
(D1)--(D5) of Definition~\ref{def:compliant-arch} subject to SSB
review, at the cost of eliminating leverage. This section
quantifies the economic cost of that trade-off under a common
jump-diffusion price process, using a Monte Carlo simulation
across four underlyings, four holding horizons, and four leverage
regimes.

\emph{Framing note.} The comparison in this section is
\emph{not} a claim that the \emph{shar\={\i}kat} architecture is
intrinsically superior to a leveraged CFD as a financial product:
the two architectures deliver different economic exposures on
identical client capital and are therefore not directly comparable
on a like-for-like basis. Concretely, a client with capital $K_c$
takes exposure $K_c \cdot L / S_{t_0}$ units of the underlying
under a CFD at leverage $L$, versus $K_c / S_{t_0}$ units under the
\emph{shar\={\i}kat} at compulsory unit leverage. The comparison
should therefore be read as a quantification of the
\emph{compliance cost} of the Shariah-imposed leverage cap:
relative to a matched-capital leveraged CFD, the compliant
architecture forgoes proportional leveraged upside and leveraged
downside, and eliminates margin-call risk. Whether this trade-off
is desirable for a given participant depends on their risk
preferences and religious commitments, not on any intrinsic
superiority claim. All figures and tables in this section derive
from
the reference implementation released with the paper
(Section~\ref{sec:code-availability}).

\subsection{Simulation design}
\label{sec:sim-design}

The simulation is organised as a $4 \times 4 \times 4$ grid over
underlyings, holding horizons, and leverage regimes, with
$n = 10{,}000$ paths per cell.

\paragraph{Underlyings.} Three major FX pairs (EUR/USD, USD/JPY,
AUD/USD) and one equity-index proxy (S\&P~500) span the volatility
range typical of retail CFD flow. Each is modelled under the Bates
(\citeyear{Bates1996}) jump-diffusion of
Assumption~\ref{ass:bates}; illustrative parameter calibrations,
consistent with published FX and equity implied-volatility surfaces
\citep{Andersen2009,Bates1996}, are recorded in
Table~\ref{tab:calibration}.

\begin{table}[!ht]
\centering
\small
\caption{Bates model calibration presets used in the Section~6
simulations. Long-run volatility
$\sigma_{\text{LR}} = \sqrt{\theta}$; the initial variance is set
to $v_0 = \theta$ (spot start at long-run equilibrium); the drift
is set to $\mu = 0$ for all underlyings to isolate the compliance
cost from directional views; the mean-reversion speed is
$\kappa_v = 2.0$~yr$^{-1}$ and the vol-of-vol is
$\sigma_v = 0.3\sqrt{\theta}$ for all underlyings, chosen to
satisfy the strict Feller condition
$2\kappa_v\theta > \sigma_v^{2}$. FX interbank overnight rates
$r^{\text{CCY1}}$, $r^{\text{CCY2}}$ used in the CFD Tom-Next
computations are the illustrative fixed values SOFR $= 5.25\%$,
ESTR $= 3.50\%$, SONIA $= 5.00\%$, TONA $= 0.10\%$, RBA cash rate
$= 4.35\%$, held constant across the simulation horizon. The
S\&P~500 case in Section~\ref{sec:sim-results} is included as a
mathematical counterfactual for isolating the compliance-cost
differential across asset-class volatility regimes: a
Shariah-compliant equity implementation would use a Standard
No.~21-screened basket or ETF rather than the index itself.
Complete numerical values are also machine-readable in
\texttt{src/bates.py} of the reference implementation.}
\label{tab:calibration}
\begin{tabular}{@{}lrrrrr@{}}
\toprule
\textbf{Underlying}
    & $\sigma_{\text{LR}}$ & $\rho$
    & $\lambda$ (jumps/yr) & $\mu_J$ & $\sigma_J$ \\
\midrule
EUR/USD & 8\%  & $-0.10$ & 1.5 & $-0.002$ & 0.012 \\
USD/JPY & 9\%  & $+0.05$ & 2.0 & $+0.001$ & 0.015 \\
AUD/USD & 11\% & $-0.20$ & 2.5 & $-0.003$ & 0.018 \\
S\&P~500\textsuperscript{$\dagger$} & 18\% & $-0.65$ & 3.0 & $-0.015$ & 0.035 \\
\bottomrule
\end{tabular}
\end{table}

\paragraph{Holding horizons.} Four horizons span the retail
trading spectrum, expressed in trading days:
1~day (intraday exit), 5~days (one trading week), 21~days
(approximately one trading month), and 63~days (one trading
quarter). All time-in-year conversions in Section~6 use the
trading-day convention $\tau_{\text{sim}}(h) = h / 252$; this
differs from the calendar/actual-365 convention $\tau = 1/365$
used in the Tom-Next derivation of
Section~\ref{sec:tomnext}, which follows industry practice
for overnight financing. The simulation thus mixes conventions
across two logically distinct operations: the Bates path
generator and the \emph{shar\={\i}kat} ujrah accrual use
trading-day-per-year, while the CFD Tom-Next fee accrual uses
calendar-day-per-year. Both are correctly implemented in the
reference code (see \texttt{src/simulate.py}) and the mixing has
been checked to leave the Section~6 comparisons unaffected in
direction, though absolute per-day figures differ between the
two operations by the ratio $252/365 \approx 0.69$; the
\emph{shar\={\i}kat} 63-day ujrah figure of approximately
USD~25 in Figure~\ref{fig:cost-of-carry} follows from
$K_c \cdot u \cdot \tau_{\text{sim}}(63) = 10{,}000 \cdot 0.01
\cdot (63/252) = 25$.

\paragraph{Leverage regimes.} Four regimes are considered:
$L \in \{1, 5, 10, 30\}$. The maximum $L = 30$ corresponds to the
ASIC 2020 product-intervention cap on retail \emph{major-FX} CFDs
specifically (\citealp{ASIC2020}); the same 30$\times$ ceiling
applies to retail major-FX CFDs under the analogous ESMA
(\citealp{ESMA2018}) and FCA (\citealp{FCA2022}) product
intervention frameworks. Lower caps apply to other retail CFD
categories under all three regimes (typically 20$\times$ for
non-major FX and gold, 10$\times$ for other commodities and major
equity indices, 5$\times$ for individual equities, and 2$\times$
for cryptocurrencies). The $L = 30$ regime in this simulation
should therefore be read as the ceiling that applies to the FX
subset of the underlyings studied here; the same regime applied
to the S\&P~500 exceeds what any of the three regulators permit
for retail index CFDs, and is included in the grid as a stress
case rather than as a representation of permitted retail exposure.

\paragraph{Contract parameters.} Client capital is fixed at
$K_c = \text{USD}\,10{,}000$. The CFD baseline uses a maintenance
margin coefficient $\alpha = 0.5$, and stylised common baseline
transaction-cost assumptions of a 2~bps bid-ask spread per side
and a 100~bps per annum broker markup on the interbank
differential; these are illustrative values applied uniformly
across the four underlyings for comparability and are not intended
as representative universal retail-CFD pricing (FX, equity
indices, and single-name equities have materially different
spread and funding structures in practice). The
\emph{shar\={\i}kat al-'in\={a}n}
architecture uses a client-side ujrah of 100~bps per annum, flat
per unit of holding period, independent of interbank rates.
Interbank rate presets $(r^{\text{CCY1}}, r^{\text{CCY2}})$ are
set to representative 2025 levels for each pair.

\paragraph{Numerical scheme.} The Bates system is discretised by
Euler--Maruyama on the log-price with full-truncation of the
variance process to preserve positivity, following
\citet{Andersen2009}. All path simulations use fixed seeds derived
from a session-wide base seed, so the reported results are
exactly reproducible via the released code.

\subsection{Metrics}
\label{sec:sim-metrics}

For each grid cell we compute:
\begin{itemize}
    \item \emph{All-in holding cost:} the sum of expected
    financing and per-entry-and-exit transaction costs (spread
    plus commission) over the holding period, in USD. The
    naming distinguishes this quantity from cost-of-carry in
    its strict sense, since the spread component accrues at
    entry and exit rather than continuously.
    \item \emph{Margin-call trigger probability:} the empirical
    frequency $\PP(\tau_{\text{MC}} < t_1)$ estimated as the
    fraction of paths that cross the barrier
    of~\eqref{eq:mc-explicit-barrier} prior to the horizon.
    \item \emph{Return-distribution moments:} mean, standard
    deviation, skewness, and excess kurtosis of the net client
    P\&L.
    \item \emph{Tail-risk measures:} Value-at-Risk and
    Conditional Value-at-Risk at the 95\% level; per-path
    maximum drawdown.
    \item \emph{Risk-adjusted performance:} annualised Sharpe
    and Sortino ratios.
    \item \emph{Loss frequency:} the empirical probability of a
    negative net client P\&L over the holding period.
\end{itemize}
\emph{Note on the equity-index case.} The Tom-Next financing
formula of Section~\ref{sec:tomnext-fx} is derived under FX
covered interest rate parity, with $r^{\text{CCY1}}$ and
$r^{\text{CCY2}}$ interpreted as base- and quote-currency
overnight rates. For the S\&P~500 case in the simulation grid,
the natural financing structure is instead
$r^{\text{funding}} - q^{\text{dividend}}$, where
$r^{\text{funding}}$ is the client's funding rate on the
leveraged position and $q^{\text{dividend}}$ is the index
dividend yield. In the simulation we take
$r^{\text{funding}} \approx r^{\text{USD}}$ and
$q^{\text{dividend}} = 0$ as an illustrative simplification
solely to isolate the pure leverage effect from dividend-carry
effects; a realistic S\&P~500 dividend yield is of order 1.3\%,
so a full implementation would apply
$r^{\text{funding}} - q^{\text{dividend}}$ with that value.
The simplification does not affect the direction of the
compliance-cost comparison but does reduce the absolute CFD
financing cost on the equity leg by that fraction. These are
not empirical calibrations and produce financing
figures broadly comparable in magnitude to a USD-quoted FX
pair, which is sufficient for the compliance-cost comparison
that is the section's object. A production implementation would
require an equity-specific financing derivation, deferred to
subsequent work.

\subsection{Statistical inference}
\label{sec:sim-inference}

Percentile bootstrap 95\% confidence intervals are computed for
all mean-P\&L estimates using 500 bootstrap resamples per cell.
For the binary margin-call-probability estimate at each cell,
Monte Carlo standard error is approximately
$\sqrt{\hat{p}(1-\hat{p})/n}$ with $n = 10{,}000$; a
representative point estimate $\hat{p} = 0.48$ carries a
standard error of approximately 0.005 (i.e.\ half a
percentage point). Bootstrap confidence intervals on the tail
metrics VaR$_{95}$ and CVaR$_{95}$ are computed in the code but
are not tabulated here for brevity; they are of order
$\pm 3\%$--$5\%$ of the point estimates and do not affect the
qualitative conclusions. The
Kolmogorov--Smirnov two-sample statistic is reported as a
simulation-based distributional-distance measure between CFD and
\emph{shar\={\i}kat} return distributions; a small KS value
indicates limited simulation-based distributional separation but is
not, and is not interpreted as, formal evidence of
distributional equality. $p$-values are adjusted for multiple
comparisons using the Benjamini--Hochberg false-discovery-rate
procedure at $q = 0.05$ (\citealp{Benjamini1995}).

\subsection{Machine learning extension: what drives margin-call
    events?}
\label{sec:sim-ml}

To characterise the price-path features that drive CFD tail losses
we train an XGBoost gradient-boosted classifier
(\citealp{ChenGuestrin2016}) on path-level features to predict
whether a given path will trigger a margin call over a 21-day
horizon at $L = 10$ on the S\&P~500 underlying. Features are
constructed from early-path information, first-day log return,
absolute log return, and realised variance; 3-day log return and
mean variance, so the classifier answers the operationally
relevant question of whether a margin call is predictable from
early observations of the path.

The classifier achieves out-of-sample area under the ROC curve
of $\text{AUC} = 0.728$ on a 25\% test partition (class balance
27.7\% margin-called). Figure~\ref{fig:feature-importance}
displays the gain-based feature importances; the 3-day
cumulative log return carries the largest gain in the model,
followed by 3-day mean variance, with first-day features
contributing marginally. This is model-consistency evidence:
within the stylised simulated environment, a simulated
margin-call event is partially predictable from simulated early
returns and realised variance, and the classifier's largest
gains come from features aggregating information over several
days rather than any single-day realisation. The result does
not identify a live-market empirical relationship; the exercise
demonstrates conditional predictability of a simulated
barrier-crossing event from simulated early-path information,
consistent with the barrier structure of
Proposition~\ref{prop:leverage-sensitivity}.

\subsection{Results}
\label{sec:sim-results}

\paragraph{Cost of carry.} Figure~\ref{fig:cost-of-carry} shows
the expected cost of carry for EUR/USD over the four horizons at
leverages $L \in \{5, 10, 30\}$ for the CFD, alongside the flat
ujrah cost of the \emph{shar\={\i}kat} architecture. The CFD's
cost scales approximately linearly in leverage and in horizon; at
$L = 30$ and a 63-day horizon it reaches USD~1{,}220, compared to
USD~25 for the \emph{shar\={\i}kat} at the same horizon
(a factor of approximately 49$\times$).

\begin{figure}[!ht]
\centering
\includegraphics[width=0.72\linewidth]{figures/fig1_cost_of_carry.pdf}
\caption{Expected cost of carry as a function of holding horizon
for the EUR/USD underlying under the standard retail CFD at three
leverage regimes and under the \emph{shar\={\i}kat al-'in\={a}n}
architecture (which is leverage-invariant). Both axes are
logarithmic. The \emph{shar\={\i}kat} cost line is unique, whereas
the CFD produces one line per leverage regime. The 100-bps annual
ujrah accrues linearly with horizon; the CFD cost accrues linearly
with the leveraged notional, giving the two lines' different
elevations.}
\label{fig:cost-of-carry}
\end{figure}

\paragraph{Margin-call frequency.}
Figure~\ref{fig:margin-call} presents the margin-call trigger
probability $\PP(\tau_{\text{MC}} < t_1)$ across the full grid.
Three qualitative features emerge:
(i) at $L = 1$, no margin calls were observed in any of the
10{,}000 simulated paths (the barrier lies approximately 50\%
below the entry price and, under the parameters of
Table~\ref{tab:calibration}, an excursion of that magnitude is
extremely unlikely within the horizons studied, though not
mathematically impossible);
(ii) at $L = 30$ and horizon 63 days, margin-call frequency
exceeds 60\% for every underlying and reaches 74.5\% for the
S\&P~500 (the highest-volatility underlying);
(iii) the S\&P~500 dominates every other underlying at every
$(L, h)$ combination beyond the first row/column, driven by its
higher long-run volatility and stronger leverage effect $\rho$.

\begin{figure}[!ht]
\centering
\includegraphics[width=\linewidth]{figures/fig2_margin_call.pdf}
\caption{Margin-call trigger probability
$\PP(\tau_{\text{MC}} < t_1)$ estimated over $n = 10{,}000$
Monte Carlo paths per grid cell, across the four underlyings of
Table~\ref{tab:calibration}. Cell shading indicates the estimated
probability; numerical annotations give the point estimate.}
\label{fig:margin-call}
\end{figure}

\paragraph{Return-distribution comparison.}
Figure~\ref{fig:return-cdf} contrasts the empirical CDFs of net
client P\&L for the AUD/USD underlying at a 21-day horizon.
The \emph{shar\={\i}kat} CDF is narrowly concentrated around
zero, reflecting the unleveraged exposure to spot AUD/USD
returns over one month. The leveraged CFD CDF at $L = 10$ is
markedly wider, with pronounced left-tail mass driven by
margin-call events. Empirical 95\% VaR statistics from the full
grid, reported in Table~\ref{tab:tail-metrics}, formalise this
visual comparison.

\begin{figure}[!ht]
\centering
\includegraphics[width=0.72\linewidth]{figures/fig3_return_cdf.pdf}
\caption{Empirical cumulative distribution functions of net
client P\&L: leveraged CFD (orange) at $L = 10$ versus the
\emph{shar\={\i}kat al-'in\={a}n} architecture (blue) at $L = 1$
on AUD/USD over a 21-day horizon. Sample size
$n = 20{,}000$ paths per architecture.}
\label{fig:return-cdf}
\end{figure}

\paragraph{Tail-risk decomposition: jumps vs.\ diffusion.}
To quantify the contribution of the jump component to CFD tail
losses, Figure~\ref{fig:tail-decomp} compares the left tail
(below the 20th percentile) of the CFD P\&L distribution on
S\&P~500 at $L = 10$ under the full Bates model against the
same simulation with the jump intensity set to zero
($\lambda = 0$). The full-model tail exhibits a pronounced spike
at the negative-balance-protection floor (loss of USD~10{,}000,
the initial margin), which the diffusion-only counterpart does
not produce.

\begin{figure}[!ht]
\centering
\includegraphics[width=0.72\linewidth]{figures/fig4_tail_decomposition.pdf}
\caption{Left-tail distribution of net CFD client P\&L on the
S\&P~500 underlying at $L = 10$ and 21-day horizon, comparing
the full Bates model (with jumps) against the diffusion-only
counterpart (no jumps). Both simulations use the same random
seeds and identical Brownian innovations, differing only in
the presence of the Poisson jump component. The spike at
$-10{,}000$~USD (initial margin) under the full model reflects
margin-call events that discontinuous price movements
contribute to; the no-jump distribution does not produce a
comparable spike over the same paths, indicating that
jumps make a substantial contribution to the extreme left tail
under the illustrative parameterisation used, without implying
that jumps are the sole cause of extreme losses.}
\label{fig:tail-decomp}
\end{figure}

\paragraph{Machine-learning feature importances.}

\begin{figure}[!ht]
\centering
\includegraphics[width=0.72\linewidth]{figures/fig5_feature_importance.pdf}
\caption{XGBoost feature importances for the margin-call
prediction task described in Section~\ref{sec:sim-ml}. Test-set
AUC of $0.728$ on the S\&P~500, 21-day horizon, $L = 10$
configuration. Under the classifier trained on this stylised
simulation, the 3-day cumulative log return carries the
largest gain-based feature importance, followed by 3-day mean
variance. Values reflect the classifier's use of the features
within the simulated environment; they are model-specific
importances, not causal or empirical structural discoveries.}
\label{fig:feature-importance}
\end{figure}

\paragraph{Tabular summary.} Table~\ref{tab:tail-metrics}
consolidates the tail-risk metrics across the four underlyings
at the 63-day horizon. The \emph{shar\={\i}kat} 95\%~VaR does
not vary with the CFD leverage parameter $L$ (the
\emph{shar\={\i}kat} has no leverage parameter), and ranges from
USD~687 (EUR/USD) to USD~1{,}554 (S\&P~500 counterfactual, see
Table~\ref{tab:calibration} footnote) driven by the underlying
volatilities. The CFD 95\%~VaR grows approximately
linearly in leverage and reaches the initial-margin floor at
$L = 30$ for the S\&P~500. Loss frequency for the CFD remains
near, but need not equal, 50\% at unit leverage; the observed
$L=1$ figures (approximately 45--57\% across the four
underlyings) reflect return skewness, jump asymmetry, financing
carry and transaction costs, none of which the zero-drift
condition eliminates. The frequency rises to 69--76\% at $L=30$
(dominated by margin-call events), whereas the
\emph{shar\={\i}kat} loss frequency remains close to 50\%
across all leverage regimes (again, because the
\emph{shar\={\i}kat} has no leverage regime).

\begin{table}[!ht]
\centering
\small
\caption{Tail-risk and loss-frequency metrics at the 63-day
holding horizon across the four underlyings. VaR reported in USD
(client capital $K_c = 10{,}000$). \emph{Shar\={\i}kat} figures
are shown once per underlying since they are leverage-invariant.
Simulation size: $n = 10{,}000$ paths per cell.
\textsuperscript{$\dagger$}The \emph{shar\={\i}kat} S\&P~500 row
is a mathematical counterfactual benchmark used to isolate the
effect of the underlying's volatility regime; cash-settled
indices are outside the eligible-underlying set of
Definition~\ref{def:sharika}, so this row does not represent an
eligible implementation, which would require a
Shariah-screened deliverable equity basket or equivalent
eligible underlying.}
\label{tab:tail-metrics}
\begin{tabular}{@{}llrrrrr@{}}
\toprule
\textbf{Underlying}
    & \textbf{Architecture}
    & $L$
    & $\PP(\tau_{\text{MC}} < t_1)$
    & VaR$_{95\%}$
    & P(loss)
    & CVaR$_{95\%}$ \\
\midrule
EUR/USD & \emph{shar\={\i}kat} & 1  & --      & 687   & 0.52 & 917 \\
        & CFD                  & 5  & 1.5\%   & 3{,}689 & 0.57 & 4{,}589 \\
        & CFD                  & 10 & 18.7\%  & 5{,}882 & 0.57 & 6{,}217 \\
        & CFD                  & 30 & 61.1\%  & 8{,}080 & 0.69 & 8{,}867 \\[3pt]
USD/JPY & \emph{shar\={\i}kat} & 1  & --      & 742   & 0.52 & 995 \\
        & CFD                  & 5  & 2.4\%   & 3{,}379 & 0.45 & 4{,}402 \\
        & CFD                  & 10 & 22.5\%  & 5{,}201 & 0.46 & 5{,}529 \\
        & CFD                  & 30 & 65.5\%  & 7{,}182 & 0.67 & 8{,}052 \\[3pt]
AUD/USD & \emph{shar\={\i}kat} & 1  & --      & 951   & 0.52 & 1{,}252 \\
        & CFD                  & 5  & 7.0\%   & 5{,}255 & 0.53 & 5{,}492 \\
        & CFD                  & 10 & 31.4\%  & 6{,}130 & 0.55 & 6{,}630 \\
        & CFD                  & 30 & 69.9\%  & 8{,}812 & 0.74 & 9{,}589 \\[3pt]
S\&P~500\textsuperscript{$\dagger$} & \emph{shar\={\i}kat} & 1  & --     & 1{,}554 & 0.43 & 2{,}140 \\
         & CFD                 & 5  & 23.2\% & 6{,}107 & 0.49 & 6{,}582 \\
         & CFD                 & 10 & 48.2\% & 7{,}578 & 0.58 & 8{,}479 \\
         & CFD                 & 30 & 74.5\% & 10{,}000 & 0.76 & 10{,}000 \\
\bottomrule
\end{tabular}
\end{table}

\paragraph{Kolmogorov--Smirnov distributional tests.} The KS
two-sample statistic between CFD and \emph{shar\={\i}kat} return
distributions ranges from 0.02--0.09 at $L = 1$ (distributions
statistically indistinguishable in economic magnitude, though
the very large sample sizes yield tiny $p$-values in most cells)
to 0.47--0.75 at $L = 30$ (distributions substantially different
in both location and shape, with $p$-values below machine
precision). The magnitude of the KS statistic thus tracks the
leverage regime almost monotonically: at $L = 1$ the
architectures are near-equivalent in their return distribution;
at $L = 30$ they are fundamentally different objects.

\subsection{Interpretation: the shape of the compliance-cost curve}
\label{sec:sim-interpretation}

We reiterate the framing note at the start of this section: the
CFD and the \emph{shar\={\i}kat} deliver different economic
exposures on identical client capital, so the results below are
interpreted as characterising the shape of the compliance-cost
curve, not as a superiority claim.

\emph{Result 1: at $L = 1$, the two architectures produce
closely similar simulated return distributions under unit
exposure.} The small KS statistics at $L = 1$
(0.02--0.09 across all underlyings and horizons) indicate
limited simulation-based distributional separation between the
\emph{shar\={\i}kat} at unit leverage and an unlevered CFD.
The two contracts are structurally distinct, the CFD is
cash-settled synthetic exposure, the \emph{shar\={\i}kat} is
proportional joint ownership, but at unit exposure their
simulated payoff distributions are close. They differ in
all-in holding cost structure (broker ujrah versus CFD spread
plus financing). At this leverage regime the compliance cost is
minimal.

\emph{Result 2: the cost-of-carry crossover at approximately
$L \geq 5$ on EUR/USD at any horizon above one week.} Because the
CFD's spread and financing scale with the leveraged notional
while the \emph{shar\={\i}kat} ujrah does not, the CFD's expected
carrying cost overtakes the \emph{shar\={\i}kat}'s at moderate
leverage. This is a genuine efficiency observation about the fee
structure: even a client indifferent to religious constraint would
find the CFD costly to hold across settlement cutoffs at
$L \geq 5$ over multi-day horizons, because the spread--financing
combination compounds. Note, however, that the client is receiving
$L\times$ the exposure for that fee, so the per-unit-of-exposure
cost comparison is more nuanced than the total-cost line suggests.

\emph{Result 3: exposure-driven tail-risk differential at $L \geq
10$.} At $L = 10$ on the S\&P~500 over 63 days, the CFD has a 48\%
probability of hitting the margin-call barrier and a 95\%~VaR of
USD~7{,}578, versus a \emph{shar\={\i}kat} VaR of USD~1{,}554. The
approximately five-fold differential in tail risk reflects two
compounding effects: the CFD delivers 10$\times$ the notional
exposure, and it carries a maintenance-margin close-out mechanism
that the unleveraged \emph{shar\={\i}kat} structurally does not.
The first is a direct consequence of the different exposures; the
second is architecture-specific and represents a genuine risk
feature (not merely a scaling of exposure). Note that the
absence of a margin-call mechanism in the \emph{shar\={\i}kat}
architecture does not mean the client is exposed to no forced-sale
risk of any kind: the partnership deed will typically specify
predefined exit or wind-down conditions (loss thresholds, term
limits, mutual-consent liquidations), and pooled physical assets
may need to be sold to satisfy withdrawal requests. What is
absent is the leverage-driven, price-triggered close-out that
occurs when maintenance-margin buffer is exhausted, which is a
specific mechanism proper to leveraged margin trading and is the
source of the tail-risk asymmetry documented in
Section~\ref{sec:sim-results}.
At $L = 30$ the CFD
saturates the negative-balance protection floor with probability
75\%, producing a large point mass in the return distribution at
the model-imposed maximum-loss threshold.

The three results together characterise the compliance-cost trade
under the Shariah leverage cap: a Muslim participant who chooses
the compliant architecture over a leveraged CFD accepts
proportionally reduced exposure (both upside and downside), pays
a flat ujrah in place of leverage-scaled financing, and eliminates
the maintenance-margin close-out mechanism. Whether this trade is desirable for a given
participant depends on their risk preferences, capital, and
religious commitments. What the results \emph{do} demonstrate is
that the trade is quantitatively well-defined and that the
compliant architecture, at holding horizons of one week or
longer, achieves risk-adjusted properties (lower tail risk per
dollar of client capital, no price-triggered maintenance-margin
liquidation) that a secular prudential-regulatory framework has
independently addressed through leverage caps and negative-balance
protection (\citealp{HeimerSimsek2019}; \citealp{ASIC2020}).

\subsection{Sensitivity analysis}
\label{sec:sim-sensitivity}

To verify robustness, we vary the three most influential model
parameters, jump intensity $\lambda$, vol-of-vol $\sigma_v$,
and the correlation $\rho$ between the price and variance
Brownian innovations, over the ranges
$\lambda \in \{0, 1, 3, 5\}$ per year, $\sigma_v \in \{0.1, 0.2,
0.4\}$, and $\rho \in \{-0.8, -0.4, 0, +0.4\}$, on the reference
S\&P~500 configuration at $L = 10$, 21-day horizon. The
qualitative pattern of Section~\ref{sec:sim-results} is
preserved across the sensitivity grid: the CFD's tail-risk
profile worsens monotonically in $\lambda$, mildly in
$\sigma_v$, and increasingly as $\rho$ becomes more negative
(the equity-style leverage effect that amplifies downside).
The \emph{shar\={\i}kat} architecture's return distribution
naturally shifts with the same parameters, it depends on
$S_{t_1}$ which depends on $(\lambda, \sigma_v, \rho)$, but
its \emph{margin-call} and \emph{maximum-loss} properties are
invariant to the leverage parameter $L$ (which is fixed at
unity by construction). This is a leverage-invariance property,
not an underlying-process-invariance property; both should be
kept clearly distinguished. Full sensitivity tables are
reproduced in the code repository
(Section~\ref{sec:code-availability}).

\subsection{Code and reproducibility}
\label{sec:code-availability}

A reference Python implementation supporting every figure,
table, and numerical claim in this section is released with the
paper under the MIT License. The repository is hosted at
\begin{center}
\url{https://github.com/Muhammad-Taha-Asif/Shariah-CFE-1}
\end{center}
and the tagged initial release (v0.1.0) is archived at Zenodo
with the persistent identifier
\begin{center}
DOI: \href{https://doi.org/10.5281/zenodo.22998850}%
     {\texttt{10.5281/zenodo.22998850}}
\end{center}
The repository includes: (i) the Bates model simulator; (ii) the
CFD and \emph{shar\={\i}kat} contract mechanics; (iii) the
metrics and statistical-inference modules; (iv) the XGBoost
classifier and sensitivity-analysis code; and (v) the
figure-generation code that produces the PDFs used in this paper.
All random-number seeds are fixed for exact reproducibility.
Dependencies are listed in the repository's
\texttt{requirements.txt}; the code runs on Python~3.11 and later
with the standard scientific stack (NumPy, SciPy, pandas,
scikit-learn, XGBoost, matplotlib).

\section{Discussion and Limitations}
\label{sec:discussion}

This section situates the paper's contributions against their
proper limits: what the results say, what they do not say, and
what an implementation of the proposed architectures would require
beyond the formal analysis conducted here.

\subsection{Scope of the paper: what it does and does not}
\label{sec:disc-scope}

The paper makes four positive claims, delivered in
Sections~\ref{sec:cfd-mechanics}--\ref{sec:simulation}
respectively: (i) a formal mathematical characterisation of the
standard retail CFD architecture under a Bates jump-diffusion;
(ii) formal propositions documenting five structural
incompatibilities or concerns with the adopted AAOIFI/OIC
constraint framework;
(iii) construction of three candidate alternative architectures,
\emph{wa'd}-based, \emph{tawarruq}-based, and
\emph{shar\={\i}kat al-'in\={a}n}, with the third established
as formally consistent with the operational design constraints
(D1)--(D5), subject to SSB review; and (iv) a Monte Carlo
simulation quantifying the compliance cost of the candidate
\emph{shar\={\i}kat} architecture against the conventional
leveraged CFD baseline.

The paper also makes explicit non-claims that bound its
interpretation. It is \emph{not} a work of fiqh: no proposition
here constitutes a scholarly ruling on the permissibility of any
existing or proposed instrument. It is not a claim about Islamic
finance more broadly: the analysis is restricted to retail
CFD/FX and the alternative architectures for equivalent
market-access use cases. It is not a claim about any specific
commercial implementation: the compliance status of a real
product depends on custody arrangements, execution mechanics,
Shariah supervisory review, and jurisdictional regulation,
none of which are within the paper's scope. And it is not an
empirical calibration exercise: the Bates parameter presets used
in Section~\ref{sec:simulation} are illustrative values
consistent with published FX and equity volatility studies, not
fits to any proprietary dataset.

\subsection{Scholarly review as a precondition for implementation}
\label{sec:disc-scholarly}

Any commercial deployment of the \emph{shar\={\i}kat
al-'in\={a}n} architecture proposed in
Section~\ref{sec:alt-sharika}, or of the \emph{wa'd}-based
alternative of Section~\ref{sec:alt-wad}, requires prior review
by a constituted Shariah supervisory board (SSB) drawn from
qualified scholars. Reputable independent SSB providers include
Amanie Advisors, Shariyah Review Bureau, Islamic Finance
Council~UK, and the internal SSBs of the major Islamic banking
institutions. Certification under AAOIFI Shariah Standards,
the reference framework adopted throughout this paper, would
be the natural target, with jurisdiction-specific supplementary
review where local regulators (Bank Negara Malaysia, the
Securities Commission Malaysia, the DFSA in Dubai) operate their
own frameworks.

The role this paper aims to play in that scholarly-review process
is preparatory: by making the mathematical structure of both the
target architecture and the constraint set explicit, the analysis
reduces the interpretive burden on the reviewing scholars and
supplies them with a documented technical foundation on which
their \emph{ijtih\={a}d} can operate. It does not, and cannot,
substitute for that review.

\subsection{Regulatory feasibility}
\label{sec:disc-regulatory}

Any commercial implementation of the
\emph{shar\={\i}kat al-'in\={a}n} architecture would require
jurisdiction-specific legal and regulatory analysis; the
sketch below identifies the frameworks most likely to be
relevant in each of several major jurisdictions, without
purporting to be a legal opinion.
\begin{itemize}
    \item In \textbf{Australia}, candidate legal forms include
    the Managed Investment Scheme or Individually Managed
    Account frameworks under Chapter~5C of the Corporations
    Act; product-specific structuring, an Australian Financial
    Services Licence with appropriate authorisations, and
    ASIC engagement would all be required
    (\citealp{ASIC2020}).
    \item In the \textbf{United Kingdom}, Part~4A authorisation
    under the FCA Handbook would apply; the FCA has engaged
    with Islamic finance product categories through
    accommodating guidance and prior policy statements
    \citep{FCA2022}.
    \item In the \textbf{United Arab Emirates}, the DFSA
    describes itself as a Shariah-\emph{systems} regulator
    rather than a Shariah regulator, and requires firms
    conducting Islamic financial business to maintain a
    Shariah supervisory board and governance framework. The
    DFSA thus provides a dedicated framework for Islamic
    financial business but does not itself certify Shariah
    compliance; legal and regulatory classification of the
    proposed structure would require jurisdiction-specific
    analysis and DFSA engagement.
    \item In \textbf{Malaysia}, Bank Negara Malaysia's Shariah
    Advisory Council has previously ruled on
    \emph{shar\={\i}kah}- and \emph{wa'd}-based structures.
    Product-specific engagement with BNM and the Securities
    Commission Malaysia would be required.
    \item In the \textbf{United States}, the applicable
    regulatory regime depends on the product's legal form:
    off-exchange retail FX activity is regulated by the CFTC
    and NFA under the Retail Foreign Exchange Dealer
    framework; equity brokerage falls under SEC/FINRA
    supervision; investment-adviser activity under the
    Investment Advisers Act. Shariah-compliance
    representations would be treated as marketing statements
    subject to the anti-fraud provisions of the relevant
    regime.
\end{itemize}
Full regulatory analysis for any specific product is beyond
the paper's scope. Implementation constraints are substantial
(Section~\ref{sec:disc-practical}) and are not resolved by the
paper's mathematical treatment.

\subsection{Practical implementation barriers}
\label{sec:disc-practical}

The \emph{shar\={\i}kat al-'aqd} structure of
Definition~\ref{def:sharika} imposes real infrastructure
requirements on any commercial implementation:
\begin{itemize}
    \item \emph{Custody.} The segregated custody of partnership
    assets (clause~3 of Definition~\ref{def:sharika}) requires
    a nominee-custodian arrangement satisfying the
    operational-\emph{qab\d{d}} criteria (i)--(iv) established
    in the (D2) step of Proposition~\ref{prop:sharika-comply},
    under the applicable safe-custody regime (in Australia, the
    ASIC provisions on custodial services including RG~133; in
    the United Kingdom, FCA CASS~6). Standard broker-dealer
    custody plumbing supports segregated custody in most
    jurisdictions but requires bespoke legal documentation
    identifying the partnership as the beneficial account
    holder and enforceable partnership rights against the
    custodian.
    \item \emph{Settlement.} Settlement cycles vary by asset
    and jurisdiction: US equities moved to T+1 on 28~May~2024;
    UK, EU, and Australian equities remain predominantly on T+2
    at the time of writing; spot FX settlement is
    pair- and convention-dependent (typically T+2 for major
    pairs, T+1 for USD/CAD). No exotic settlement mechanics are
    required for any of these. Physical commodity underlyings
    (allocated gold, LME warrants) require additional custody
    arrangements but remain tractable.
    \item \emph{Liquidity provisioning.} The broker's capital
    contribution $K_b$ must be warehoused for the duration of
    the venture. This is economically analogous to prime-brokerage
    balance-sheet commitment but structured contractually as
    equity rather than debt, and requires the broker to hold
    sufficient regulatory capital to support the aggregate
    contribution across its client book.
    \item \emph{Client onboarding.} Additional documentation
    (Shariah-compliance attestations, disclosure of joint
    ownership) increases onboarding friction relative to the
    click-through CFD account. This is a commercial constraint,
    not a technical one.
\end{itemize}
None of these barriers is fundamental. All are recognisable
extensions of existing broker infrastructure and legal practice.

\subsection{Model limitations}
\label{sec:disc-model}

The Bates jump-diffusion adopted in Assumption~\ref{ass:bates} is
a well-tested workhorse of quantitative finance, sufficient for
the paper's purpose, namely, comparing architectures under a
common realistic-enough price process, but not the last word
on empirical price dynamics. Model features not captured include:
\begin{itemize}
    \item \emph{Regime switching.} Discrete shifts between low-
    and high-volatility regimes, well-documented in FX and
    equity markets \citep{Hamilton1989}, are absent from the
    single-regime Bates specification.
    \item \emph{Order-book microstructure.} The paper's
    execution model treats spread as an exogenous parameter;
    real-world execution involves adverse selection, latency,
    and inventory dynamics analysed in the microstructure
    literature (\citealp{Kyle1985,GlostenMilgrom1985}).
    \item \emph{Liquidity dry-up.} The Bates model assumes
    continuous execution at posted prices; in crisis conditions
    the bid-ask spread widens materially and depth vanishes.
    \item \emph{Parameter non-stationarity.} The illustrative
    calibrations of Table~\ref{tab:calibration} are held
    constant; real markets exhibit parameter drift that a
    production implementation would need to accommodate
    through regular recalibration.
    \item \emph{Alternative jump distributions.} The
    lognormal-jump specification is standard but restrictive;
    heavier-tailed alternatives (double-exponential,
    variance-gamma) may better capture observed left-tail
    behaviour on some underlyings.
\end{itemize}
Extension of the analysis to any of these model refinements is
straightforward given the released code base and is a natural
subject for follow-up work.

\subsection{Ethical considerations}
\label{sec:disc-ethics}

A Shariah-compliant leveraged retail product would raise, in
principle, the same investor-protection concerns as its
conventional counterpart. The \emph{shar\={\i}kat al-'in\={a}n}
architecture of Definition~\ref{def:sharika}, however, does not
admit leverage within its joint-ownership construction
(Remark~\ref{rem:sharika-no-leverage}). This is not a design
choice imposed for prudential reasons; it is a property of the
specific joint-ownership construction under the paper's adopted
constraints. The absence of leverage under this candidate
architecture is associated with the empirical evidence in
Section~\ref{sec:motivation} that leveraged retail participation
is associated with substantial realised losses among retail
participants (\citealp{HeimerSimsek2019};
\citealp{Barber2014Taiwan}); readers may draw their own
conclusions about the wider implications.

At a broader level, the availability of Shariah-compliant
financial products supports religious autonomy for observant
participants. The programme of research this paper initiates
aims to describe and formalise architectures consistent with
the constraint set that a substantial share of the global
population regards as binding.

\section{Conclusion and Future Research}
\label{sec:conclusion}

\subsection{Summary of findings}
\label{sec:concl-summary}

The paper's central negative documentation, developed in
Section~\ref{sec:shariah-analysis} as a sequence of five formal
propositions, is that the standard retail
Contract-for-Difference exhibits five structural
incompatibilities or concerns under the adopted AAOIFI/OIC
constraint framework: an interest-rate-linked and credit-plus-
increment financing structure (Proposition~\ref{prop:riba});
substantial \emph{gharar} concerns arising from cash-settled
contingent payoff without ownership or delivery, forced-
liquidation risk, and asymmetric execution
(Proposition~\ref{prop:gharar}); a zero-sum counterparty
structure on the internalised fraction of notional under B-book
execution (Proposition~\ref{prop:maysir}); non-fit with the
classical nominate contract categories, including absence of
\emph{qab\d{d}} (Proposition~\ref{prop:qabd}); and a
reciprocal-deferred-contingent obligation structure that engages
the \emph{bay' al-k\={a}li' bi'l-k\={a}li'} prohibition on the
broad reading (Proposition~\ref{prop:bay-kali}). The
``swap-free Islamic account'' variant marketed by conventional
brokers, analysed in Section~\ref{sec:swap-free}, addresses at
most the first of these concerns and only in part.

The paper's central positive contribution, developed in
Section~\ref{sec:alt-sharika}, is a market-access architecture
built on the classical \emph{shar\={\i}kat al-'in\={a}n}
joint-equity partnership contract, together with a proof that
this architecture is formally consistent with the operational
design constraints (D1)--(D5) adopted in this paper, subject to
review by qualified Shariah supervisory boards
(Proposition~\ref{prop:sharika-comply}). The architecture
provides the client with genuine economic exposure to a chosen
deliverable underlying at unit leverage, with the broker acting
as an agent-manager on the same side of the venture at the
gross level rather than as an opposing counterparty
(Remark~\ref{rem:sharika-equivalent-spot}). The absence of
leverage is a structural property of this specific
joint-ownership construction under the paper's adopted
constraints (Remark~\ref{rem:sharika-no-leverage}); the paper
does not claim that all Shariah-compliant financing anywhere is
incapable of leverage.

The paper's central simulation-based quantitative claim,
established in Section~\ref{sec:simulation} through Monte Carlo
simulation under illustrative Bates parameterisations across
four underlyings,
four holding horizons, and four leverage regimes, is a
quantification of the compliance cost of the Shariah leverage
cap. Because the compliant architecture cannot be leveraged, its
comparison with a leveraged CFD is not a like-for-like exposure
comparison and does not support a superiority claim in the
ordinary financial sense. What the simulation \emph{does}
establish is: (i) at unit exposure, the underlying-return
component of the two architectures is closely matched, before
architecture-specific fees, execution mechanics, and cash-flow
definitions;
(ii) the CFD's fee structure scales with the leveraged notional
whereas the \emph{shar\={\i}kat} ujrah does not, producing a
crossover in absolute cost of carry at approximately $L \geq 5$
for holding horizons above one week; and (iii) at higher
leverage regimes, the CFD introduces margin-call forced-liquidation
risk that the unleveraged \emph{shar\={\i}kat} structurally
cannot exhibit. At $L = 10$ on the S\&P~500 over a 63-day
horizon, the leveraged CFD triggers margin-call events on 48\%
of paths; no maintenance-margin calls were observed in the
10{,}000 simulated \emph{shar\={\i}kat} paths, and more
fundamentally the \emph{shar\={\i}kat} architecture contains no
price-triggered maintenance-margin close-out mechanism. At
$L = 30$ the CFD saturates the
negative-balance protection floor on 75\% of paths
(Table~\ref{tab:tail-metrics}). These figures describe the
trade the Muslim participant faces when choosing between the
architectures: reduced exposure in exchange for the elimination
of the maintenance-margin close-out mechanism and formal
consistency with the constraint set.

\subsection{Contributions restated}
\label{sec:concl-contributions}

The five contributions announced in
Section~\ref{sec:contributions} have been developed in this
paper:
\begin{enumerate}
    \item A formal characterisation of the standard CFD/FX
    architecture as a coupled stochastic system
    (Section~\ref{sec:cfd-mechanics}).
    \item Formal propositions documenting five structural
    incompatibilities or concerns under the adopted AAOIFI/OIC
    constraint framework (Section~\ref{sec:shariah-analysis}).
    \item Three first-principles alternative architectures with
    explicit payoff dynamics, one of which
    (\emph{shar\={\i}kat al-'in\={a}n}) is formally consistent
    with the operational design constraints adopted in this
    paper (Section~\ref{sec:alternatives}).
    \item A simulation-based quantitative comparison of all-in
    holding cost, tail risk, and margin-call frequency across
    architectures, framed as a compliance-cost analysis
    (Section~\ref{sec:simulation}).
    \item An open-source Python reference implementation
    released under the MIT License
    (Section~\ref{sec:code-availability}).
\end{enumerate}

\subsection{Directions for future research}
\label{sec:concl-future}

The title of this paper, ``Shariah Computational Financial
Engineering~I'', signals its role as the opening entry in a
planned series. Subsequent papers in the series will address the
following directions, each of which extends the framework of
Definitions~\ref{def:shariah-valid}
and~\ref{def:compliant-arch} into a new instrument class or
methodological setting.

\emph{Paper II: Shariah-compliant portfolio optimisation under
AAOIFI screening constraints.} The classical Markowitz
mean-variance framework requires modification to accommodate
(i)~sector exclusion constraints, (ii)~financial-ratio screening
thresholds under AAOIFI Standard~No.~21, and
(iii)~\emph{purification} of incidental impermissible income.
The resulting optimisation problem is convex under natural
formulations but requires purpose-built solver architecture; the
theoretical analysis and empirical performance comparison
against conventional benchmarks constitute the intended subject
of the next paper.

\emph{Paper III: \emph{\d{S}uk\={u}k} pricing under stochastic
profit-rate models.} The equivalent of interest-rate modelling
for the Shariah-compliant fixed income market. The analogue of
the Heath--Jarrow--Morton framework under a profit-rate rather
than interest-rate driver, calibrated against the growing
\emph{\d{s}uk\={u}k} market and applicable to profit-sharing
investment accounts as well as issued
\emph{\d{s}uk\={u}k}.

\emph{Paper IV: \emph{Tak\={a}ful} risk modelling and
Shariah-compliant actuarial reserving.} Cooperative insurance
under the \emph{mud\={a}raba} or \emph{wak\={a}la} structure
requires actuarial treatment that departs from the standard
mutual-of-strangers pooling model. Regulatory capital and
reserving under a genuinely cooperative structure is an open
technical question with substantial practical stakes.

Beyond the planned series, several natural extensions of the
present paper's approach suggest themselves:
\begin{itemize}
    \item An empirical study of observant-Muslim retail-trader
    behaviour on Shariah-compliant versus conventional accounts,
    to test whether the leverage-cap feature of the
    \emph{shar\={\i}kat} architecture improves realised
    outcomes as the theoretical analysis predicts.
    \item A deep-learning approach to Shariah-compliance
    classification for equity screening at scale, addressing
    the operational bottleneck faced by every Islamic asset
    manager.
    \item Formal analysis of systematic hedging strategies
    under the \emph{wa'd}-based framework of
    Section~\ref{sec:alt-wad}, including cost-of-carry
    comparisons against conventional NDFs and Islamic FX Forward
    products.
    \item Extension of the market-access architectures to
    decentralised-finance settings, where the joint-ownership
    structure of Definition~\ref{def:sharika} maps naturally
    onto smart-contract-based tokenised custody.
\end{itemize}

\subsection{Closing remark}
\label{sec:concl-closing}

The observant-Muslim participant base in global capital markets
is large, rising, and structurally underserved by the current
product landscape. Sophisticated tools for their genuine
participation, under structures that respect the religious
constraints that a substantial share of the world's population
regards as binding, represent both an intellectual challenge
and a real contribution to financial infrastructure. This paper
takes one small step toward that broader research programme.
Much more remains to be done.

% ========================================================================
%  APPENDICES
% ========================================================================
\appendix

\section{Full mathematical derivations}
\label{app:derivations}

This appendix collects the derivations underlying the propositions
of Sections~\ref{sec:cfd-mechanics}--\ref{sec:alternatives}. All
notation is as in the main text.

\subsection{Conditional distribution of the log-price under Bates}
\label{app:bates-density}

Let $X_t \coloneqq \ln(S_t / S_0)$ denote the log-price process
under Assumption~\ref{ass:bates}. Its characteristic function
$\phi_X(u; t) \coloneqq \EE[\exp(i u X_t) \mid \FF_0]$ admits the
semi-closed decomposition
\begin{equation}
    \phi_X(u; t)
    = \phi^{\text{Heston}}_X(u; t) \cdot \phi^{\text{jump}}_X(u; t),
    \label{eq:app-cf-decomposition}
\end{equation}
where the diffusive component follows the \citet{Heston1993}
closed form
\begin{align}
    \phi^{\text{Heston}}_X(u; t)
    &= \exp\!\bigl( C(u; t) + D(u; t) v_0 + i u \mu t \bigr),
    \label{eq:app-heston-cf} \\
    D(u; t)
    &= \frac{b_2 - d}{\sigma_v^{2}}
       \left( \frac{1 - e^{-d t}}{1 - g e^{-d t}} \right),
    \nonumber \\
    C(u; t)
    &= \frac{\kappa_v \theta}{\sigma_v^{2}}
       \left[ (b_2 - d) t
         - 2 \ln\!\left(\frac{1 - g e^{-d t}}{1 - g}\right) \right],
    \nonumber \\
    b_2 &= \kappa_v - i \rho \sigma_v u, \nonumber \\
    d   &= \sqrt{ b_2^{\,2} + \sigma_v^{2}\,(u^2 + i u) }, \nonumber \\
    g   &= (b_2 - d) / (b_2 + d), \nonumber
\end{align}
and the jump component is
\begin{equation}
    \phi^{\text{jump}}_X(u; t)
    = \exp\!\Bigl( \lambda t \bigl[
        e^{i u \mu_J - u^2 \sigma_J^{2} / 2} - 1
        - i u \kappa
      \bigr] \Bigr).
    \label{eq:app-jump-cf}
\end{equation}
The conditional density of $X_t$ is recovered by Fourier
inversion,
\begin{equation}
    f_{X_t}(x)
    = \frac{1}{2\pi}
      \int_{-\infty}^{\infty}
      e^{-i u x}\, \phi_X(u; t) \, \dd u,
    \label{eq:app-fourier-inversion}
\end{equation}
which admits numerical evaluation via the Carr--Madan FFT
scheme in settings where an analytic density is required
(\citealp{Bates1996,Bjork2009}). This Fourier machinery is
included here for analytical completeness only. The Monte Carlo
path simulation used in Section~\ref{sec:simulation} integrates
the discretised SDE of Assumption~\ref{ass:bates} directly and
does not invoke the density or the FFT inversion.

\subsection{Margin-call barrier hitting probability}
\label{app:barrier}

We derive the closed-form margin-call probability quoted in
Proposition~\ref{prop:leverage-sensitivity}. Suppose $S_t$
follows the geometric Brownian motion special case of
Assumption~\ref{ass:bates} (i.e.\ $\lambda = 0$,
$v_t \equiv \sigma^{2}$). The log-price
$X_t = \ln(S_t / S_{t_0})$ is a Brownian motion with drift:
\begin{equation}
    X_t = \left( \mu - \tfrac{1}{2}\sigma^{2} \right)(t - t_0)
          + \sigma W_{t - t_0},
    \qquad t \geq t_0.
    \label{eq:app-log-price-bm}
\end{equation}
For a long CFD position, the margin-call barrier of
\eqref{eq:mc-explicit-barrier} translates into the log-price
barrier
\begin{equation}
    b \coloneqq \ln\!\left( 1 - \frac{1 - \alpha}{L} \right)
    = -\frac{1 - \alpha}{L} + O\!\bigl( L^{-2} \bigr),
    \label{eq:app-log-barrier}
\end{equation}
so that $\tau_{\text{MC}} = \inf\{ t > t_0 : X_t \leq b \}$.

The standard first-passage distribution for Brownian motion with
drift (\citealp{Shreve2004}, Theorem~3.7.6) gives
\begin{equation}
    \PP(\tau_{\text{MC}} < h)
    = \Phi\!\left( \frac{b - \nu h}{\sigma \sqrt{h}} \right)
    + e^{2 \nu b / \sigma^{2}}
      \Phi\!\left( \frac{b + \nu h}{\sigma \sqrt{h}} \right),
    \label{eq:app-first-passage}
\end{equation}
where $\nu = \mu - \sigma^{2}/2$ and $h = t_1 - t_0$. Setting
$\nu = 0$ (leading-order approximation on short horizons) and
substituting the leading term of \eqref{eq:app-log-barrier}
yields
\begin{equation}
    \PP(\tau_{\text{MC}} < h)
    \approx 2\, \Phi\!\left(
      -\frac{1 - \alpha}{L \sigma \sqrt{h}}
    \right),
    \label{eq:app-mc-approx}
\end{equation}
which is the approximation quoted in
Proposition~\ref{prop:leverage-sensitivity}. Writing
$x \coloneqq (1-\alpha)/(L\sigma\sqrt{h})$, the leverage
sensitivity follows from the standard Taylor expansion of the
standard normal CDF near the origin,
\begin{equation}
    \Phi(-x) = \tfrac{1}{2} - \frac{x}{\sqrt{2\pi}}
             + O(x^3), \qquad x \to 0,
    \label{eq:app-phi-expansion}
\end{equation}
which under $L \to \infty$ (hence $x \to 0$) gives
\begin{equation}
    \PP(\tau_{\text{MC}} < h)
    = 1 - \sqrt{\tfrac{2}{\pi}}
        \cdot \frac{1 - \alpha}{L \sigma \sqrt{h}}
      + O\!\bigl( L^{-3} \bigr)
    \;\longrightarrow\; 1 \quad \text{as } L \to \infty.
    \label{eq:app-mc-scaling}
\end{equation}
Note: the more familiar Mills-ratio asymptotic
$\Phi(-x) \sim \phi(x)/x$ holds as $x \to \infty$; here we are in
the opposite regime ($x \to 0$) and the expansion of
\eqref{eq:app-phi-expansion} is the correct one. The limit
$\PP \to 1$ is intuitive: at infinite leverage the barrier
coincides with the entry price and any downward Brownian
excursion within horizon $h$ triggers a margin call.

For the full Bates dynamics ($\lambda > 0$),
\eqref{eq:app-first-passage} is \emph{not} in general a strict
bound in either direction: the log-jump-size distribution
$\ln J_t \sim \mathcal{N}(\mu_J, \sigma_J^2)$ admits both
$J_t > 1$ and $J_t < 1$, so a jump event can move the path
either towards or away from a downward barrier. The direction
and magnitude of the effect depend on
$(\lambda, \mu_J, \sigma_J)$ and the compensator adjustment
$-\lambda\kappa$; the net contribution is quantified numerically
in Section~\ref{sec:simulation} (Figure~\ref{fig:tail-decomp}).

\subsection{Wa'd payoff derivation and cost-of-carry decomposition}
\label{app:wad}

Under Definition~\ref{def:wad-fx}, at delivery time $T$ the
client and broker execute a spot FX transaction at the pre-agreed
rate $K$. The client's cash flows in CCY2 are:
\begin{itemize}
    \item At $t_0$: pay ujrah $u^{\text{wa'd}}$ (if front-loaded)
    or accrue it to $T$.
    \item At $T$: pay $N \cdot K$ units of CCY2, receive $N$
    units of CCY1.
    \item Immediately after $T$: sell $N$ units of CCY1 at spot
    rate $S_T$, receive $N \cdot S_T$ units of CCY2.
\end{itemize}
Netting yields $\Pi^{\text{wa'd}}_c(t_0, T) = N (S_T - K)
- u^{\text{wa'd}}$, which is
equation~\eqref{eq:wad-payoff}.

Under the risk-neutral measure $\PP^{*}$ implied by covered
interest rate parity, the discounted forward rate satisfies
$\EE^{*}[S_T \mid \FF_{t_0}] = F^{\text{CIP}}_0$, so
\begin{equation}
    \EE^{*}\bigl[ \Pi^{\text{wa'd}}_c(t_0, T) \mid \FF_{t_0} \bigr]
    = N \bigl( F^{\text{CIP}}_0 - K \bigr) - u^{\text{wa'd}}.
    \label{eq:app-wad-riskneutral}
\end{equation}
The economic cost of the wa'd relative to a conventional NDF
struck at $K = F^{\text{CIP}}_0$ is thus exactly $u^{\text{wa'd}}$
in expected value. Under the physical measure $\PP$, the
comparison depends on the market's risk premium, which is not
constrained by the wa'd mechanics but only by the joint
calibration decision on $K$ and $u^{\text{wa'd}}$ discussed in
Remark~\ref{rem:wad-debate}.

\subsection{\emph{Shar\={\i}kat al-'in\={a}n} payoff, extended
derivation}
\label{app:sharika}

We extend the proof of Proposition~\ref{prop:sharika-payoff} to
the general case of unequal profit and loss shares under the
AAOIFI Standard~No.~12 \S3/1/5/2 constraint that loss shares
equal capital contribution ratios.

Let $\pi_c, \pi_b$ denote the profit-sharing ratios with
$\pi_c + \pi_b = 1$, and let $\ell_c = K_c / (K_c + K_b)$,
$\ell_b = K_b / (K_c + K_b)$ denote the loss-sharing ratios.
The joint venture's terminal P\&L is
$\Delta \coloneqq N \cdot (S_{t_1} - S_{t_0})$, where
$N = (K_c + K_b) / S_{t_0}$. The client's share is
\begin{equation}
    \Pi_c(\Delta)
    = \begin{cases}
        \pi_c \cdot \Delta - u^{\text{shr}}_c
            & \text{if } \Delta \geq 0, \\
        \ell_c \cdot \Delta - u^{\text{shr}}_c
            & \text{if } \Delta < 0.
      \end{cases}
    \label{eq:app-sharika-general}
\end{equation}
Under the AAOIFI-compliant proportional case
$\pi_c = \ell_c = K_c / (K_c + K_b)$, the two branches of
\eqref{eq:app-sharika-general} collapse to the single expression
\begin{equation}
    \Pi_c
    = \frac{K_c}{K_c + K_b} \cdot N \cdot (S_{t_1} - S_{t_0})
      - u^{\text{shr}}_c
    = K_c \cdot \left( \frac{S_{t_1}}{S_{t_0}} - 1 \right)
      - u^{\text{shr}}_c,
    \label{eq:app-sharika-proportional}
\end{equation}
recovering equation~\eqref{eq:sharika-payoff}. In the general
case with $\pi_c > \ell_c$ (higher profit share than loss share
for one party, compensating that party for taking on
disproportionate management responsibility or for
non-capital contributions to the venture, per AAOIFI SS~No.~12
\S3/1/5), the payoff is a piecewise-linear call-like structure
that may remain permissible under the applicable partnership
rules, provided the loss-sharing constraint $\ell_c = K_c /
(K_c + K_b)$ is satisfied (an AAOIFI SS~No.~12 constraint that
forbids losses being borne otherwise than in proportion to
capital) and all other contractual requirements are met,
subject to SSB review; the client's expected profit is enhanced
relative to the proportional case, with the excess taken from
the broker's profit share.

\section{Fiqh sources catalogued}
\label{app:fiqh-sources}

Table~\ref{tab:fiqh-catalogue} consolidates the primary fiqh
sources cited across the paper. AAOIFI Shariah Standards
are cited by number and, where relevant, by section reference;
OIC Fiqh Academy resolutions by number and session; contemporary
scholarly works by author and section reference within the cited
text. This catalogue is intended as a navigation aid for
scholarly review of the paper's compliance claims.

\begin{table}[!p]
\centering
\scriptsize
\setlength{\tabcolsep}{4pt}
\renewcommand{\arraystretch}{1.05}
\caption{Primary fiqh sources cited in this paper. AAOIFI SS
denotes AAOIFI Shariah Standard; OIC Res. denotes International
Islamic Fiqh Academy Resolution.}
\label{tab:fiqh-catalogue}
\begin{tabular}{@{}p{4.3cm}p{4.3cm}p{4.3cm}@{}}
\toprule
\textbf{Source} & \textbf{Subject} & \textbf{Cited in} \\
\midrule
AAOIFI SS No.~1
    & Trading in Currencies; possession requirements, deferment
      prohibition, restrictions on forward/futures currency
      transactions and on binding bilateral currency promises
      (\S2/9)
    & Prop.~\ref{prop:riba} (in support of the general \emph{rib\={a}}
      argument), Prop.~\ref{prop:wad-comply} (constraint on
      Architecture A), Rem.~\ref{rem:riba-foreshadow} \\
AAOIFI SS No.~8
    & \emph{Mur\={a}ba\d{h}a} and credit-sale increments;
      prohibitions related to \emph{bay' al-k\={a}li' bi'l-k\={a}li'}
    & Prop.~\ref{prop:bay-kali},
      Table~\ref{tab:violation-summary} \\
AAOIFI SS No.~12 \S3/1/5/2
    & \emph{Shar\={\i}kah}; proportional loss-sharing rule
    & Def.~\ref{def:sharika}, App.~\ref{app:sharika} \\
AAOIFI SS No.~18
    & \emph{Qab\d{d}} (constructive possession) requirement
    & Prop.~\ref{prop:qabd},
      Table~\ref{tab:violation-summary} \\
AAOIFI SS No.~20
    & Sale of commodities in organised markets
    & Section~\ref{sec:lit-derivatives} \\
AAOIFI SS No.~21
    & Financial-ratio screening for Shariah-compliant equities
    & Section~\ref{sec:alt-sharika}, Section~\ref{sec:concl-future} \\
AAOIFI SS No.~27
    & Indices
    & Section~\ref{sec:lit-derivatives} \\
AAOIFI SS No.~30 \S5
    & Organised \emph{tawarruq}
    & Prop.~\ref{prop:tawarruq-comply} \\
AAOIFI SS No.~31 \S5/1
    & Controls on \emph{gharar} in financial transactions
    & Prop.~\ref{prop:gharar},
      Table~\ref{tab:violation-summary} \\
OIC Res.~63 (1/7), 1992
    & Financial markets; definition of \emph{maysir}
    & Prop.~\ref{prop:maysir},
      Table~\ref{tab:violation-summary} \\
OIC Res.~179 (19/5), 2009
    & Organised \emph{tawarruq}; impermissibility ruling
    & Prop.~\ref{prop:tawarruq-comply},
      Rem.~\ref{rem:tawarruq-criticism} \\
OIC Res.~238 (24/2), 2019
    & Mutual binding commitments in future currency exchange
    & Rem.~\ref{rem:wad-debate},
      Prop.~\ref{prop:wad-comply} \\
AAOIFI SS No.~49
    & Unilateral and Bilateral Promise
      (\emph{al-wa'd} and \emph{al-muw\={a}'adah})
    & Rem.~\ref{rem:wad-debate} \\
AAOIFI Governance
    Standard No.~1
    & Shariah Supervisory Board: appointment, composition, and
      report
    & Rem.~\ref{rem:sharika-governance} \\
\d{H}ad\={\i}th of Ibn Umar
    & Prophetic prohibition against
      \emph{bay' al-k\={a}li' bi'l-k\={a}li'}
    & Prop.~\ref{prop:bay-kali} \\
Qur'\={a}n 2:275
    & Prohibition of \emph{rib\={a}}
    & Section~\ref{sec:lit-fiqh} \\
Qur'\={a}n 5:90
    & Prohibition of \emph{maysir}
    & Section~\ref{sec:lit-fiqh} \\
\citet{Usmani2002}, ch.~3
    & Contemporary treatment of \emph{rib\={a}}
    & Prop.~\ref{prop:riba},
      Rem.~\ref{rem:marketing-vs-structural} \\
\citet{Kamali2000}, ch.~5
    & Contemporary treatment of \emph{gharar} thresholds
    & Prop.~\ref{prop:gharar} \\
\citet{ElGamal2006}
    & Islamic finance law, economics, and practice
    & Rem.~\ref{rem:wad-debate},
      Rem.~\ref{rem:marketing-vs-structural} \\
\bottomrule
\end{tabular}
\end{table}

\section{Python code documentation}
\label{app:code-doc}

This appendix documents the module structure of the reference
Python implementation released at
\url{https://github.com/Muhammad-Taha-Asif/Shariah-CFE-1}. Each
module is self-contained and can be imported independently.

\paragraph{\texttt{src/bates.py}}
Implements Assumption~\ref{ass:bates}. Exports:
\begin{itemize}
    \item \texttt{BatesParams}: a frozen dataclass encoding the
    ten parameters of the Bates specification, with helper
    properties for the mean jump size and the Feller condition
    check.
    \item \texttt{simulate\_paths(params, n\_paths, horizon\_years,
    n\_steps, seed)}: vectorised Monte Carlo simulator using
    Euler--Maruyama on the log-price with full-truncation of the
    variance process to preserve positivity, following
    \citet{Andersen2009}. Returns paired arrays of price and
    variance paths.
    \item \texttt{preset(name)}: returns the illustrative
    calibration presets of Table~\ref{tab:calibration} for
    ``EURUSD'', ``USDJPY'', ``AUDUSD'', ``SPX''.
\end{itemize}

\paragraph{\texttt{src/contracts.py}}
Implements the CFD contract mechanics of
Section~\ref{sec:cfd-mechanics} and the \emph{shar\={\i}kat}
architecture of Section~\ref{sec:alt-sharika}. Exports:
\begin{itemize}
    \item \texttt{CFDParams}, \texttt{SharikaParams}: dataclasses
    encoding contract parameters.
    \item \texttt{cfd\_payoff\_with\_margin\_call}: computes CFD
    gross P\&L, accumulated Tom-Next financing, and net client
    P\&L accounting for margin-call forced liquidation and
    negative-balance protection.
    \item \texttt{margin\_call\_time}: returns the first-hitting
    index of the barrier of equation~\eqref{eq:mc-explicit-barrier}
    for each simulated path.
    \item \texttt{sharika\_payoff}: computes the client's net
    payoff under Proposition~\ref{prop:sharika-payoff}.
\end{itemize}

\paragraph{\texttt{src/metrics.py}}
Statistical machinery for Section~\ref{sec:sim-metrics}. Exports
\texttt{moments}, \texttt{var\_cvar}, \texttt{sharpe},
\texttt{sortino}, \texttt{calmar}, \texttt{bootstrap\_ci},
\texttt{ks\_test}, and \texttt{summarise}. All operate on
one-dimensional P\&L arrays with the sign convention that
positive values are gains.

\paragraph{\texttt{src/simulate.py}}
Master grid runner for Section~\ref{sec:simulation}. Iterates
over four underlyings, four horizons, four leverages, and
runs \texttt{run\_cell()} on each combination, persisting the
results to \texttt{tables/results\_grid.csv}. All random seeds
are derived from a session-wide \texttt{SEED\_BASE = 20260928}
for exact reproducibility.

\paragraph{\texttt{src/figures.py}}
Publication-quality figure generation for
Section~\ref{sec:sim-results}. Produces
\texttt{fig1\_cost\_of\_carry.pdf} through
\texttt{fig4\_tail\_decomposition.pdf} with typography configured
to match the paper's Latin Modern Roman body text.

\paragraph{\texttt{src/ml\_analysis.py}}
XGBoost margin-call classifier of
Section~\ref{sec:sim-ml}. Trains the classifier on early-path
features and generates
\texttt{fig5\_feature\_importance.pdf}.

Dependencies are declared in \texttt{requirements.txt} at the
repository root. The full simulation grid runs in approximately
two minutes on commodity laptop hardware.

\section{Data sources and calibration protocol}
\label{app:data}

The Monte Carlo results in Section~\ref{sec:simulation} are
generated under the illustrative Bates parameter presets of
Table~\ref{tab:calibration}. These presets are consistent in
sign and order of magnitude with published FX and equity
implied-volatility studies \citep{Andersen2009,Bates1996} and
with contemporaneous option-market observations for the four
underlyings, but are not the result of a proprietary
calibration.

For a production implementation of the analysis, the
recommended data pipeline is as follows.

\paragraph{Tick-level FX data.} Free public tick data for the
three FX pairs is available from Dukascopy Bank SA (Geneva) via
their historical data feed. Data covering the period 2015--2025
is required to span the low-volatility regime (2015--2018),
the moderate-volatility period (2019--2022), and the elevated
regime (2022--2025). Data acquisition should respect Dukascopy's
terms of service, which permit academic research use but not
onward redistribution; a download script for the recommended
snapshots is included in the repository (\texttt{data/}
directory) but no raw data are checked in.

\paragraph{Equity and index data.} Historical minute-bar data
for the S\&P~500 index and its ETF proxies (SPY, VOO) is
available from the CBOE DataShop or, for daily data,
Yahoo!~Finance via the \texttt{yfinance} Python package. Bid-ask
spread data at the tick level requires a paid data vendor
(TAQ, Bloomberg B-Pipe).

\paragraph{Interbank overnight rates.} Data for SOFR, ESTR,
SONIA, TONA, and the RBA cash rate is available via the
respective central bank public data portals at daily frequency;
intraday granularity, where needed for cross-checking with
Tom-Next swap rates, is available from the OIS market data
providers.

\paragraph{Data cleaning protocol.} A production calibration
would apply the following cleaning steps in order:
(i) exclude weekend and public-holiday timestamps;
(ii) exclude tick-level records with zero volume, zero spread,
or a bid--ask cross;
(iii) apply a \emph{jump-preserving} tick-cleaning procedure,
such as the \citet{BrownleesGallo2006} algorithm or a
comparable rule that removes microstructure noise (mistaken
ticks, feed errors, quotes with implausible spreads) without
deleting genuine economic jumps, a blanket $\pm 5\sigma$
filter is expressly not applied at this stage, since the
Bates specification aims to \emph{estimate} jump-intensity and
jump-size parameters and a naive tail-cut would bias
$\lambda$, $\mu_J$, and $\sigma_J$ downward;
(iv) resample from tick to five-minute bars to reduce
microstructure noise while preserving jump signatures;
(v) construct daily aggregate return, realised variance, and
realised skewness series for calibration input; the
distinction between bad-tick outliers and genuine jumps in the
resulting series is then handled inside the Bates calibration
itself rather than by pre-filtering.

\paragraph{Parameter calibration.} The Bates parameters are
recoverable from either (a) time-series estimation using
particle filtering or MCMC on the coupled SDE system, or (b)
cross-sectional calibration against traded option-market implied
volatility surfaces via Fourier pricing on the characteristic
function of \ref{app:bates-density}. The latter approach
is preferred where an option market exists for the underlying;
for underlyings without a liquid option market, time-series
calibration is the default. Regular recalibration (weekly or on
material market events) is recommended.

Full implementation of the calibration pipeline is deferred to
Paper~II of this series (Section~\ref{sec:concl-future}) and is
not required for the architectural comparison that is the
subject of the present paper.

% ---------- Bibliography ----------
\bibliographystyle{elsarticle-harv}
\bibliography{references}

\end{document}
